% Options for packages loaded elsewhere
\PassOptionsToPackage{unicode}{hyperref}
\PassOptionsToPackage{hyphens}{url}
%
\documentclass[
  10pt,
]{scrreprt}
\usepackage{amsmath,amssymb}
\usepackage{setspace}
\usepackage{iftex}
\ifPDFTeX
  \usepackage[T1]{fontenc}
  \usepackage[utf8]{inputenc}
  \usepackage{textcomp} % provide euro and other symbols
\else % if luatex or xetex
  \usepackage{unicode-math} % this also loads fontspec
  \defaultfontfeatures{Scale=MatchLowercase}
  \defaultfontfeatures[\rmfamily]{Ligatures=TeX,Scale=1}
\fi
\usepackage{lmodern}
\ifPDFTeX\else
  % xetex/luatex font selection
  \setmainfont[]{DejaVu Serif}
  \setsansfont[]{DejaVu Sans}
  \setmonofont[]{DejaVu Sans Mono}
\fi
% Use upquote if available, for straight quotes in verbatim environments
\IfFileExists{upquote.sty}{\usepackage{upquote}}{}
\IfFileExists{microtype.sty}{% use microtype if available
  \usepackage[]{microtype}
  \UseMicrotypeSet[protrusion]{basicmath} % disable protrusion for tt fonts
}{}
\makeatletter
\@ifundefined{KOMAClassName}{% if non-KOMA class
  \IfFileExists{parskip.sty}{%
    \usepackage{parskip}
  }{% else
    \setlength{\parindent}{0pt}
    \setlength{\parskip}{6pt plus 2pt minus 1pt}}
}{% if KOMA class
  \KOMAoptions{parskip=half}}
\makeatother
\usepackage{xcolor}
\usepackage[margin=22mm]{geometry}
\usepackage{longtable,booktabs,array}
\usepackage{calc} % for calculating minipage widths
% Correct order of tables after \paragraph or \subparagraph
\usepackage{etoolbox}
\makeatletter
\patchcmd\longtable{\par}{\if@noskipsec\mbox{}\fi\par}{}{}
\makeatother
% Allow footnotes in longtable head/foot
\IfFileExists{footnotehyper.sty}{\usepackage{footnotehyper}}{\usepackage{footnote}}
\makesavenoteenv{longtable}
\setlength{\emergencystretch}{3em} % prevent overfull lines
\providecommand{\tightlist}{%
  \setlength{\itemsep}{0pt}\setlength{\parskip}{0pt}}
\setcounter{secnumdepth}{5}
\usepackage{amsmath,amssymb,mathtools,fvextra,xurl,microtype,needspace,longtable,booktabs,array}
\usepackage{scrlayer-scrpage}
\clearpairofpagestyles
\ihead{Split-Zero / Benzel P7}\ohead{Cumulative research record}\cfoot{\pagemark}
\setkomafont{pageheadfoot}{\normalfont\small}
\setkomafont{disposition}{\normalfont\bfseries}
\setlength{\emergencystretch}{3em}
\DefineVerbatimEnvironment{Highlighting}{Verbatim}{breaklines,breakanywhere,commandchars=\\\{\},fontsize=\footnotesize}
\RecustomVerbatimEnvironment{verbatim}{Verbatim}{breaklines,breakanywhere,fontsize=\footnotesize}
\fvset{breaklines=true,breakanywhere=true,fontsize=\footnotesize}
\AtBeginEnvironment{longtable}{\footnotesize}
\ifLuaTeX
  \usepackage{selnolig}  % disable illegal ligatures
\fi
\usepackage{bookmark}
\IfFileExists{xurl.sty}{\usepackage{xurl}}{} % add URL line breaks if available
\urlstyle{same}
\hypersetup{
  pdftitle={Split-Zero cohomology and Benzel Problem 7},
  hidelinks,
  pdfcreator={LaTeX via pandoc}}

\title{Split-Zero cohomology and Benzel Problem 7}
\usepackage{etoolbox}
\makeatletter
\providecommand{\subtitle}[1]{% add subtitle to \maketitle
  \apptocmd{\@title}{\par {\large #1 \par}}{}{}
}
\makeatother
\subtitle{Cumulative source record --- through turn 9}
\RedeclareSectionCommand[tocnumwidth=2.8em]{chapter}
\makeatletter
\renewcommand{\@pnumwidth}{2em}
\makeatother
\author{}
\date{17 September 2026}

\begin{document}
\maketitle

\chapter*{Note on this reading edition}
This edition includes the new turn-9 construction for
\(q=3k+1\), \(k\ge4\), \(m\ge3k+3\), and \(d\ge m\),
following the recovered proof editions through turn 8. The new manuscript
remains a research proof pending independent mathematical review; it does
not assert a complete solution of Propp Problem 7.

The earlier reading guide and proof bodies below are retained as historical
editions. Their descriptions of then-current coverage, remote integration,
and further work refer to their own dates. The final part supplies the
turn-9 continuation. Current entry points and integration provenance are
at \url{https://github.com/KokunoYumeto/benzel-workbench}.

The supplied cumulative archive contained editable LaTeX and the older
through-turn-8 PDF, but its recorded turn-9 PDF build had failed. This
reading edition repairs the obsolete font commands in the appended
TeX, updates its chapter hierarchy and the subtitle, adjusts contents
spacing, and adds this note. Original source files,
historical variants, and the unchanged supplied archive are preserved
separately. Compilation is not mathematical certification.
\clearpage

{
\setcounter{tocdepth}{0}
\tableofcontents
}
\setstretch{1.06}
\chapter{Reading guide and mathematical
state}\label{reading-guide-and-mathematical-state}

\textbf{Cumulative snapshot through focused turn 8; assembled 17
September 2026.}

This volume preserves the complete recovered proof editions, their
original parameter domains, and their verification boundaries. It does
not declare the full original Propp Problem 7 resolved. Independent
mathematical review, priority certification, and a new Lean build are
not asserted. Exact computation and source compilation retain their
narrower scopes.

The archive contains nine unchanged delivered ZIPs, extracted sources,
independently uploaded variants, mathematical checkers, saved
certificates, earlier patches, and the original non-zeta portfolio. The
source catalogue identifies the original path and SHA-256 of every
reproduced proof. The following chapters retain historical next-step
statements even where a later chapter advances beyond them. Repeated
arguments are separate editions, not additional discoveries.

\textbf{Edition boundary.} The concurrent turn-2 third-collar proof and
the later fourth-collar delivery are both preserved. Turn 7 was
delivered as code, parameter tables, and a detailed conversation report,
without a completed standalone manuscript. Its code is included and its
mathematical report is identified below; no fictitious recovered
manuscript is supplied. Remote integration commit
\path|038ed02bb524e60e9c2bd20c3a0aa02365f4001d| additionally contains
\texttt{TURN2.md} and an integer 9-by-27 matrix replay. Their full raw
files were not materialized in this cumulative handback; the exact
remote locators are recorded in the README. The complete delivered
integral quotient proofs and later full comparison-kernel calculations
are included.

The latest previously read source head is
\path|20cb022dcfda7f37b3adec13b5619be8f13be3e8|, the draft PR26
contribution. This is a cumulative delivery archive, not a complete Git
clone or a claim that the local-only editions are integrated remotely.

\section{Original objects}\label{original-objects}

The original cell is \((x,y)\in\mathbb Z^2\), with

\[u=y-x,\qquad v=1-x-2y,\qquad w=2x+y-1.\]

The benzel \(V(a,b)\) consists of the cells with
\(1-a\le u,v,w\le b-1\). The retained integer lane is

\[t_d=\binom d2,\qquad W_h(d)=V(d+3h,2d+3h),\qquad\Delta=t_d-h,\]

with inverse \(d=b-a\), \(h=(2a-b)/3\). The right stone and original
bones have offsets

\[\begin{array}{c|ccc}
R&(0,0)&(1,0)&(0,1)\\
H&(0,0)&(1,0)&(2,0)\\
V&(0,0)&(0,1)&(0,2)\\
D&(0,0)&(1,-1)&(2,-2).
\end{array}\]

The original rotation \(\rho(x,y)=(y,1-x-y)\) has inverse \(\rho^2\).
Reflection \(F(x,y)=(x,1-x-y)\) has inverse \(F\). Their actions on the
difference coordinates are \((u,v,w)\mapsto(v,w,u)\) and
\((u,v,w)\mapsto(-w,-v,-u)\). The source gives

\[|W_h(d)|=3\Delta+9h(h+d).\]

\section{Actual Split-Zero objects}\label{actual-split-zero-objects}

The coefficient semiring is \(G(\mathbb Z)=\mathbb Z\sqcup\{\tau\}\).
Its supported ring zero is \(e=0_{\mathbb Z}\); its global zero is
\(\tau\). The explicit pair isomorphism sends \(\tau\) to \((0,0)\) and
\(n\) to \((n,1)\) in

\[\{(0,0)\}\cup(\mathbb Z\times\{1\})\subset\mathbb Z\times\mathbb B.\]

For original support-indexed modules \(C_A^i\) and original generator
transitions \(j_{AB}\), reconstruction gives

\[M^i=\coprod_A C_A^i,\quad (A,x)+(B,y)=(A\vee B,jx+jy),\]
\[e(A,x)=(A,0),\qquad\tau(A,x)=(\bot,0).\]

A separate bottom is retained where the empty release set has a nonempty
cell support. The incidence differential sends each original placed tile
to its three original cell generators. The internal quotient coequalizes
that incoming boundary and its support-preserving zero companion.

The actual affine filling equation is kept in an augmented linear
complex:

\[\widehat d_A(r,n)=\partial n-rv_A,\qquad v_A=\mathbf1_{\Lambda_A}.\]

For newly released original bones \(\eta_{AB}\), the partition identity
and cochain map are

\[v_B=jv_A+\partial\eta_{AB},\qquad
\widehat j_{AB}(r,n)=(r,jn+r\eta_{AB}),\]
\[\widehat d_B\widehat j_{AB}=j\widehat d_A.\]

Charge-one nonnegative integral cycles are exactly the original positive
filling fibre: an all-ones target forces one-fold cell coverage. Its
inclusion into the integer cycle fibre retains all original tile
coefficients. The proofs calculate signed lifts, original dual cochains,
and positive replacements on these same objects.

\section{Collars and the evaluated support
delay}\label{collars-and-the-evaluated-support-delay}

The recovered sources construct \(h=1,2,3\) for every \(d\ge3\), and
\(h=4,5\) for \(d\ge4\). The nine-cell relation compares the two
positive partitions

\[\begin{aligned}
\{R(x,y),H(x+2,y),H(x+1,y+1)\}
\longleftrightarrow
\{H(x,y),H(x,y+1),R(x+3,y)\}.
\end{aligned}\]

The signed tile difference has zero incidence; replacing one disjoint
patch by the other gives inverse maps on the specified original tiling
fibres.

Turn 4 studies a fixed fifth-repair dictionary of sixteen original
bones. Original coefficient evaluation gives a split copy of
\(\mathbb Z\) through prefix 14. A displayed integral signed chain kills
the distinguished class at prefix 15. A positive filling occurs at
prefix 16. Sixteen integer dual cochains exclude all 65,535 proper
release subsets. The quantified scope is the stated source tiling,
stone, translation and dictionary, not global optimization over other
configurations.

On that actual rank-one image, the Rees inclusion is

\[T^{15}\mathbb Z[T]w\hookrightarrow\mathbb Z[T]w.\]

Its defect has integral basis \(w,Tw,\ldots,T^{14}w\) and an
antidiagonal residue pairing with entries \(\delta_{i+j,14}\). The
calculated boundary character is

\[1-z^{15}=(1-z)(1+z+\cdots+z^{14}),\qquad
-\left.z\frac{d}{dz}(1-z^{15})\right|_{z=1}=15.\]

\section{Global quotient, finite jet, and actual
kernels}\label{global-quotient-finite-jet-and-actual-kernels}

The integral global bone quotient is

\[Q=\mathbb Z[X^{\pm1},Y^{\pm1}]/(1+X+X^2,\ 1+Y+Y^2,\ X^2+XY+Y^2).\]

The source gives inverse ring maps

\[Q\cong\mathbb Z[\omega]\ltimes\mathbb F_3\sigma,\]
\[\omega^2+\omega+1=0,\quad \sigma^2=0,\quad(\omega-1)\sigma=3\sigma=0,\]
\[X\mapsto\omega,\quad Y\mapsto\omega^2+\sigma;
\qquad\omega\mapsto[X],\quad\sigma\mapsto[1+X+Y].\]

Multiplication on the target is \((a,c)(b,e)=(ab,\bar a e+\bar b c)\),
with \(\bar a=a(1)\bmod3\). The finite jet

\[\mathbb F_3[\varepsilon,\eta]/(\varepsilon^2,\varepsilon\eta,\eta^2)\]

is precisely \(Q/3Q\) through

\[a+b\omega+c\sigma\mapsto(a+b)+(b+c)\varepsilon+c\eta.\]

The inverse sends \(\varepsilon\) to \(\omega-1\) and \(\eta\) to
\(\omega^2+\sigma-1\). A global or finite-jet zero is not used in place
of a support-contained positive preimage.

For two actual matchings \(A,B\) on \(\Lambda\), set
\(C=\mathbb Z[\Lambda]\) and let \(B_A,B_B\) be their incidence images.
The original representative comparison gives

\[\ker(C/B_A\to C/(B_A+B_B))\cong B_B/(B_A\cap B_B).\]

The inverse takes the new-boundary part of a killed representative; two
choices differ precisely in \(B_A\cap B_B\). The full integer kernel is
computed using the bipartite graph on actual old and new bones, with an
edge for each shared original cell. Singly owned cells force zero
coefficients. Coefficients propagate equally along shared cells. Thus
only closed components supply intersection relations, one primitive
integer relation per closed component. Deleting one new-bone pivot per
closed component gives a full integral kernel basis. Subtracting pivot
coefficients gives its coordinate map; restoring zero pivots is the
inverse. The forest duals are original integer cell cochains, not merely
a rational nullspace basis.

\section{Stable residual and the larger invariant
families}\label{stable-residual-and-the-larger-invariant-families}

The source's variable-shift packing gives

\[\partial\Gamma_\Delta(d,h)+\mathbf1_{\Omega_\Delta}=\mathbf1_{W_h(d)},
\quad|\Gamma_\Delta|=3h(h+d),\quad|\Omega_\Delta|=3\Delta.\]

The exact original-cell rank bijection is

\[(r,s,p)\mapsto\rho^p(r-s,-s),\quad r\ge1,\ 0\le s<r,\ p=0,1,2,
\qquad\operatorname{rank}=t_r+s+1.\]

Its inverse chooses the unique minimum of \(u,v,w\), rotates it to
\(u\), and recovers \(r=-u,s=-y\). The residual consists of exactly the
ranks at most \(\Delta\), independently of the exterior parameter \(d\).

At \(\Delta=t_m\), the residual is the original \(W_0(m)\) with its
explicit all-stone tiling. This includes the complete one-stone ray

\[V\left(\frac{3d^2-d-6}{2},\frac{3d^2+d-6}{2}\right),\qquad d\ge3.\]

The retained proof also gives a positive band homotopy that moves its
three separated source holes to the original central right stone,
including inverse slides. The first-collar core supplies
\(\Delta=t_m-1\) on its recorded domain.

\section{Turn 7: labelled report
reconstruction}\label{turn-7-labelled-report-reconstruction}

This paragraph reconstructs the mathematical report accompanying turn
7's preserved original code; it is not presented as a recovered proof
manuscript. The tables cover \(q\in\{1,2,3,4,5,6,9,12\}\) for
\(m\ge q+2\), \(d\ge m\), \(\Delta=t_m-q\), \(h=t_d-t_m+q\). Their
positive identity is

\[\partial B+\partial\mathcal R=\mathbf1_{\Omega_\Delta}+\partial A.\]

The receiving-tail index maps are

\[j_{\mathrm{tail}}=j_{\mathrm{stable}}+q-y,\qquad j_{\mathrm{stable}}=j_{\mathrm{tail}}+y-q.\]

They hold on the original anchors because

\[y-3q-m+1+3(q-1-\ell)=y-m-2-3\ell.\]

The full tiling adjoins \(\Gamma_\Delta\setminus A\). The report
combines the original tables and earlier certified cases to cover
\(W_6(d)\) for \(d\ge4\), \(W_9(d)\) for \(d\ge5\) and \(W_{12}(d)\) for
\(d\ge6\). The report explicitly did not complete its final consolidated
replay; the archive does not turn that limitation into a new PASS.

\section{The largest unbounded construction at this
snapshot}\label{the-largest-unbounded-construction-at-this-snapshot}

Turn 8 constructs the original positive tiling for every

\[k\ge1,\quad m\ge3k+2,\quad d\ge m,\]
\[q=3k,\quad\Delta=t_m-3k,\quad h=t_d-t_m+3k.\]

The original released and replacement source sets have size
\(3km+6k^2-3\) and satisfy

\[\partial B_{m,k}+\partial\mathcal R_{m,k}
=\mathbf1_{\Omega_\Delta}+\partial A_{m,k}.\]

The same literal table produces the constructor and a six-variable
integer Laurent certificate before either parameter is specialized.
Denominators involve only the original cell variables; the specified
injective localization returns the identity to the original cell module.
Positivity of tile coefficients and exact indicator equality establish a
partition.

The two other deletion residues retain signed channels and finite
endpoint repairs, not a completed uniform positive endpoint theorem. The
integer rotational norm is \(N=1+\rho+\rho^2\), \(N^2=3N\). Only the
central right stone is fixed, so the symmetry restriction is related to
the full positive fibre by its actual inclusion; asymmetric tilings
remain available.

\chapter{Methodology}\label{edition-1}

Source:
\path|sources/splitzero-nonzeta-workbench/splitzero-nonzeta-workbench/METHODOLOGY.md|.

SHA-256:
\path|4a59947eba9669708d904a89169e83f0a03650a5d6bb3947431f4b903745686b|.

\section{Pinned mathematical sources}\label{edition-1-section-1}

Repository: \path|KokunoYumeto/zeta-function-research-reader|, commit
\path|1c8ec52c85c173adc9f8403a8a26914955e2f5a9|.

The principal sources read are
\path|workbenches/splitzero-tandem/tex/split_zero_carriers.tex|,
\path|workbenches/splitzero-tandem/tex/support_diagrams.tex|,
\path|formal/splitzero/DERIVED_MATHEMATICS.md| and
\path|workbenches/split-support-rees-trace/RESEARCH_NOTE.md|, together
with the programme and workbench guides. UPSTREAM.lock.json supplies the
exact paths and two returned Git blob identities. These references
retain the full proofs; this document does not replace them.

The transfer programme below uses the explicit algebraic constructions,
followed by actual target-specific morphisms. No arithmetic theta
estimate is transported to another problem by changing its name. No full
analytic or Lean re-verification of the upstream repository was
performed here.

\section{1. Scalar support with its actual
projections}\label{edition-1-section-2}

For a commutative ring R, the source constructs G(R)=R disjoint union
\{tau\}. The ordinary ring zero e stays in the supported copy R. Tau is
the additive identity and multiplicatively absorbing. Supported sums and
products use the ring operations, including those whose result is e.

The map

\begin{verbatim}
eta(tau)=(0,0), eta(r)=(r,1)
\end{verbatim}

is a semiring isomorphism onto \path|{(0,0)} union (R x {1})| inside
\texttt{R\ x\ Boolean}. Addition of two supported pairs gives
\texttt{(r+s,1)} and multiplication gives \texttt{(rs,1)}; operations
involving \texttt{(0,0)} give the stated \path|identity/absorption|
laws. The inverse sends \texttt{(0,0)} to tau and \texttt{(r,1)} to r,
proving the inverse identities.

Projection p to R sends both tau and e to 0. Projection chi to Boolean
sends tau to 0 and every supported element, including e, to 1. These
formulas explicitly specify the information retained by each map. The
supported inclusion j:R-\textgreater G(R) sends ordinary zero to e,
while the ambient semiring zero is tau.

For a matrix with explicitly absent entries, evaluate the determinant's
indexed permutation sum using this lift. Applying p gives the ordinary
signed determinant with absent entries evaluated at zero. Applying chi
gives the Boolean sum of the permitted permutation terms. For the
all-ones 2-by-2 matrix, the signed lifted sum is e: its two present
terms cancel. For a matrix with an entirely absent row, every
permutation term is tau and the sum is tau. This proves the precise
behavior rather than identifying term presence with nonzero signed
amplitude.

\textbf{Transfer:} parity, permanent and immanant tasks retain a free
module on the actual indexed combinatorial objects before the final
character/sign evaluation. G(R) records supported cancellation, while
the free source retains which objects contributed. The map on each basis
element is its original signed or character-weighted term. A
nonvanishing claim still requires evaluating that map or proving a
property of its image.

\section{2. Support diagrams and
reconstruction}\label{edition-1-section-3}

Use a join-semilattice L with bottom, R-modules V\_l and coherent linear
transports rho\_lk for l\textless=k. The total carrier is the disjoint
union of the V\_l. Its addition is

\begin{verbatim}
(l,x)+(k,y)=(l join k, rho_l,l-join-k(x)+rho_k,l-join-k(y)).
\end{verbatim}

The supported scalar r acts by \path|(l,x)->(l,rx)| and tau acts by the
bottom-fibre zero. The bottom fibre is not required to vanish. A
morphism is \texttt{(alpha,f\_l)} with alpha preserving bottom and joins
and with \path|f_k rho_lk = rho'_alpha(l),alpha(k) f_l|.

The inverse reconstruction in the source takes a G(R)-semimodule M to
\texttt{L=eM}, \path|V_l={m:em=l}| and \path|rho_lk(m)=m+k|. The
comparison \texttt{(l,m)-\textgreater{}m} has inverse
\texttt{m-\textgreater{}(em,m)}. Absorption \texttt{m+em=m} proves
additivity; the source gives the complete natural comparison, including
changes of support index. The local zero at l is l itself in the
reconstructed module.

\textbf{Transfer:} fix a tiling region and index the allowed move sets
by inclusion, a join-semilattice. Adding allowed moves is an explicit
inclusion on edge generators with identity on tiling generators. This
supplies a coherent diagram without assuming that every pair of
differently sized regions has a canonical extension map. Region
enlargement is a separate construction task. Congruence-depth quotients,
chain filtrations and labelled boundary-state systems likewise require
their own displayed transitions.

\section{3. Actual cohomology and the comparison
kernel}\label{edition-1-section-4}

For a fixed cochain map f:C-\textgreater D, write
\texttt{Z\_C=ker\ d\_C}, \path|B_C=im d_C(previous)| and

\begin{verbatim}
Krep={z in Z_C : f(z) is in B_D}.
\end{verbatim}

The map \path|Krep/B_C -> ker H(f)| sends the class of z to its original
source cohomology class. Its well-definedness follows from the fact that
f carries source boundaries into target boundaries. Its prequotient
kernel is exactly B\_C. Every class in ker H(f) has a cycle
representative in Krep. These three statements prove the quotient
isomorphism. An intertwining chain action preserves these submodules and
the same representative map intertwines the induced action.

At a support label, a boundary maps to its own supported zero. For
\path|F=T(alpha,f)|, the full inverse image of supported target zeros is
\path|T(L,ker f_l)|; the inverse image of the global zero restricts to
the labels with \path|alpha(l)=bottom|. Their comparison is the
inclusion of that smaller label set. This gives the exact relationship
of the two constructions without dropping killed classes at other
labels.

\textbf{Transfer:} knot-homology tasks keep the actual integral
differentials and comparison chain maps. Local-move tasks use the edge
differential \path|[t->t'] -> [t']-[t]| and augmentation
\path|epsilon([t])=1|. Reduced H0 is \path|ker epsilon / im d|.
Componentwise augmentation identifies H0 with one copy of R per
component: choose a vertex in each component for the inverse and use
path boundaries to relate every vertex to its chosen representative. The
reduced quotient has one relation that the component coefficients sum to
zero. A finite move graph with c components therefore has reduced rank
c-1. An all-region connectivity result still requires a proof uniform in
the region, with realizability retained.

\section{4. Filtered images, Rees defects and retained
jets}\label{edition-1-section-5}

The source starts with a filtered comparison c and its actual image I.
It keeps both filtrations \path|Q^a I=c(F^a V)| and
\path|P^a I=I intersect G^a W|. The inclusion of their Rees lattices
\path|L_Q subset L_P| gives the defect \texttt{D\_c=L\_P/L\_Q}. The
quotient measures the precise filtration discrepancy of that same map.
Its residue duality uses the coefficient of T\^{}(-1) in evaluation; it
does not supply a positive norm by itself.

An explicit fixture is \path|T^2 k[T] subset k[T]|. Its quotient basis
is \texttt{1,T}, Hilbert polynomial \texttt{1+z}, and
\path|(1-z)(1+z)=1-z^2|. The negative derivative of the numerator at z=1
is 2, its exact length. The dual residue pairing in the corresponding
monomial bases is the invertible anti-diagonal 2-by-2 matrix. The tests
check these identities.

The finite-jet source retains \texttt{R{[}t{]}/(h)} including repeated
roots and nilpotent classes. For ideal I in a polynomial algebra A, the
map

\begin{verbatim}
I/I^2 -> Omega_A tensor_A A/I, [f] -> df mod I
\end{verbatim}

is well-defined because \path|d(uv)=u dv+v du| has image zero modulo I
for u,v in I. Its kernel is retained. For \path|A=Q[x], I=(x)|,
\path|[ax+bx^2+cx^3]| maps to \texttt{a\ dx}. The original coefficients
are not set to convenient values before this map is constructed.

\textbf{Transfer:} characteristic-coefficient maps use every permitted
matrix entry, compute their actual polynomial image and differential,
then study explicit finite jets or relation modules. A local jet
statement is recorded at its base point and order. To obtain a global
coefficient-surjectivity conclusion, the global image needs its own
proof.

\section{5. Source-derived metrics and comparison
control}\label{edition-1-section-6}

The upstream workbench computes attainable metrics from actual source
corrections and distinguished orthogonal representatives. It keeps local
units, multiplicities and the original source measure in its
\path|Christoffel/Gram| construction. The non-zeta task must provide its
own source norm and observation map before invoking these constructions.

For the concrete map \texttt{q:Q\^{}2-\textgreater{}Q},
\texttt{q(x,y)=x+y}, with the Euclidean source inner product, the
section of 1 orthogonal to ker q is \texttt{(1/2,1/2)}. Its norm squared
is 1/2. Every other preimage is \path|(1/2+t,1/2-t)| with norm squared
\texttt{1/2+2t\^{}2}, proving the minimality directly. This is the Gram
induced by the actual map and norm.

\textbf{Transfer:} quantum frame problems retain vector coordinates and
the map sending the frame to its Gram matrix. Covariance or symmetry
restrictions come with explicit inclusions into that vector-coordinate
source. A certificate within one restricted image remains attached to
that inclusion; a full existence or nonexistence statement requires the
corresponding full-domain argument.

\section{6. Why this ranking follows the
maps}\label{edition-1-section-7}

A-ranked scopes already expose finite relation matrices, integral
complexes, quotient towers, local moves or polynomial maps to which the
displayed constructions can be applied immediately after source
certification. B scopes have a concrete source/target route but require
a new comparison, a substantial combinatorial construction or a
quantitative estimate. C scopes retain a named major unresolved bridge:
geometric realization, analytic compactness, integral comparison, an
inverse structural theorem or passage from finite stages to an infinite
endpoint.

This is a ranking of the written routes and available checks, not a
claim that an unlisted mathematical connection cannot exist. The exact
first task and outstanding full-problem obligation are recorded
separately for every candidate in BATCHES.md. Independent review may
move any record between tiers or remove a stale target.

\chapter{Initial exact interface proofs}\label{edition-2}

Source:
\path|sources/splitzero-nonzeta-workbench/splitzero-nonzeta-workbench/notes/verified-interfaces.md|.

SHA-256:
\path|97742cd29d95eb933df723d7d90d20a5bc3e3b4fa8e05ca5b3b24300ede93f8d|.

These are source-interface demonstrations and exact fixtures, not claims
of novelty or solutions to the listed open problems.
\path|tests/test_interfaces.py| contains their reproducible checks. The
broader source relationships are in METHODOLOGY.md.

\section{A. A triangle blocks one naive deletion/contraction
recipe}\label{edition-2-section-1}

For the triangle C3, delete an edge to obtain the three-vertex path, and
contract that edge to obtain two vertices joined by two parallel edges.
Their reduced Laplacians are respectively

\begin{verbatim}
[[2,-1],[-1,2]],  [[1,-1],[-1,2]],  [[2]].
\end{verbatim}

The respective cokernel orders are 3,1,2. More explicitly, Smith
reduction yields critical groups Z/3, the trivial group and Z/2. Every
group homomorphism Z/2-\textgreater Z/3 sends the generator to an
element x with 2x=0; the only such element is zero. Thus it cannot
supply an injective first arrow of a short exact sequence with these
groups. Their orders also fail the required multiplicative equality: 3
is not 2 times 1.

AIM chip-firing Question 10 asks about an algebra-level
\path|deletion/contraction| construction. The group calculation does not
dispose of that question. It tells the agent exactly which naive
group-level arrow fails.

There is a coefficient-sensitive test at the algebra level. Over Q, use
the augmentation

\begin{verbatim}
Q[C3]=Q[z]/(z^3-1) -> Q,   z -> 1.
\end{verbatim}

The Chinese-remainder map \path|[f] -> (f(1), [f] mod (z^2+z+1))| is an
algebra isomorphism onto \path|Q x Q[z]/(z^2+z+1)|: its factors are
coprime since the second evaluates to 3 at z=1, and the two residue
classes determine a unique class modulo their product. The augmentation
kernel is therefore isomorphic, as an algebra with its own identity, to
the quadratic field Q(zeta\_3).

On the contraction side, \path|Q[C2]=Q[u]/(u^2-1)| contains the nonzero,
nonidentity idempotent \texttt{(1+u)/2}. An injective algebra
homomorphism from Q{[}C2{]} onto that augmentation kernel would give a
nontrivial idempotent in a field, contradicting \texttt{a(a-1)=0}. This
rules out the *specified* exact sequence with this augmentation and
these rational group algebras. Over C the second factor splits into two
copies of C, so this particular idempotent obstruction changes through
the displayed scalar-extension map. An admissible version of the AIM
question must retain its intended coefficients and arrows.

The workbench uses this example to prevent replacing an actual
comparison kernel by a numerically attractive tree-count identity. It is
not advertised as a new solution to AIM Question 10.

\section{B. An exact deletion-observation
map}\label{edition-2-section-2}

Fix the binary word x=(x\_0,\ldots,x\_(n-1)) and retain each bit
independently with probability p, with 0\textless p\textless1. Its
coefficient polynomial is \path|P_x(z)=sum_i x_i z^i|. The mean
polynomial of the observed trace is

\begin{verbatim}
M_x(w)=p P_x(1-p+pw).
\end{verbatim}

For bit i to contribute, it survives with probability p.~Its output
index is the number of earlier surviving bits. Independence gives the
generating function \texttt{(1-p+pw)\^{}i} for that count. Multiplying
by p x\_i and summing proves the formula.

The coefficient matrix on polynomials of degree below n is triangular:
the coefficient of w\^{}i contributed by z\^{}i is p\^{}(i+1), and
higher output degrees do not occur. Its determinant is
\path|p^(n(n+1)/2)|, nonzero for this fixed p.~This proves finite
injectivity on coefficient vectors. The same determinant becomes small
as n grows; the calculation supplies no uniform statistical sample
bound. The sample-complexity task therefore retains the quantitative
comparison rather than ending at finite injectivity.

The tests enumerate every deletion mask for every binary word of length
at most 5 at p=1/2 and compare the exact expected polynomial with the
displayed formula. They also check the triangular determinant through
n=6.

\section{C. A realizable subsequence-deck kernel
element}\label{edition-2-section-3}

Let F be the free integer module on binary words of length 4, with basis
{[}x{]}, and let D\_2 map {[}x{]} to its multiplicity vector of length-2
subsequences indexed by 00,01,10,11. Subsequence positions are
increasing, not necessarily contiguous.

For 0110 the vector is \texttt{(1,2,2,1)}. For 1001 it is also
\texttt{(1,2,2,1)}. Direct enumeration of the six position pairs proves
both evaluations. Thus \path|[0110]-[1001]| is in ker D\_2 and comes
from two actual distinct source words. This is a concrete calibration
collision, not an all-n result on the k-deck reconstruction threshold.

General linear relations in ker D\_k remain indexed by their
coefficients. A task seeking a collision of two words must exhibit a
difference of two actual basis vectors, as this example does. The map
from source words into the free module is explicit and its image is
retained.

\section{D. The fourth moment retains the complete
autocorrelation}\label{edition-2-section-4}

For signs a\_0,\ldots,a\_(N-1), put \path|P(z)=sum_i a_i z^i| and
\path|c_k=sum_i a_i a_(i+k)| for k\textgreater=0. Direct multiplication
gives

\begin{verbatim}
P(z)P(z^-1)=N+sum_(k=1)^(N-1) c_k(z^k+z^-k).
\end{verbatim}

The constant term of its square is

\begin{verbatim}
N^2 + 2 sum_(k=1)^(N-1) c_k^2.
\end{verbatim}

Only exponents k and -k pair to give zero; this proves the identity
without discarding any lag. It is the exact algebraic entry for the
retained merit-factor variant. The tests check all sign vectors through
N=6. An asymptotic merit-factor conclusion requires estimates across the
original sequence family, not a change to independent correlation
coordinates.

\section{E. A killed class remains in its original
source}\label{edition-2-section-5}

Use a source window with degree-one space Q and zero adjacent maps. In
the target window, the incoming boundary is
\texttt{a-\textgreater{}(a,0)} in Q\^{}2 and the outgoing differential
is zero. The quotient map is \texttt{(a,b)-\textgreater{}b}.

The chain map \texttt{f(a)=(a,0)} is nonzero on representatives but zero
after that quotient; \texttt{g(a)=(0,a)} survives. The original source
class {[}1{]} is nonzero. The comparison kernel for f is all Q, and for
g it is zero. The identity-on-representatives map
\path|Krep/B_source -> ker H(f)| records this directly. This reproduces
the upstream example's maps, not only its dimensions.

\chapter{Reviewed A-tier continuation}\label{edition-3}

Source:
\path|sources/splitzero-nonzeta-reviewed-20260914/splitzero-nonzeta-reviewed-20260914/review/notes/proofs.md|.

SHA-256:
\path|9d81eb801f1f1adf1bfefec7678010632e6b251c0c2b8fdfec6b84c9b3831136|.

14 September 2026. These are written derivations and reproduced results,
with exact finite certificates. No new full resolution of P7, P15--P19,
Green 92 or AIM's coefficient-characterization question is asserted.
Published results used below retain their attribution. Independent
specialist review and upstream Lean reconstruction have not been
performed.

\section{1. Original benzel carriers and coefficient
convention}\label{edition-3-section-1}

The cell carrier is the integer triples v=(i,j,k) of sum one with all
three differences j-i, k-j, i-k in {[}1-a,b-1{]}, on the original domain
a,b\textgreater=2, a\textless=2b, b\textless=2a. Write e1,e2,e3 for the
standard basis. Right stones are u+\{e1,e2,e3\}, with sum(u)=0; left
stones are u-\{e1,e2,e3\}, with sum(u)=2. Bones are \{v-d,v,v+d\} for
d=e1-e2, e1-e3, e2-e3. These are Propp's five translated prototiles, not
the excluded bent trihex.

The map to axial coordinates is (i,j,k)-\textgreater(i,j), with inverse
(i,j)-\textgreater(i,j,1-i-j). It transforms the region differences into
j-i, 1-i-2j, 2i+j-1. For a left stone with base u, translate axial
coordinates by -(u\_i-1,u\_j); its offsets are \{(0,0),(1,0),(1,-1)\}. A
right stone has offsets \{(0,0),(1,0),(0,1)\} after translating by
-(u\_i,u\_j). Bones give straight triples in directions
(1,-1),(1,0),(0,1). This supplies both inverse carrier maps and the
explicit tile transport, including the literal carrier used in the P4
formalization.

Let R(T), L(T), B(T) count the right stones, left stones and bones of an
actual tiling. Use the stone-count invariant \path|Delta=R(T)-L(T).|
Propp's area-unit invariant C and this value satisfy C=3 Delta: each
stone consists of three cells. With s=(a+b) mod 3, the published formula
is

\begin{verbatim}
6 Delta = 3(a-b)^2-a-b                         (s=0)
        = -a^2+4ab-b^2-a-b+2                  (s=1)
        = 3(a-b)^2+a+b-2                      (s=2).
\end{verbatim}

The distinction of units is fixed by that multiplication map, not
silently discarded. Source: Defant--Li--Propp--Young,
arXiv:2209.05717v2, introduction and Theorem 1.1; Propp's problem paper,
Section 1.

\section{2. An explicit construction reproducing the resolved P14
family}\label{edition-3-section-2}

For every original admissible (a,b) with a+b=1 mod 3, there is exactly
one stones-and-bones tiling. This is the published DLPY Theorem 1.1. The
following direct residue construction supplies another fully specified
route to the same known endpoint.

Put A=1-a, B=b-1, and r=2-a mod 3. For each cell v, consider the three
possible anchors u=v-e\_j of right stones containing v. Their first
differences u\_j-u\_i occupy the three different residue classes modulo
3. Choose the unique anchor whose first difference is r.

All differences of a sum-zero anchor are congruent:
\path|(u_k-u_j)-(u_j-u_i)=-3u_j.| The three differences of the chosen
anchor lie in {[}A-1,B+1{]}, since the differences of v lie in
{[}A,B{]}. The endpoints A+1=2-a and B-1=b-2 are both congruent to r.
None of A-1,A,B,B+1 has that residue. Consequently every anchor
difference belongs to {[}A+1,B-1{]}. Across its right stone, each of the
three differences takes the three values delta(u)-1,delta(u),delta(u)+1.
The entire chosen stone therefore lies in the original benzel.

Every cell in this stone selects the same anchor: that anchor has
residue r and the selection is unique. Thus the stones partition the
cell carrier. More explicitly, for the selected anchor set U, the map

\begin{verbatim}
U x {1,2,3} -> V,  (u,j) -> u+e_j
\end{verbatim}

is bijective, with inverse obtained by the cell's unique selected anchor
and its offset index. This proves exhaustiveness and disjointness, not
only containment.

For uniqueness of a right-stone tiling of any finite region, choose a
cell of maximal first coordinate. In any right stone containing it, this
cell must be the e1-offset; the e2- or e3-offset would put another cell
one unit farther in the first coordinate. Its tile is therefore forced.
Removing that tile and repeating proves uniqueness by induction on the
number of cells.

Finally the published invariant and cell-area formulas give
\textbar V\textbar=3 Delta in this residue family. In any
stones-and-bones tiling,

\begin{verbatim}
|V|/3 = R+L+B = Delta = R-L,
\end{verbatim}

so 2L+B=0. Nonnegative integer tile counts force L=B=0. The partition
constructed above is consequently the unique tiling using all five
original prototiles. No extra symmetry or selected tiling family was
imposed.

The checker verifies the partition directly for all 150 ordered
admissible pairs with 2\textless=a,b\textless=30 and a+b=1 mod 3. This
computation accompanies the all-parameter proof; it is not its
replacement. The old author's-page open label is retained in provenance
but no longer controls the workbench's P14 routing.

\section{3. Exact residual scope for P7}\label{edition-3-section-3}

Retain the original request for a type-103 tiling whenever
Delta\textgreater=0. The involution (i,j,k)-\textgreater(i,k,j) negates
and permutes the three difference coordinates, so it transports the
(a,b) carrier to (b,a). It preserves the right/left stone sets and
permutes bone directions. It is its own inverse. We can therefore work
on a\textless=b with this specific reflection available to return every
construction.

The s=1 family is covered by Section 2. For s=2, generalized compression
Theorem 1.1 gives a positive product counting the original
\path|right-stone/two-bone| tilings, with

\begin{verbatim}
n=b+1-a, k=(2a-b-1)/3,
(a,b)=(n+3k,2n+3k-1).
\end{verbatim}

Both coordinate transformations are inverse; k\textgreater=0 on this
residue family. Each tiling counted there is also a type-103 tiling
under the literal inclusion of the permitted tile sets. No change to its
placed cells occurs. This disposes of that P7 subfamily using a
published result, without asserting a bijection with the larger tiling
set.

For s=0 define d=b-a and h=(2a-b)/3. The inverse is

\begin{verbatim}
a=d+3h, b=2d+3h, Delta=binom(d,2)-h.
\end{verbatim}

Substitution proves all three identities. Nonnegative invariant gives
\path|0<=h<=binom(d,2).| At h=0 the original boundary benzel is the same
cell set as (a,2a-2), already covered by the unique-right-stone family.
To check that cell equality directly, the largest difference of a cell
bounded below by 1-a is at most 2a-2 because the three differences sum
to zero. Equality 2a-2 forces the other two to be 1-a, whose
reconstruction has a coordinate with denominator three; thus that value
is unattainable for a sum-one integer cell. The actual upper bound is
2a-3, exactly the (a,2a-2) bound. At h=binom(d,2), the parameters are
\path|d(3d-1)/2,d(3d+1)/2,| the bone-only family of Kim--Propp
(arXiv:2206.04223).

The remaining new-construction scope after these published subfamilies
is therefore

\begin{verbatim}
d>=3, 1<=h<=binom(d,2)-1,
(a,b)=(d+3h,2d+3h).
\end{verbatim}

The h=1 part is also the P6 claim family: n=d+3 gives (a,b)=(n,2n-3). It
remains claim-audit material here rather than being removed on a claimed
Lean status alone. The checker supplies witnesses for all 38 unordered
admissible nonnegative-invariant pairs with
2\textless=a\textless=b\textless=12. That finite window includes
boundary and interior cases and is not a proof for the displayed
infinite residual family.

\section{4. Original weighted counts factor through finite
jets}\label{edition-3-section-4}

This calculation applies to the original all-five-tiles counts in
P15--P18. Let T be the entire finite tiling set of a specified benzel.
Put c=\textbar Delta\textbar{} and \path|j(T)=min(R(T),L(T)),| and
retain the integer polynomial

\begin{verbatim}
P(u)=sum_T u^j(T),  F(z)=sum_T z^(R(T)+L(T)).
\end{verbatim}

The integer identity R+L=\textbar R-L\textbar+2 min(R,L), with R-L=Delta
fixed, gives the exact equality

\begin{verbatim}
F(z)=z^c P(z^2).                                      (1)
\end{verbatim}

The original indexed map is Z{[}T{]}-\textgreater Z{[}u{]},
{[}T{]}-\textgreater u\^{}j(T). Its kernel consists of the integer
combinations with zero coefficient sum on every j-fibre. The map
P-\textgreater z\^{}c P(z\^{}2) is an injective Z-linear map with
inverse on its image obtained by reading the coefficients of
z\^{}(c+2j). Thus (1) retains the complete weighted polynomial, not just
its parity or its value at a selected weight.

Writing a\_j=\#\{T:j(T)=j\} and M\_m=sum\_j a\_j binom(j,m), the
binomial theorem proves the finite exact expansion

\begin{verbatim}
F(3)=3^c sum_(m>=0) 8^m M_m.                         (2)
\end{verbatim}

For precision r\textgreater=1, let N=ceil(r/3) and define the actual
ring maps

\begin{verbatim}
J_r: Z[u] -> (Z/2^r)[epsilon]/(epsilon^N), u->1+epsilon,
E_r: (Z/2^r)[epsilon]/(epsilon^N) -> Z/2^r, epsilon->8.
\end{verbatim}

E\_r is well-defined because 8\^{}N=0 modulo 2\^{}r. Its composite with
J\_r is evaluation at u=9 followed by reduction modulo 2\^{}r. Therefore

\begin{verbatim}
F(3) mod 2^r = 3^c E_r(J_r(P))
             = 3^c sum_(0<=m<N) 8^m M_m mod 2^r.    (3)
\end{verbatim}

The truncation uses the actual nilpotent epsilon and records its full
surviving coefficients. In particular F(3)=3\^{}c F(1) mod 8.
Equivalently, multiplying by the actual inverse of the odd unit 3\^{}c
in Z/8 gives 3\^{}(-c)F(3)=F(1) mod 8. This is a proved first-three-bits
comparison of the weighted and unweighted observations.

For the peripheral families, direct substitution gives

\begin{verbatim}
Delta(n,2n-3)=(n-5)(n-2)/2,
Delta(n,2n-4)=(n^2-7n+14)/2.
\end{verbatim}

The absolute value in c is retained, including n=4 in P16/P18. The new
assignments request the original minority-stone factorial moments M\_m
at the needed finite precision. Controlling those moments uniformly as n
varies remains mathematical work; (3) does not assert the requested
2-adic continuity.

Every ring map above lifts through G by G(f)(tau)=tau and
\path|G(f)(x-supported)=f(x)-supported.| Composition follows on these
two exhaustive cases. It sends supported zero to supported zero,
including after an evaluation vanishes. For a fixed maximum precision,
order the precision labels by reverse integer order and use the quotient
maps as transitions. This is a finite join-semilattice with bottom at
maximum precision and join=min. Its bottom fibre is allowed to be
nonzero, exactly as in the upstream support-diagram construction. No
unconstructed identification with the theta source is used.

There are 408 checked instances of (3): all 17 unordered admissible
benzel pairs with 2\textless=a\textless=b\textless=7 and precisions 1
through 24. The independent exact-cover enumeration also reproduces
their full polynomials, including the coefficient-by-coefficient
invariant parity constraint.

\section{5. P19: integral contraction certificates on the finite benzel
window}\label{edition-3-section-5}

Propp's full P19 concerns arbitrary simply-connected tileable regions.
Here the certified family is explicitly limited to the 17 benzels above,
with 6,492 tilings and 28,314 move edges in total. In particular the
(7,7) case has 5,766 tilings and 26,127 permitted edges.

For a fixed region, let C\_0 be the free Z-module on all its actual
tilings and C\_-1 the free Z-module on the allowed moves of Figure 5.
Send an oriented move {[}T-\textgreater U{]} to {[}U{]}-{[}T{]}, and let
epsilon:C\_0-\textgreater Z send every tiling to 1. Then epsilon
partial=0. We use exactly the two permitted two-tile changes: opposite
stones versus parallel bones; and a stone with a bone versus a stone of
the same orientation with a differently oriented bone, with the
identical six covered cells retained.

The emitted certificate lists every tiling, a root T0, and a parent for
every other tiling. The checker validates coverage and disjointness,
checks each parent edge against the actual two-move types, and verifies
that every parent chain reaches T0. Let h({[}T{]}) be the signed sum
along that chain from T0 to T, and let s(1)={[}T0{]}. Telescoping gives,
on every original generator,

\begin{verbatim}
partial h = id_(C_0)-s epsilon,  epsilon s=id_Z.
\end{verbatim}

For z in ker(epsilon) this becomes z=partial h(z); conversely every move
boundary has augmentation zero. This proves
\path|ker(epsilon)=im(partial)| over Z, so the reduced H\_0 vanishes for
every certified region. It retains explicit integer boundary
representatives, not merely a floating-point rank.

The checker constructs tiles in two ways: translated triples in the
barycentric carrier, and exhaustive classification of all three-cell
subsets against axial offsets. It also compares a memoized
fewest-options counting recurrence with a separate fixed-first-cell
enumeration. These are independently implemented finite checks, not
independent specialist review.

To attach the result to Split-Zero, let L be the power set of the
permitted edge set, ordered by inclusion and joined by union. At each
label A, retain the same C\_0, the edge submodule Z{[}A{]}, and
\path|H_A=C_0/im(partial_A).| For A subset B, identity on C\_0 and edge
inclusion commute with partial. The resulting map H\_A-\textgreater H\_B
sends a class to the class of the same representative, and

\begin{verbatim}
ker(H_A->H_B) = im(partial_B)/im(partial_A)
\end{verbatim}

by that representative map, whose kernel and surjectivity are immediate
from the two submodules. It is the upstream K\_rep/B\_C comparison with
all original boundaries retained. Reconstructing the diagram sends a
killed class to the zero at label B, not to absence. This is the
concrete method application. Extending these contractions to every
simply-connected region remains outside the finite certificate.

\section{6. Green 92: an exact original-deck map and a reproduced
collision family}\label{edition-3-section-6}

Let A=Z\textless X0,X1\textgreater{} be the noncommutative polynomial
algebra and I=(X0,X1). For a binary word w=w1\ldots wn define

\begin{verbatim}
S_w=(1+X_w1)...(1+X_wn).
\end{verbatim}

Choosing either 1 or the letter in each factor proves that the
coefficient of X\_v1\ldots X\_vj is the number binom(w,v) of
index-selected occurrences of v as a subsequence. Concatenation
satisfies S\_(uv)=S\_u S\_v. Linear extension from the free abelian
group on original words gives the coefficient observation map; modulo
I\^{}(k+1) it records all decks through length k.

For fixed n\textgreater=k and \textbar v\textbar=j\textless=k, count
pairs consisting of a selected occurrence of v and a k-index set
containing it. Counting first the occurrence and first the k-index set
gives respectively the right and left sides of

\begin{verbatim}
sum_(|u|=k) binom(w,u) binom(u,v)
   = binom(n-j,k-j) binom(w,v).                     (4)
\end{verbatim}

Thus the exact k-deck determines every smaller deck by division by the
displayed positive integer. That division is only used on the realizable
image, where (4) proves integrality. This explicitly relates Green's
original k-deck question to the truncated algebra; no arbitrary linear
kernel vector is called a collision of two source words.

For a concrete all-parameter family, define U0=0,V0=1 and

\begin{verbatim}
U_(k+1)=U_k V_k,  V_(k+1)=V_k U_k.
\end{verbatim}

These are the classical complementary Thue--Morse blocks; the k-binomial
literature already studies their equal decks (Lejeune--Leroy--Rigo,
arXiv:1812.07330). The following derivation is a reproduction, without a
novelty claim.

Put D\_k=S\_Uk-S\_Vk. Since D0=X0-X1 and

\begin{verbatim}
D_(k+1)=D_k(S_Vk-1)-(S_Vk-1)D_k,
\end{verbatim}

induction gives D\_k in I\^{}(k+1). The two distinct words of length
2\^{}k therefore have the same subsequence counts through length k. To
check the first surviving layer, for k\textgreater=1 let c\_k be the
coefficient of X0\^{}k X1 in D\_k. We have c1=1. The linear part of
S\_Vk-1 is 2\^{}(k-1)(X0+X1). The coefficient of the all-zero monomial
in D\_k vanishes because U\_k and V\_k have equal zero counts.
Consequently the displayed commutator gives

\begin{verbatim}
c_(k+1)=-2^(k-1)c_k,
c_k=(-1)^(k-1) 2^((k-1)(k-2)/2), which is nonzero.
\end{verbatim}

This proves that the difference dies under truncation through degree k
but survives in degree k+1, with its actual coefficient and source words
retained. The code replays k=1,\ldots,7 and performs 32,738 complete
word/deck comparisons with a second checker that enumerates index
subsets directly.

The exhaustive windows n\textless=12, k\textless=4 give first collision
lengths 2,4,7,12 respectively. Witnesses are

\begin{verbatim}
k=1: 01 / 10
k=2: 0110 / 1001
k=3: 0110001 / 1000110
k=4: 011000100110 / 100010110001.
\end{verbatim}

The scan includes all smaller eligible lengths, so the finite minimality
statement is certified by the declared enumeration. It is not a new
general bound for Green 92; the finite-jet difference and realizable
two-word certificates are reusable attack inputs.

\section{7. AIM coefficient maps: a complete small sign-chamber
construction}\label{edition-3-section-7}

Consider all real matrices {[}{[}a,b{]},{[}c,d{]}{]} with
a,b\textgreater0 and c,d\textless0. For each such matrix use
S=diag(1,1/b). Direct multiplication gives

\begin{verbatim}
S^(-1) A S = [[a,1],[z,d]], z=bc<0.
\end{verbatim}

The full coordinate map is (a,b,c,d)-\textgreater(a,d,z,b); its inverse
is \path|(a,d,z,b)->(a,b,z/b,d),| with a,b\textgreater0 and
d,z\textless0. Thus b has been retained as an invertible coordinate,
rather than an unjustified entry replacement.

The characteristic-coefficient map in the displayed chart is

\begin{verbatim}
Phi(a,d,z)=(-a-d,ad-z).
\end{verbatim}

For every desired pair (u,v) in R\^{}2 take
a=1+\textbar u\textbar+\textbar v\textbar, d=-u-a, z=ad-v. Then
a\textgreater0 and a+u\textgreater=1+\textbar v\textbar, so d\textless0
and a(a+u)\textgreater\textbar v\textbar. Hence z=-a(a+u)-v\textless0.
Direct substitution yields \path|Phi(a,d,z)=(u,v).| This is a complete
section proving spectral arbitrariness of this particular 2x2 full sign
pattern. It is a control instance, not a characterization of all
patterns in AIM Section 2 Q1.

The Jacobian minor in (d,z) is {[}{[}-1,0{]},{[}a,-1{]}{]} with
determinant 1. At \path|(a,d,z)=(1,-1,-1),| the actual matrix squares to
zero. These observations specify finite-jet inputs at the original
nilpotent point. The checker verifies 441 integer target pairs and three
rational b-values for each.

\section{8. Source-resolution controls and review
boundary}\label{edition-3-section-8}

AIM chip-firing Q14 is addressed by Perkinson--Yang--Yu,
arXiv:1309.2201v2, Theorem 3, with inverse Algorithms 1--2 and the
\path|degree/kappa-inversion| identity. Our exact reproduction verifies
both inverse laws against independently enumerated parking functions and
spanning trees: K2 through K6 and all 38 connected labelled four-vertex
graphs, 43 graph cases and 1,569 \path|graph/parking-function|
instances. Root 0 and descending neighbor order are fixed, as in the
source. This is a replay of published mathematics.

AIM matrix-spectrum Section 3 Q6 is the Grone--Merris assertion proved
by Hua Bai, arXiv:0911.2172v4. The status update uses that matching
primary theorem. The quick checks construct the original incidence
matrix B, retain L=BB\^{}T, verify exact characteristic polynomials for
29 \path|complete/complete-bipartite| graph cases, and test the
resulting integer majorizations. They do not audit Bai's full proof.

Shitov's arXiv:1612.01783v2 supplies a counterexample to the complex
zero-pattern 2n analogue: order 708, 1,415 nonzeros. The workbench
records this precise related result, not a disposition of the real
sign-pattern or arbitrary-positive-characteristic questions. We
reproduced the paper's two explicit 8x8 matrices, with characteristic
polynomials t\^{}8 and (t-1)\^{}8, using exact trace recurrences and an
independent rational-elimination determinant check at nine values of t.
Agreement at nine points certifies equality of the degree-eight
polynomials. This checks those two matrices, not the whole
708-dimensional generic-selection proof. The source's 8x8 block itself
is not claimed spectrally arbitrary: its coefficient identity is
phi4=-x4 phi7, as direct determinant expansion verifies.

P4 and P6 have pinned public formal sources, not merely an index
mention. The P4 carrier, invariant numerator and publication endpoint
were inspected. Its integer invariant numerator is exactly 6 Delta, and
its target tile carrier matches Section 1. Nine finite P4 witnesses and
seven P6 formula values were independently recalculated, the latter for
n=5 through 11. The exact P4 commit is
f8f91033216ba16fa851adbb9d6a9cf005509619; P6 is
5d1983de4592261dd0b915987ebcc9c0c8ae11b8. No Lean executable is
installed in this runtime. Neither the full source dependency closure
nor the trust-zero build was replayed. Both entries retain
\path|CLAIMED_UNVERIFIED| and request that specific remaining review.

\section{Sources and reproduction}\label{edition-3-section-9}

Primary problem sources and theorem locators are in
\path|../catalog/status-delta.json.| Upstream Split-Zero support
diagrams are pinned at
\href{mailto:zeta-function-research-reader@1c8ec52c85c173adc9f8403a8a26914955e2f5a9}{\path|zeta-function-research-reader@1c8ec52c85c173adc9f8403a8a26914955e2f5a9|},
\path|workbenches/splitzero-tandem/tex/support_diagrams.tex,| especially
D1--D8.

Run from this continuation's root:

\begin{verbatim}
python tools/check_progress.py
python -O tools/check_progress.py
\end{verbatim}

The full certificate and receipt files are generated under evidence/. No
network, paid computation, downstream agent launch or public-resolution
announcement occurs. A finite receipt cannot promote a literature claim
or our written argument to independently reviewed status.

\chapter{Turn 1: first and second collars}\label{edition-4}

Source: \path|sources/benzel-p7-turn1/benzel-p7/PROOF.md|.

SHA-256:
\path|c0655d3122ce84f3e6cebcd417ba6ca17fdcbc0ec5b908246ef9cff07c1a58b7|.

\textbf{Sprint turn 1 --- 15 September 2026.} Written proof and
executable finite replay. This does not resolve the full Problem 7. The
all-parameter construction below is derived in this continuation; its
priority relative to every earlier special-family construction has not
been established. In particular, the first collar overlaps the P6 family
and carries no claim of first discovery. No independent specialist
review or Lean build is claimed.

\section{1. Target, source coordinates, and exact
scope}\label{edition-4-section-1}

Propp's Problem 7 asks for tilings by right stones and all three bone
orientations throughout the original admissible domain with nonnegative
Conway--Lagarias invariant. Use its integer cells \texttt{(i,j,k)} with
\texttt{i+j+k=1} and all three differences in \texttt{{[}1-a,b-1{]}}.
The axial map and inverse are

\begin{verbatim}
(i,j,k) -> (i,j),        (i,j) -> (i,j,1-i-j).
\end{verbatim}

Write

\begin{verbatim}
D0=j-i, D1=1-i-2j, D2=2i+j-1,
V(a,b)={(i,j):1-a<=D0,D1,D2<=b-1}.
\end{verbatim}

The original right stone at axial anchor \texttt{(i,j)} has cells

\begin{verbatim}
S(i,j)={(i,j),(i+1,j),(i,j+1)}.
\end{verbatim}

Let \texttt{D(i,j)}, \texttt{H(i,j)}, \texttt{V(i,j)} denote
respectively the three bones with cells

\begin{verbatim}
D: (i+u,j-u), H: (i+u,j), V: (i,j+u), u=0,1,2.
\end{verbatim}

Thus every tile below is an actual translated original prototile. In
tile tables \texttt{V} denotes a vertical bone; \texttt{V(a,b)} with
parameters denotes the region.

Put

\begin{verbatim}
W_h(d)=V(d+3h,2d+3h),  d>=3, h>=0.
\end{verbatim}

The source cell-count formula (DLPY, equation (1)) gives

\begin{verbatim}
|W_h(d)|=3 binom(d,2)+(9d-3)h+9h^2.                 (1)
\end{verbatim}

The source invariant in stone-count units is

\begin{verbatim}
Delta(W_h)=binom(d,2)-h.                             (2)
\end{verbatim}

These are literal substitutions \path|a=d+3h,b=2d+3h|, with inverse
\path|d=b-a,h=(2a-b)/3| on the residue-zero family. No region is changed
by this parameter map.

\textbf{Constructed result.} For every integer
\texttt{d\textgreater{}=3}, the regions \texttt{W\_1(d)} and
\texttt{W\_2(d)} have explicitly specified type-103 tilings. Their
numbers of right stones and bones are, respectively,

\begin{verbatim}
W_1: binom(d,2)-1 right stones, 3d+3 bones;
W_2: binom(d,2)-2 right stones, 6d+12 bones.          (3)
\end{verbatim}

In original parameters the second family is \path|(a,b)=(n,2n-6)|,
\texttt{n\textgreater{}=9}, and its reflected family. This is an
infinite subfamily of the original target, not a renamed replacement for
Problem 7.

\section{2. Base tiling and its two reserved
stones}\label{edition-4-section-2}

For integers \texttt{r,s\textgreater{}=0}, \texttt{r+s\textless{}=d-2},
place the right stone with anchor

\begin{verbatim}
q(r,s)=(d-2-2r-s,r-s).
\end{verbatim}

The associated sum-zero barycentric anchor has third coordinate
\texttt{2-d+r+2s}. Its difference coordinates are

\begin{verbatim}
2-d+3r, 2-d+3s, 2-d+3(d-2-r-s).
\end{verbatim}

Each lies in \texttt{{[}2-d,2d-4{]}}. A right stone varies each
difference by \texttt{-1,0,1}, so each placed cell belongs to
\texttt{W\_0(d)}.

For any cell there are exactly three possible right-stone anchors
containing it. Their first differences occupy the three distinct
residues modulo three. Every displayed anchor has first difference
\texttt{2-d} modulo three, so two displayed stones cannot share a cell.
There are \texttt{binom(d,2)} anchors. Their disjoint cells have
cardinality \path!3 binom(d,2)=|W_0|! by (1), proving coverage.

Call this tiling \texttt{T\_0}. It contains the two distinct stones

\begin{verbatim}
Q0=S(0,2-d)       (r=0,s=d-2),
Q1=S(2-d,d-2)     (r=d-2,s=0).                       (4)
\end{verbatim}

They are distinct for \texttt{d\textgreater{}=3}.

\section{3. First collar: an explicit positive solution of the residual
equation}\label{edition-4-section-3}

Translate \texttt{W\_0} and \texttt{T\_0} by

\begin{verbatim}
t1=(-2,1).
\end{verbatim}

The difference-coordinate increments are \texttt{(3,0,-3)}. Consequently
the translated inner region \texttt{I1=t1+W\_0} has bounds

\begin{verbatim}
D0 in [4-d,2d+2],
D1 in [1-d,2d-1],
D2 in [-d-2,2d-4],                                 (5)
\end{verbatim}

and is contained in \texttt{W\_1}, whose three bounds are
\texttt{{[}-d-2,2d+2{]}}. The translated \texttt{Q0} is

\begin{verbatim}
S1=S(-2,3-d).
\end{verbatim}

Here is the complete bone family \texttt{C1}:

\begin{longtable}[]{@{}llll@{}}
\toprule\noalign{}
Label & Kind & Anchor & Index range \\
\midrule\noalign{}
\endhead
\bottomrule\noalign{}
\endlastfoot
A & D & \texttt{(-2,3-d)} & one tile \\
B & D & \texttt{(-2,4-d)} & one tile \\
C & D & \texttt{(-1,1-d)} & one tile \\
F\_r & H & \path|(-d-1+2r,d+1-r)| &
\texttt{0\textless{}=r\textless{}=d-1} \\
G\_r & H & \texttt{(r,3-d+r)} &
\texttt{0\textless{}=r\textless{}=d-2} \\
E & V & \texttt{(1,-d)} & one tile \\
J\_r & V & \texttt{(r+2,r-d)} &
\texttt{0\textless{}=r\textless{}=d-1} \\
\end{longtable}

\subsection{Containment and the exact intersection with the inner
source}\label{edition-4-section-4}

Substitution gives the following difference coordinates for the cell
with offset \texttt{u} in each row:

\begin{longtable}[]{@{}llll@{}}
\toprule\noalign{}
Label & D0 & D1 & D2 \\
\midrule\noalign{}
\endhead
\bottomrule\noalign{}
\endlastfoot
A & \texttt{5-d-2u} & \texttt{2d-3+u} & \texttt{-d-2+u} \\
B & \texttt{6-d-2u} & \texttt{2d-5+u} & \texttt{-d-1+u} \\
C & \texttt{2-d-2u} & \texttt{2d+u} & \texttt{-d-2+u} \\
F\_r & \texttt{2d+2-3r-u} & \texttt{-d-u} & \texttt{-d-2+3r+2u} \\
G\_r & \texttt{3-d-u} & \texttt{2d-5-3r-u} & \texttt{2-d+3r+2u} \\
E & \texttt{-d-1+u} & \texttt{2d-2u} & \texttt{1-d+u} \\
J\_r & \texttt{-d-2+u} & \texttt{2d-1-3r-2u} & \texttt{3r+3-d+u} \\
\end{longtable}

Every entry is in \texttt{{[}-d-2,2d+2{]}}: it is affine in
\texttt{r,u}, so the minimum and maximum occur at the displayed index
endpoints and \texttt{u=0,2}; substitution with
\texttt{d\textgreater{}=3} gives the bounds. This proves outer
containment for every integer parameter.

Against (5), all cells of C, G, E and J violate the lower D0 bound; all
cells of F violate the lower D1 bound. Row A has only its \texttt{u=0}
cell inside I1. Row B has only its \texttt{u=0,1} cells inside I1. Those
three cells are exactly S1. Thus

\begin{verbatim}
covered(C1) subset (W_1\I1) union S1.               (6)
\end{verbatim}

\subsection{Disjointness}\label{edition-4-section-5}

The three D-bones have distinct \texttt{k=1-i-j} coordinates
(\texttt{d,d-1,d+1}). H-bones F have distinct rows
\texttt{j\textgreater{}=2}; H-bones G have distinct rows
\texttt{j\textless{}=1}, so all H-bones are mutually disjoint. V-bones
occupy distinct columns \texttt{i=1} and \texttt{i=2,...,d+1}.

D versus F: every D cell has \texttt{j\textless{}=1}, whereas F has
\texttt{j\textgreater{}=2}. D versus G: G has \texttt{i\textgreater{}=0}
and \texttt{j\textgreater{}=3-d}. The possible D cells at
\texttt{i\textgreater{}=0} have \texttt{j\textless{}3-d}. D versus V:
rows A and B have \texttt{i\textless{}=0}; row C has only one cell with
\texttt{i=1}, at \texttt{j=-d-1}, below E. The other V columns begin at
\texttt{i=2}.

E versus H: E has \texttt{j\textless{}=2-d}, below G and F. A cell of J
in column \texttt{i} has \texttt{j\textless{}=i-d}. A G-cell in that
column satisfies \texttt{r\textgreater{}=i-2}, hence
\path|j=3-d+r>=i+1-d|. An F-cell has \texttt{j\textgreater{}=2} and
\texttt{i\textless{}=d+3-2j}; combining this with
\texttt{j\textless{}=i-d} would give \texttt{j\textless{}=1}, a
contradiction. This checks every pair of orientation families.

There are \path|3+d+(d-1)+1+d=3d+3| bones. Their disjoint cells number
\texttt{9d+9}. Equation (1) gives

\begin{verbatim}
|W_1|-|W_0|+|S1|=9d+6+3=9d+9.
\end{verbatim}

Together with (6) this proves the exact partition

\begin{verbatim}
covered(C1)=(W_1\I1) disjoint-union S1.              (7)
\end{verbatim}

Remove S1 from the translated source tiling and insert C1. The result

\begin{verbatim}
T1=(t1+T0)\{S1} union C1                            (8)
\end{verbatim}

is a genuine type-103 tiling of W\_1. The other reserved source stone Q1
survives and now has anchor \texttt{(-d,d-1)}.

\section{4. Second collar: another exact positive
lift}\label{edition-4-section-6}

Translate T1 and W\_1 by

\begin{verbatim}
t2=(1,1).
\end{verbatim}

The difference-coordinate increments are \texttt{(0,-3,3)}. The
translated inner region \texttt{I2=t2+W\_1} is contained in W\_2 and has
bounds

\begin{verbatim}
D0 in [-d-2,2d+2],
D1 in [-d-5,2d-1],
D2 in [1-d,2d+5].                                  (9)
\end{verbatim}

The surviving reserved stone becomes

\begin{verbatim}
S2=S(1-d,d).
\end{verbatim}

Use this complete bone family C2:

\begin{longtable}[]{@{}llll@{}}
\toprule\noalign{}
Label & Kind & Anchor & Index range \\
\midrule\noalign{}
\endhead
\bottomrule\noalign{}
\endlastfoot
A & D & \texttt{(-d-2,d)} & one tile \\
B & D & \texttt{(-d-2,d+1)} & one tile \\
C\_r & D & \path|(-d-1+r,d-2-2r)| &
\texttt{0\textless{}=r\textless{}=d-1} \\
F & H & \texttt{(-d-3,d+2)} & one tile \\
G & H & \texttt{(-d-2,d+3)} & one tile \\
H & H & \texttt{(-d-1,d+1)} & one tile \\
J & H & \texttt{(-d,d)} & one tile \\
K\_r & V & \path|(1-d+r,d-3-2r)| &
\texttt{0\textless{}=r\textless{}=d-1} \\
L & V & \texttt{(1,-d-2)} & one tile \\
M & V & \texttt{(2,-d-3)} & one tile \\
N\_r & V & \path|(3+r,-d-2+r)| & \texttt{0\textless{}=r\textless{}=d} \\
\end{longtable}

\subsection{Containment and inner
intersection}\label{edition-4-section-7}

\begin{longtable}[]{@{}llll@{}}
\toprule\noalign{}
Label & D0 & D1 & D2 \\
\midrule\noalign{}
\endhead
\bottomrule\noalign{}
\endlastfoot
A & \texttt{2d+2-2u} & \texttt{3-d+u} & \texttt{-d-5+u} \\
B & \texttt{2d+3-2u} & \texttt{1-d+u} & \texttt{-d-4+u} \\
C\_r & \texttt{2d-1-3r-2u} & \texttt{6-d+3r+u} & \texttt{-d-5+u} \\
F & \texttt{2d+5-u} & \texttt{-d-u} & \texttt{-d-5+2u} \\
G & \texttt{2d+5-u} & \texttt{-d-3-u} & \texttt{-d-2+2u} \\
H & \texttt{2d+2-u} & \texttt{-d-u} & \texttt{-d-2+2u} \\
J & \texttt{2d-u} & \texttt{1-d-u} & \texttt{-d-1+2u} \\
K\_r & \texttt{2d-4-3r+u} & \texttt{6-d+3r-2u} & \texttt{-d-2+u} \\
L & \texttt{-d-3+u} & \texttt{2d+4-2u} & \texttt{-d-1+u} \\
M & \texttt{-d-5+u} & \texttt{2d+5-2u} & \texttt{-d+u} \\
N\_r & \texttt{-d-5+u} & \texttt{2d+2-3r-2u} & \texttt{3-d+3r+u} \\
\end{longtable}

The endpoint substitution used above proves that every entry belongs to
W\_2's bounds \texttt{{[}-d-5,2d+5{]}}.

Against (9), A, B, C and K are below the D2 lower bound; F and G exceed
the D0 upper bound; L exceeds the D1 upper bound; M and N are below the
D0 lower bound. H has only its \texttt{u=2} cell in I2. J has only its
\texttt{u=1,2} cells in I2. The three exceptional cells are exactly S2.
Thus

\begin{verbatim}
covered(C2) subset (W_2\I2) union S2.               (10)
\end{verbatim}

\subsection{Disjointness}\label{edition-4-section-8}

The D-bones have distinct \texttt{k} levels \path|3,2,4,...,d+3|. The
four H-bones have distinct rows \path|d+2,d+3,d+1,d|. The V-bones have
distinct columns \path|1-d,...,0,1,2,3,...,d+3|.

D versus H: H's rows \texttt{d+2,d+3} exceed all D rows. In row
\texttt{d+1}, D contributes only the point with column \texttt{-d-2}, to
the left of H's columns. In row \texttt{d}, D contributes only columns
\texttt{-d-2,-d-1}, to the left of J's columns starting at \texttt{-d}.
The indexed D-family has rows at most \texttt{d-2}.

H versus V: K has row at most \texttt{d-1}, below all H rows. L, M and N
have columns at least 1; all H cells have column at most
\texttt{2-d\textless{}=-1}.

D versus V: A and B have columns at most \texttt{-d}, to the left of K's
first column \texttt{1-d}. C has columns at most 0, to the left of L, M
and N. For a possible C\_r/K\_s collision at offsets u,v, equality of
first coordinates gives \texttt{s=r+u-2}; equality of second coordinates
then gives \texttt{u-v=3}. This is impossible for
\texttt{0\textless{}=u,v\textless{}=2}. These arguments exhaust all
pairs.

There are \path|2+d+4+d+2+(d+1)=3d+9| bones, covering \texttt{9d+27}
disjoint cells. Equation (1) gives

\begin{verbatim}
|W_2|-|W_1|+|S2|=9d+24+3=9d+27.
\end{verbatim}

Together with (10), this proves

\begin{verbatim}
covered(C2)=(W_2\I2) disjoint-union S2.             (11)
\end{verbatim}

Therefore

\begin{verbatim}
T2=(t2+T1)\{S2} union C2                            (12)
\end{verbatim}

tiles W\_2 by precisely the original right stones and bones. No tiling
search or unproved induction step is used in (8) or (12). Equations (3)
follow by counting the two consumed stones and the two added bone
families.

The reflection \path|(i,j,k)->(i,k,j)|, or axially
\path|(i,j)->(i,1-i-j)|, is its own inverse, sends V(a,b) to V(b,a),
preserves both stone orientations and permutes bone directions. Applying
it to every placed tile gives the reflected family with the same
certificate.

\section{5. A finite-jet obstruction and its explicit
cancellation}\label{edition-4-section-9}

The following calculation identifies the exact class canceled by the
collar lift. It keeps the original lattice translations and gives a
nonzero cohomology observation for an unmodified collar.

Let

\begin{verbatim}
A=F_3[e,f]/(e^2,ef,f^2),
phi: Z[X^+-1,Y^+-1] -> A,
X -> 1+e, Y -> 1+f.
\end{verbatim}

The displayed images are units, with inverses \texttt{1-e,1-f}, so this
is a defined ring homomorphism. For every integer i,j,

\begin{verbatim}
phi(X^i Y^j)=1+i e+j f.                             (13)
\end{verbatim}

Map a cell generator \texttt{(i,j)} to (13) and extend linearly. A bone
with direction \texttt{(p,q)} maps to the sum over u=0,1,2:

\begin{verbatim}
3 + (3i+3p)e + (3j+3q)f = 0.
\end{verbatim}

A right stone maps to \texttt{e+f}; a left stone maps to
\texttt{2e-f=-(e+f)}. The class \texttt{e+f} is nonzero, since
\texttt{1,e,f} is a basis of A. Multiplication by any translation unit
\texttt{1+i\ e+j\ f} fixes \texttt{e+f}.

For any W\_h, cyclic permutation of the barycentric coordinates permutes
the finite cell set. Thus the three coordinate sums are equal, and each
is \texttt{\textbar{}W\_h\textbar{}/3}. Its cell observation is
consequently

\begin{verbatim}
phi(1_W_h)=(|W_h|/3)(e+f)
          =(binom(d,2)-h)(e+f),                   (14)
\end{verbatim}

where (1) gives the second equality modulo three. Translating W\_h
leaves this observation unchanged because its cell count is divisible by
three.

For the raw residual \path|r_h=1_W_(h+1)-J 1_W_h| of a containing
translation J, (14) gives

\begin{verbatim}
phi(r_h)=-(e+f) != 0.                              (15)
\end{verbatim}

Every integral sum of bone boundaries has phi-value zero. Equation (15)
therefore detects an actual obstruction to filling that collar with
bones while leaving the entire source tiling untouched. Adding the
boundary of one actual right stone cancels this observation. For h=0 and
h=1, the tables prove the stronger integral identities

\begin{verbatim}
partial beta_1 = 1_W1 - J1 1_W0 + partial S1,
partial beta_2 = 1_W2 - J2 1_W1 + partial S2.         (16)
\end{verbatim}

Here beta\_i is the sum of the original placed bone generators, each
with coefficient +1. The positive coefficients and the exact partition
checks establish the actual tiling lifts. Equation (15) by itself
supplies no such positive preimage.

\section{6. The exact Split-Zero cohomology
interface}\label{edition-4-section-10}

For a finite cell support lambda let C\_B\^{}0(lambda) be the free
Z-module on placed bones contained in lambda; let C\_103\^{}0(lambda)
additionally include the placed right stones; let C\^{}1(lambda) be the
free Z-module on its cells. Define the actual two-term differentials by

\begin{verbatim}
partial[tile]=sum_(cell in tile)[cell].
\end{verbatim}

Finite supports, ordered by inclusion and joined by union, form a
semilattice with bottom the empty set. Inclusions of actual tile and
cell generators give coherent maps of these complexes. Translation sends
every original cell and tile to its translate and commutes with partial
on generators; its inverse is translation by the negative vector.

The inclusion from the bone complex to the type-103 complex is identity
in degree one. On first cohomology its kernel has the explicit
description

\begin{verbatim}
im(partial_103)/im(partial_B)
    -> ker(C^1/im(partial_B) -> C^1/im(partial_103)),
[z] -> [z].                                       (17)
\end{verbatim}

The forward map is defined because im(partial\_B) is a submodule of
\path|im(partial_103).| Its prequotient kernel is exactly
im(partial\_B); every killed class has a representative in
\path|im(partial_103).| These facts give the inverse and prove (17).

The cell observation phi kills bone boundaries, hence descends to the
bone cohomology. Equations (15)-(16) retain the collar class and its
opposite stone class there before displaying their bone-boundary lift.
This is the particular comparison-kernel computation used here.

Apply the pinned upstream reconstruction to the full support diagrams: a
reconstructed class is \path|(lambda,[z])|, a transition preserves its
actual representative, the supported-zero scalar sends it to
\texttt{(lambda,0)}, and the absent scalar sends it to the bottom zero.
The natural reconstruction maps and inverse are the upstream D2-D5 maps.
The quotient uses the D6 coequalizer, and (17) is D7 specialized to
these displayed complexes. The observation target may be the constant
A-diagram on the same support labels; its bottom fibre is allowed by the
source construction.

No infinite theta-source estimate or arithmetic weight bound is
imported. The argument uses these explicit support, quotient, jet and
positive-boundary maps. Both the small jet observation and the positive
lift are recorded, so an agent cannot treat disappearance of the jet
obstruction as an automatically supplied tiling.

\section{7. Scope of the next step}\label{edition-4-section-11}

The full unresolved construction range after the previously
source-matched families is \path|d>=3,1<=h<=binom(d,2)-1|. This note
constructs h=1,2. The remaining interior values begin at h=3 where
present; the invariant-zero endpoint retains its published Kim--Propp
construction.

The exact fixed-source test at d=3,h=2 attempts to retain every old
bone, remove only the remaining right stone and fill the next collar
with bones. For every one of the 13 containing integer translations
there is a cell of that patch belonging to no permitted bone contained
in the patch. The code records such cells directly. Thus this particular
frozen-source recurrence stops, including at the already tileable
invariant-zero endpoint. It neither refutes P7 nor rules out a lift
after moving old bones.

The next research task is to construct a support-changing boundary
repair with old bone generators retained in the comparison, then obtain
a repeatable positive lift through the full h-range. No universal repair
or remaining induction is assumed in this note.

\section{Sources and verification}\label{edition-4-section-12}

\begin{itemize}
\tightlist
\item
  James Propp, *Trimer Covers in the Triangular Grid: Twenty Mostly Open
  Problems*, Sections 1 and 6, Problem 7:
  \href{https://www.samuelfhopkins.com/OPAC/files/proceedings/propp.pdf}{\path|https://www.samuelfhopkins.com/OPAC/files/proceedings/propp.pdf|}
\item
  Author status page rechecked 2026-09-15; still labels P7 open, with a
  stated June-2025 status horizon:
  \href{https://faculty.uml.edu/jpropp/benzels.html}{\path|https://faculty.uml.edu/jpropp/benzels.html|}
\item
  Defant, Li, Propp, Young, *Tilings of Benzels via the Abacus
  Bijection*, arXiv:2209.05717v2, equation (1), invariant formula and
  Theorem 1.1:
  \href{https://arxiv.org/html/2209.05717v2}{\path|https://arxiv.org/html/2209.05717v2|}
\item
  The preceding workbench continuation, especially the residual
  parameter maps:
  \href{https://github.com/KokunoYumeto/mathematics-commons-pilot/blob/ea0cf71f7af213f6e25dd0669e2e67cc024bd0dc/workbenches/splitzero-nonzeta/review/notes/proofs.md}{\path|https://github.com/KokunoYumeto/mathematics-commons-pilot/blob/ea0cf71f7af213f6e25dd0669e2e67cc024bd0dc/workbenches/splitzero-nonzeta/review/notes/proofs.md|}
\item
  Split-Zero support reconstruction and cohomology, upstream commit
  \path|1c8ec52c85c173adc9f8403a8a26914955e2f5a9|,
  \path|workbenches/splitzero-tandem/tex/support_diagrams.tex|, D1-D8:
  \href{https://github.com/KokunoYumeto/zeta-function-research-reader/blob/1c8ec52c85c173adc9f8403a8a26914955e2f5a9/workbenches/splitzero-tandem/tex/support_diagrams.tex}{\path|https://github.com/KokunoYumeto/zeta-function-research-reader/blob/1c8ec52c85c173adc9f8403a8a26914955e2f5a9/workbenches/splitzero-tandem/tex/support_diagrams.tex|}
\end{itemize}

\path|python check.py| and \path|python -O check.py| replay explicit
tilings, independent cell-set definitions, shapes, incidence identities,
invariant counts and the fixed-source obstruction. They are bounded
implementation tests accompanying the written all-parameter proof, not a
replacement for it. Third-party proof manuscripts or formalizations are
not bundled.

\chapter{Turn 2: concurrent third-collar edition}\label{edition-5}

Source: \path|supplements/benzel-p7-turn2-package/turn2/PROOF.md|.

SHA-256:
\path|53063a773775affef03e85f0ddde1185a8a7e5d9e5edce2b1bf22d9caea48560|.

\textbf{Focused Propp-7 attempt, turn 2, 15 September 2026.} Parent:
\path|90f23b04fbde1e25d07342e2b8a8c80d919402c5|. The constructions in
§§1--4 and algebraic identifications in §§5--7 have complete written
proofs below. Independent specialist review and Lean replay have not
been performed. No full solution of Problem 7, first-discovery claim, or
empirical probability of completion is asserted.

\section{1. Original objects and constructed
result}\label{edition-5-section-1}

Keep exactly the original axial cells and tiles from
\href{../PROOF.md}{the turn-1 proof}:

\begin{itemize}
\tightlist
\item
  \texttt{D0=j-i}, \texttt{D1=1-i-2j}, \texttt{D2=2i+j-1};
\item
  \path|W_h(d)={(i,j):1-d-3h <= D0,D1,D2 <= 2d+3h-1}|;
\item
  \path|S(x,y)={(x,y),(x+1,y),(x,y+1)}|;
\item
  bones \path|D(x,y)={(x+u,y-u):u=0,1,2}|,
  \path|H(x,y)={(x+u,y):u=0,1,2}|, and \path|V(x,y)={(x,y+u):u=0,1,2}|.
\end{itemize}

These are the actual type-103 prototiles. The inverse of the axial cell
map is \path|(i,j)->(i,j,1-i-j)|. The parameter map to original benzels
is

\begin{verbatim}
(d,h) -> (a,b)=(d+3h,2d+3h),
d=b-a, h=(2a-b)/3.
\end{verbatim}

For every integer \texttt{d\textgreater{}=3}, we construct a type-103
tiling of \texttt{W\_3(d)}. It has

\begin{verbatim}
binom(d,2)-3 right stones and 9d+27 bones.
\end{verbatim}

Equivalently this tiles every original \texttt{(n,2n-9)}-benzel for
integers \texttt{n\textgreater{}=12}, and its reflected family. The
constructor is given by formulas; its final execution performs no
search. The \texttt{d=3} member is the known bone-only \texttt{(12,15)}
endpoint, retained as a control rather than a new resolution.

The previous area formula gives

\begin{verbatim}
|W_h(d)|=3 binom(d,2)+(9d-3)h+9h^2.
\end{verbatim}

It is used only with its original parameters. Source attribution and the
full original P7 target are retained in §9.

\section{2. A positive nine-cell move, with both
directions}\label{edition-5-section-2}

For any integers \texttt{x,y}, define the following two lists of actual
tiles:

\begin{verbatim}
L(x,y) = {S(x,y), H(x+2,y), H(x+1,y+1)},
R(x,y) = {H(x,y), H(x,y+1), S(x+3,y)}.                 (T1)
\end{verbatim}

Each list partitions exactly the same nine cells:

\begin{verbatim}
{(x+v,y):0<=v<=4} union {(x+v,y+1):0<=v<=3}.
\end{verbatim}

Indeed L partitions the bottom row into lengths 2 and 3 and the top row
into lengths 1 and 3. R partitions the bottom row into lengths 3 and 2
and the top row into lengths 3 and 1. These descriptions prove
containment, disjointness and coverage, not just equality of areas.

On the free integer module on placed tiles, put

\begin{verbatim}
kappa = [H(x,y)]+[H(x,y+1)]+[S(x+3,y)]
        -[S(x,y)]-[H(x+2,y)]-[H(x+1,y+1)].           (T2)
\end{verbatim}

For the original incidence differential, \path|partial kappa=0|. On the
set of tilings containing L, adding kappa replaces that exact subpatch
by R. All removed coefficients were 1; the new patch is disjoint from
every tile outside the common nine-cell support. This gives a tiling
with coefficients 0 or 1. Subtraction of kappa is the inverse on the
tilings containing R. Thus these are inverse maps between the two
specified subsets of the original tiling fibre, not an assumed map
between all tilings.

Linear extension on the free modules on those subsets gives inverse
linear maps. They commute with the cell-incidence observation, since
both send a full tiling to its original region vector. Inclusion of the
nine-cell support into any larger finite support commutes with the
replacement. Translation has its literal inverse translation, so (T1)
also supplies all translated versions with the original coordinates
attached.

\section{3. Locate the move in the actual second
tiling}\label{edition-5-section-3}

Let \texttt{T2(d)} be the complete turn-1 second tiling, and translate
it by

\begin{verbatim}
t3=(1,-2).
\end{verbatim}

Its three difference-coordinate increments are \texttt{(-3,3,0)}. Its
region I is therefore contained in \texttt{W\_3(d)} and has exact bounds

\begin{verbatim}
D0 in [-d-8,2d+2],
D1 in [-d-2,2d+8],
D2 in [-d-5,2d+5].                                  (T3)
\end{verbatim}

The translated source tiling contains

\begin{verbatim}
S(d-2,0), H(d,0), H(d-1,1).                          (T4)
\end{verbatim}

Here is their provenance inside the original construction. The three
successive translations sum to zero: \path|(-2,1)+(1,1)+(1,-2)=(0,0)|.
The stone in (T4) is the original base stone indexed by \texttt{r=s=0};
it is distinct from the two previously consumed corner stones for
\texttt{d\textgreater{}=3}. The first-collar families \texttt{F\_(d-1)}
and \texttt{G\_(d-2)}, translated by \path|(1,1)+(1,-2)=(2,-1)|, are
respectively \texttt{H(d-1,1)} and \texttt{H(d,0)}. No second-collar
step removed a bone.

Apply (T1) at \path|(x,y)=(d-2,0)|. This changes exactly the two old
bones in (T4), and moves the source stone to

\begin{verbatim}
S*=S(d+1,0).
\end{verbatim}

Write T2* for the resulting tiling of I. In particular the old
frozen-source obstruction has been addressed by an actual positive move
of old tile generators.

\section{4. Complete third-collar table and
proof}\label{edition-5-section-4}

Use the following bone family C3. Every index bound is inclusive.

\begin{longtable}[]{@{}llll@{}}
\toprule\noalign{}
Label & Kind & Anchor & Range \\
\midrule\noalign{}
\endhead
\bottomrule\noalign{}
\endlastfoot
A\_r & D & \path|(-d-5+r,d+3-2r)| &
\texttt{0\textless{}=r\textless{}=d+2} \\
B & D & \texttt{(-d-4,d+3)} & one \\
C & D & \texttt{(-d-4,d+4)} & one \\
E & D & \texttt{(-d-3,d+4)} & one \\
F & D & \texttt{(-d-2,d+4)} & one \\
G\_r & D & \path|(-d+2r,d+3-r)| &
\texttt{0\textless{}=r\textless{}=d-1} \\
J\_r & H & \path|(-d-3+2r,d+5-r)| &
\texttt{0\textless{}=r\textless{}=d+1} \\
K & H & \texttt{(d,3)} & one \\
L & V & \texttt{(d+1,0)} & one \\
M & V & \texttt{(d+2,0)} & one \\
N & V & \texttt{(d+3,1)} & one \\
P & V & \texttt{(d+4,-1)} & one \\
Q & V & \texttt{(d+5,-3)} & one \\
\end{longtable}

\subsection{4.1 Original outer containment and exact inner
intersection}\label{edition-5-section-5}

At offset \texttt{u=0,1,2}, direct substitution gives:

\begin{longtable}[]{@{}llll@{}}
\toprule\noalign{}
Label & D0 & D1 & D2 \\
\midrule\noalign{}
\endhead
\bottomrule\noalign{}
\endlastfoot
A\_r & \texttt{2d+8-3r-2u} & \texttt{-d+3r+u} & \texttt{-d-8+u} \\
B & \texttt{2d+7-2u} & \texttt{-d-1+u} & \texttt{-d-6+u} \\
C & \texttt{2d+8-2u} & \texttt{-d-3+u} & \texttt{-d-5+u} \\
E & \texttt{2d+7-2u} & \texttt{-d-4+u} & \texttt{-d-3+u} \\
F & \texttt{2d+6-2u} & \texttt{-d-5+u} & \texttt{-d-1+u} \\
G\_r & \texttt{2d+3-3r-2u} & \texttt{-d-5+u} & \texttt{2-d+3r+u} \\
J\_r & \texttt{2d+8-3r-u} & \texttt{-d-6-u} & \texttt{-d-2+3r+2u} \\
K & \texttt{3-d-u} & \texttt{-d-5-u} & \texttt{2d+2+2u} \\
L & \texttt{-d-1+u} & \texttt{-d-2u} & \texttt{2d+1+u} \\
M & \texttt{-d-2+u} & \texttt{-d-1-2u} & \texttt{2d+3+u} \\
N & \texttt{-d-2+u} & \texttt{-d-4-2u} & \texttt{2d+6+u} \\
P & \texttt{-d-5+u} & \texttt{-d-1-2u} & \texttt{2d+6+u} \\
Q & \texttt{-d-8+u} & \texttt{2-d-2u} & \texttt{2d+6+u} \\
\end{longtable}

Every entry lies in \texttt{{[}-d-8,2d+8{]}}. Here is a full interval
check for the indexed rows, whose variables attain their extrema at the
displayed endpoints:

\begin{itemize}
\tightlist
\item
  A: D0 lies in \texttt{{[}-d-2,2d+8{]}}, D1 in \texttt{{[}-d,2d+8{]}},
  D2 in \texttt{{[}-d-8,-d-6{]}}.
\item
  G: D0 lies in \texttt{{[}2-d,2d+3{]}}, D1 in \texttt{{[}-d-5,-d-3{]}},
  D2 in \texttt{{[}2-d,2d+1{]}}.
\item
  J: D0 lies in \texttt{{[}3-d,2d+8{]}}, D1 in \texttt{{[}-d-8,-d-6{]}},
  D2 in \texttt{{[}-d-2,2d+5{]}}.
\end{itemize}

The ten fixed rows have only \texttt{u=0,1,2}; their extrema as printed
lie between the same bounds for \texttt{d\textgreater{}=3}. This proves
outer containment for every parameter in the theorem.

Compare the same table with (T3). A lies below I's D2 lower bound. B at
u=0 lies below that bound, and at u=1,2 lies above I's D0 upper bound.
C,E exceed I's D0 upper bound. F at u=0,1 exceeds that bound and at u=2
lies below I's D1 lower bound. G,J,K all lie below I's D1 lower bound. N
lies below that D1 bound; P,Q exceed I's D2 upper bound. L has only
u=0,1 inside I; M has only u=0 inside I. Those three cells are precisely
S*. Therefore

\begin{verbatim}
covered(C3) subset (W_3(d)\I) union S*.             (T5)
\end{verbatim}

\subsection{4.2 Disjointness of every orientation
pair}\label{edition-5-section-6}

For D-bones, the quantity \texttt{1-i-j} is constant. The A family
occupies distinct levels \texttt{3,...,d+5}; B,C,E,F occupy
\texttt{2,1,0,-1}; G occupies \texttt{-2,...,-d-1}. All D-bones are
disjoint. The J horizontal rows are distinct and range from 4 through
d+5; K occupies row 3. The five V-bones have distinct columns d+1
through d+5. Thus each orientation family is internally disjoint.

Every D-cell has \texttt{i\textless{}=d}: the A maximum is -1, the
fixed-row maximum is at most -d, and the G maximum is d.~Every V-cell
has \texttt{i\textgreater{}=d+1}, proving D/V disjointness.

For D/H, let u denote the D offset and v the H offset. A\_r meeting J\_s
would, from the second coordinate, require \texttt{s=2+2r+u}; the first
coordinate would then require \texttt{v=-6-3r-u\textless{}0}. For
B,C,E,F meeting J\_s, substituting the equal-row equation makes the J
first coordinate exceed the D first coordinate respectively by
\texttt{5+u+v}, \texttt{3+u+v}, \texttt{2+u+v}, \texttt{1+u+v}. For G\_r
meeting J\_s, equal rows require \texttt{s=2+r+u}, and the
first-coordinate excess is \texttt{1+u+v}. Each case is impossible for
nonnegative offsets.

For K in row 3, an A\_r cell in that row has \texttt{2r+u=d}, hence
first coordinate \texttt{-5-r\textless{}d}; a G\_r cell in row 3 has
\texttt{r+u=d}, hence first coordinate \texttt{r\textless{}=d-1}. The
fixed D rows have second coordinate at least d+1, which is at least 4.
None meets K's columns d,d+1,d+2. This exhausts D/H.

For H/V, J has row at least 4 whereas all five V-bones have row at most
3. K has row 3 and columns d,d+1,d+2. L and M stop at row 2; N,P,Q start
at column at least d+3. None meets K. All orientation pairs are now
checked.

\subsection{4.3 Coverage, full tiling, and exact boundary
equations}\label{edition-5-section-7}

The number of bones is

\begin{verbatim}
(d+3)+4+d+(d+2)+1+5=3d+15.
\end{verbatim}

Their disjoint cells number \texttt{9d+45}. The original area formula
gives

\begin{verbatim}
|W_3|-|W_2|+3=(9d+42)+3=9d+45.
\end{verbatim}

By (T5), equality of finite cardinalities proves

\begin{verbatim}
covered(C3)=(W_3\I) disjoint-union S*.              (T6)
\end{verbatim}

Remove S* from T2* and insert C3. This is the promised tiling T3 of W3.
It retains \path|binom(d,2)-3| stones. The move preserves the previous
\texttt{6d+12} bones; inserting C3 gives \texttt{9d+27} bones.

Let beta3 be the sum of the C3 bone generators and J3 the original
translation. Equation (T6) is the exact integral cochain identity

\begin{verbatim}
partial beta3 = 1_W3 - J3 1_W2 + partial[S*].        (T7)
\end{verbatim}

Together with (T2), this also gives the direct positive replacement:
remove the three tiles (T4), add the two moved horizontal bones, and add
C3. The net change of the bone chain and the consumed stone is recorded,
not hidden by a quotient.

The reflection \path|(i,j)->(i,1-i-j)| is its own inverse. It exchanges
the region parameters, carries a right-stone cell set to a right-stone
cell set, and permutes the three straight bone directions. Applying it
cellwise to T3 yields the reflected tiling with identical counts. This
completes the whole all-parameter construction.

\section{5. The full translation-invariant bone
cohomology}\label{edition-5-section-8}

The finite-jet observation from turn 1 admits an exact integral
calculation. Let

\begin{verbatim}
P=Z[X^±1,Y^±1],
H=1+X+X^2, V=1+Y+Y^2, D=X^2+XY+Y^2,
Q=P/(H,V,D).
\end{verbatim}

H and V are the cell polynomials of the bones at (0,0); D is the cell
polynomial of the D-bone at (0,2). Multiplication by Laurent monomials
is the original cell/tile translation. Thus the ideal is exactly the
image of the global, finite-support integer bone-incidence map, and Q is
its first cohomology.

Use a formal generator omega with \path|omega^2+omega+1=0|, and put

\begin{verbatim}
C=Z[omega],
E=C direct-sum F_3 sigma
\end{verbatim}

with multiplication

\begin{verbatim}
(a,c)(b,e)=(ab, abar*e+bbar*c),
\end{verbatim}

where \texttt{abar} is the ring map \texttt{C-\textgreater{}F\_3}
sending omega to 1. It is well-defined because \texttt{1+1+1=0} in F\_3.
This multiplication is associative and distributive by the ring and
module laws; its unit is (1,0). Equivalently sigma\^{}2=0,
\path|(omega-1)sigma=0|, and 3sigma=0. The additive decomposition is
literal, so sigma is nonzero of order three.

There is an algebra isomorphism

\begin{verbatim}
Q -> E,
X -> omega,
Y -> omega^2+sigma.                                (T8)
\end{verbatim}

Both images are units: their inverses are omega\^{}2 and omega-sigma. H
maps to zero. V maps to \path|(1+omega+omega^2)+3sigma=0|. D maps to the
same zero. Thus the map is defined on the displayed quotient.

Here is its inverse, proved within Q. Write \texttt{s=1+X+Y}. The
defining relations give

\begin{verbatim}
XY=X+Y+2,
(X-1)s=D-V=0,
(Y-1)s=D-H=0.
\end{verbatim}

Consequently \path|3s=(1+X+X^2)s=0| and \path|s^2=(1+X+Y)s=3s=0|. Also
\texttt{Y=X\^{}2+s}. Therefore

\begin{verbatim}
E -> Q, omega -> [X], sigma -> [1+X+Y]
\end{verbatim}

is well-defined. Both composites fix all the specified generators
X,Y,omega,sigma. They are inverse algebra homomorphisms. In particular
the underlying abelian group of the full quotient is
\path|Z^2 direct-sum Z/3|, with no torsion discarded.

\section{6. The original finite jet is exactly the reduction of this
quotient}\label{edition-5-section-9}

Let \path|A=F_3[epsilon,eta]/(epsilon^2,epsilon eta,eta^2)|. The
quotient map E-\textgreater A is

\begin{verbatim}
omega -> 1+epsilon,
sigma -> epsilon+eta.                             (T9)
\end{verbatim}

Under (T8), Y maps to \path|(1+epsilon)^2+epsilon+eta=1+eta|; hence (T9)
composes with the original cell map exactly as in turn 1. For an element
\path|a+b omega+c sigma|, its image is

\begin{verbatim}
(a+b) + (b+c)epsilon + c eta, reduced modulo 3.
\end{verbatim}

Its kernel is precisely the subgroup \texttt{3C} (with zero sigma
component), which is also 3E. Modulo 3 the inverse sends epsilon to
omega-1 and eta to omega\^{}2+sigma-1. Thus

\begin{verbatim}
Q/3Q -> A
\end{verbatim}

is an explicitly proved isomorphism, and (T8)-(T9) give the requested
relationship between the integral cohomology and the finite jet.

A translated cell has E-class

\begin{verbatim}
[X^i Y^j]=omega^(i+2j)+j sigma.                     (T10)
\end{verbatim}

This follows for all integers from the unit formulas and sigma\^{}2=0.
The torsion coefficient is taken modulo three. Summing (T10) over the
actual cyclically symmetric W\_h gives zero in C: the cyclic cell
permutation \path|(i,j)->(j,1-i-j)| changes \texttt{i+2j} by 2 modulo 3,
so the three residue counts agree. The coordinate sums also agree and
their sum is \textbar W\_h\textbar. Therefore

\begin{verbatim}
[1_W_h]=(binom(d,2)-h)sigma in Q.                   (T11)
\end{verbatim}

This uses the exact area formula modulo three. Multiplication by every
translation unit fixes sigma. The raw collar class is consequently
\texttt{-sigma} in the full integral quotient, and adding S* cancels it
there. Equation (T7) supplies the positive, support-contained lift,
which is stronger information than the quotient equality alone.

\section{7. Global, finite-support and positive lifts: actual comparison
maps}\label{edition-5-section-10}

For a finite cell set lambda, let C1(lambda) be its free integer cell
module and B(lambda) the image of the bone generators entirely contained
in lambda. Cell inclusion induces

\begin{verbatim}
C1(lambda)/B(lambda) -> Q.
\end{verbatim}

Its kernel is given by the explicit representative map

\begin{verbatim}
(C1(lambda) intersect im(partial_global))/B(lambda)
      -> ker(C1(lambda)/B(lambda)->Q), [z]->[z].     (T12)
\end{verbatim}

A class is killed precisely when its original representative is a global
bone boundary. The prequotient kernel is B(lambda); choosing that same
representative gives the inverse on every kernel class. This proves
(T12), including the retained finite support.

A concrete nonzero kernel class can be written without a search. Let
lambda be the union of the right stones at (0,0), (12,0), and (24,0),
and let z be its 0/1 cell vector. No original bone is contained in
lambda, so B(lambda)=0 and {[}z{]} is nonzero. Nevertheless z is a
global bone boundary. In the original Laurent polynomials,

\begin{verbatim}
3s=sH-(X+2)(D-V),
(X^m-1)s=(1+X+...+X^(m-1))(D-V), m=12,24.
\end{verbatim}

Adding these three identities produces the exact signed bone preimage of

\begin{verbatim}
(1+X^12+X^24)s=z.
\end{verbatim}

All terms are translates of the specified H,V,D bones. Their integer
coefficients and finite global support are supplied by the checker. This
realizes, rather than merely names, the comparison kernel. A second
example, the three cells (0,0),(3,0),(6,0), has zero finite jet but full
Q-class 3 in C, exhibiting the nonzero kernel 3C of (T9).

For the actual collar proof, no such out-of-region signed substitution
is used. The nine-cell move is the explicit kernel vector (T2), acting
with nonnegative resulting coefficients on its specified tiling
subfibres; (T7) is a positive preimage contained in the actual enlarged
region. These maps describe exactly how the obstruction calculation is
returned to the original tiling problem.

Finite supports ordered by inclusion and joined by union, with their
tile and cell modules, give the original coherent support diagram. The
pinned Split-Zero reconstruction sends \path|(lambda,[z])| to its
support-labelled class, supported zero to \texttt{(lambda,0)}, and
absence to the bottom zero. The upstream D2-D5 inverse maps, D6 quotient
and D7 kernel apply to the explicitly displayed diagrams. Translation is
a support-index morphism with its actual inverse. No theta estimate or
unproved positive lifting rule is imported.

\section{8. What the next calculation now starts
from}\label{edition-5-section-11}

The third-collar failure recorded in turn 1 is resolved by (T1)-(T7).
The construction gives h=3 for every d\textgreater=3, in addition to the
existing h=1,2. It does not supply a recurrence for all remaining h.

There are also five fully specified h=4 tiling certificates for
d=4,5,6,7,8 in \path|fourth-seeds.json|. Each translates T3 by (-2,1),
removes the original stone at (-1,4-d), releases eight specifically
recorded old bones and inserts the listed replacement bones. The checker
verifies the entire tiling and the positive residual equation using
integers. These five cases are bounded results. The recorded eight-bone
release is a valid construction, not a proved minimum or an asserted
uniform family.

The next research calculation is the support-contained transport of that
next source stone through the displayed eight-bone patch, followed by a
parameter-uniform table or a repeatable positive move sequence. The five
seeds retain every original generator, so the next session can work on
this specific map rather than restart a search. The full P7 interior
still includes h\textgreater=4 where h\textless binom(d,2), and the
known invariant-zero endpoints remain controls.

Two exploratory collar shapes were also tested: concentric bone-only
differences \path|W_(h+d)(d+1)\W_h(d)| for d=2..5,h=0..4. A
mixed-integer solver returned verified fillings for h=0,1 and infeasible
statuses for h=2,3,4. Only the returned positive tilings were checked
exactly; the solver's infeasibility statuses were not independently
certified and are not theorems. These tests do not authorize an
extension theorem. The successful argument of this note uses the
explicit tables above instead.

\section{9. Sources, tests and attribution}\label{edition-5-section-12}

\begin{itemize}
\tightlist
\item
  Original problem: Propp, *Trimer Covers in the Triangular Grid: Twenty
  Mostly Open Problems*, Problem 7,
  \href{https://www.samuelfhopkins.com/OPAC/files/proceedings/propp.pdf}{\path|https://www.samuelfhopkins.com/OPAC/files/proceedings/propp.pdf|}
  . The statement was already read in the pinned parent workbench.
\item
  Author's status page:
  \href{https://faculty.uml.edu/jpropp/benzels.html}{\path|https://faculty.uml.edu/jpropp/benzels.html|}
  . It retains a June-2025 status horizon. A bounded search on
  2026-09-15 found no established later full P7 resolution; it does not
  establish exhaustive novelty coverage.
\item
  Original area/invariant conventions: Defant--Li--Propp--Young,
  \href{https://arxiv.org/html/2209.05717v2}{\path|https://arxiv.org/html/2209.05717v2|}
  . The exact parameter/cell maps are in the parent proof.
\item
  Bone-only endpoint control: Kim--Propp, *A pentagonal number theorem
  for tribone tilings*,
  \href{https://arxiv.org/html/2206.04223v6}{\path|https://arxiv.org/html/2206.04223v6|}
  . No priority claim for (12,15) is made.
\item
  Split-Zero input:
  \path|KokunoYumeto/zeta-function-research-reader@1c8ec52c85c173adc9f8403a8a26914955e2f5a9|,
  \path|workbenches/splitzero-tandem/tex/support_diagrams.tex|, D1--D8.
\end{itemize}

The construction tables were found with bounded mixed-integer
exploration and then proved for all the stated parameters. The final
constructor and checker use only the Python standard library. They
verify original cells, shapes, multiplicities, the local kernel vector,
both cochain equations, counts, reflected tilings, the
\path|full-quotient/jet| square, and the five fourth-collar
certificates. No solver return is used as the proof of the all-parameter
result, and no external formal build or independent mathematical
acceptance is claimed. Third-party manuscripts and source corpora are
not redistributed.

\chapter{Turn 2: later third/fourth-collar edition}\label{edition-6}

Source: \path|sources/benzel-p7-turn2/benzel-p7/turn2/PROOF.md|.

SHA-256:
\path|d7337a59db163251955e867973215914617d7730b08afe186767e66bfabf1f35|.

\textbf{Propp 7 sprint, turn 2; 15 September 2026.} This continuation
retains the literal cells, placed tiles and source tilings of turn 1. It
constructs two further infinite families, not a full resolution of
Problem 7. All-parameter arguments and exact rational certificates are
supplied. Independent specialist acceptance and Lean verification have
not been obtained. No priority claim is made for the algebraic
computation or the special families.

The unchanged parent is \texttt{../PROOF.md}, pinned to workbench commit
\path|90f23b04fbde1e25d07342e2b8a8c80d919402c5|. \texttt{../check.py} is
its original search-free implementation. The original Propp statement is
at
\href{https://arxiv.org/html/2206.06472v4}{\path|https://arxiv.org/html/2206.06472v4|}
(Problem 7); its maintained author page, checked this turn, still labels
it open:
\href{https://faculty.uml.edu/jpropp/benzels.html}{\path|https://faculty.uml.edu/jpropp/benzels.html|}.
This is not an exhaustive status or novelty search.

\section{1. Literal domains, tiles, and
results}\label{edition-6-section-1}

A cell is an integer pair \texttt{(i,j)}, with inverse barycentric
coordinate map \path|(i,j) -> (i,j,1-i-j)|. Write

\begin{verbatim}
D0=j-i, D1=1-i-2j, D2=2i+j-1,
V(a,b)={(i,j): 1-a<=D0,D1,D2<=b-1},
W_h(d)=V(d+3h,2d+3h).
\end{verbatim}

The parameter map retains its inverse \texttt{d=b-a},
\texttt{h=(2a-b)/3}. A right stone is

\begin{verbatim}
S(i,j)={(i,j),(i+1,j),(i,j+1)}.
\end{verbatim}

The bones \path|D(i,j), H(i,j), V(i,j)| respectively contain

\begin{verbatim}
(i+u,j-u), (i+u,j), (i,j+u),  u=0,1,2.
\end{verbatim}

Every anchor below is an integer anchor of one of these original
prototiles. The source cell count and invariant, substituted through the
displayed parameter map, are

\begin{verbatim}
|W_h| = 3 binom(d,2) + (9d-3)h + 9h^2,
Delta(W_h) = binom(d,2)-h.
\end{verbatim}

The source of these formulas is Defant--Li--Propp--Young, *Tilings of
Benzels via the Abacus Bijection*,
\href{https://arxiv.org/html/2209.05717v2}{\path|https://arxiv.org/html/2209.05717v2|}.
The original invariant and its stone-count convention remain unchanged.

\textbf{Results proved here.} There is an explicitly constructed
type-103 tiling of \texttt{W\_3(d)} for every integer
\texttt{d\textgreater{}=3}, and of \texttt{W\_4(d)} for every integer
\texttt{d\textgreater{}=4}. Their tile counts are

\begin{longtable}[]{@{}lrr@{}}
\toprule\noalign{}
Region & Right stones & Bones \\
\midrule\noalign{}
\endhead
\bottomrule\noalign{}
\endlastfoot
\texttt{W\_3(d)} & \path|binom(d,2)-3| & \texttt{9d+27} \\
\texttt{W\_4(d)} & \path|binom(d,2)-4| & \texttt{12d+48} \\
\end{longtable}

In original coordinates these are \texttt{(n,2n-9)} for
\texttt{n\textgreater{}=12} and \texttt{(n,2n-12)} for
\texttt{n\textgreater{}=16}. The original coordinate involution
\path|(i,j)->(i,1-i-j)| gives the reflected families. It is its own
inverse, preserves right stones, and exchanges D/H bones while reversing
V-bone direction. Applying it to every cell and tile proves the
reflected assertions.

\section{2. Third collar: release two original
bones}\label{edition-6-section-2}

Let \texttt{T2(d)} be the exact second tiling in the parent proof.
Translate it by

\begin{verbatim}
t3=(1,-2),  I3=t3+W_2(d).
\end{verbatim}

The difference increments are \texttt{(-3,3,0)}. The inner bounds are

\begin{verbatim}
D0 in [-d-8,2d+2],
D1 in [-d-2,2d+8],
D2 in [-d-5,2d+5].                         (3.1)
\end{verbatim}

The outer bounds of \texttt{W\_3} are \texttt{{[}-d-8,2d+8{]}} for each
coordinate, proving containment of the actual translated source.

Remove the following original tiles:

\begin{verbatim}
S3=S(d-2,0),  B31=H(d-1,1),  B32=H(d,0).    (3.2)
\end{verbatim}

These are not hypothetical source choices. The three shifts from the
base through stage 3 sum to zero. Thus S3 is precisely the base stone
\texttt{q(0,0)} from the parent construction; it was neither of the two
earlier consumed corner stones for \texttt{d\textgreater{}=3}. B31 is
the parent's first-collar \texttt{F\_(d-1)} translated by
\texttt{(2,-1)}, and B32 is its \texttt{G\_(d-2)} translated by
\texttt{(2,-1)}. The second collar removes a stone, not these bones.
This identifies all three preimages in T2 and proves their presence and
distinctness.

Insert the following family C3:

\begin{longtable}[]{@{}llll@{}}
\toprule\noalign{}
Label & Bone & Anchor & Index range \\
\midrule\noalign{}
\endhead
\bottomrule\noalign{}
\endlastfoot
A\_r & D & \path|(-d-5+r,d+3-2r)| &
\texttt{0\textless{}=r\textless{}=d+2} \\
B & D & \texttt{(-d-4,d+3)} & one \\
C & D & \texttt{(-d-4,d+4)} & one \\
E & D & \texttt{(-d-3,d+4)} & one \\
F & D & \texttt{(-d-2,d+4)} & one \\
G\_r & D & \path|(-d+2r,d+3-r)| &
\texttt{0\textless{}=r\textless{}=d-1} \\
J\_r & H & \path|(-d-3+2r,d+5-r)| &
\texttt{0\textless{}=r\textless{}=d+1} \\
K & H & \texttt{(d,3)} & one \\
L & H & \texttt{(d-2,0)} & one \\
M & H & \texttt{(d-2,1)} & one \\
P & V & \texttt{(d+1,0)} & one \\
Q & V & \texttt{(d+2,0)} & one \\
N\_r & V & \path|(d+3+r,1-2r)| & \texttt{0\textless{}=r\textless{}=2} \\
\end{longtable}

\subsection{2.1 All original coordinate
bounds}\label{edition-6-section-3}

At cell offset u, the coordinate expressions are:

\begin{longtable}[]{@{}llll@{}}
\toprule\noalign{}
Label & D0 & D1 & D2 \\
\midrule\noalign{}
\endhead
\bottomrule\noalign{}
\endlastfoot
A\_r & \texttt{2d+8-3r-2u} & \texttt{-d+3r+u} & \texttt{-d-8+u} \\
B & \texttt{2d+7-2u} & \texttt{-d-1+u} & \texttt{-d-6+u} \\
C & \texttt{2d+8-2u} & \texttt{-d-3+u} & \texttt{-d-5+u} \\
E & \texttt{2d+7-2u} & \texttt{-d-4+u} & \texttt{-d-3+u} \\
F & \texttt{2d+6-2u} & \texttt{-d-5+u} & \texttt{-d-1+u} \\
G\_r & \texttt{2d+3-3r-2u} & \texttt{-d-5+u} & \texttt{-d+2+3r+u} \\
J\_r & \texttt{2d+8-3r-u} & \texttt{-d-6-u} & \texttt{-d-2+3r+2u} \\
K & \texttt{3-d-u} & \texttt{-d-5-u} & \texttt{2d+2+2u} \\
L & \texttt{2-d-u} & \texttt{3-d-u} & \texttt{2d-5+2u} \\
M & \texttt{3-d-u} & \texttt{1-d-u} & \texttt{2d-4+2u} \\
P & \texttt{-d-1+u} & \texttt{-d-2u} & \texttt{2d+1+u} \\
Q & \texttt{-d-2+u} & \texttt{-d-1-2u} & \texttt{2d+3+u} \\
N\_r & \texttt{-d-2-3r+u} & \texttt{-d-4+3r-2u} & \texttt{2d+6+u} \\
\end{longtable}

Each expression lies in \texttt{{[}-d-8,2d+8{]}} for the displayed r
ranges and \texttt{u=0,1,2}. This follows by evaluating its affine
minimum and maximum at the r endpoints and using
\texttt{d\textgreater{}=3}. The exact endpoint computations are also
replayed in \path|verify_certificates.py|, with no upper bound on d.

Against (3.1), A is below the inner D2 lower bound; B,C,E exceed its D0
upper bound; F,G,J,K are below its D1 lower bound; N exceeds its D2
upper bound. L and M lie wholly in I3. P has exactly offsets 0,1 in I3,
and Q has exactly offset 0 there. These nine inner cells are precisely

\begin{verbatim}
S3 disjoint-union B31 disjoint-union B32.    (3.3)
\end{verbatim}

Their rows are \path|j=0, i=d-2,...,d+2| and \path|j=1, i=d-2,...,d+1|.
The table and (3.2) give these same nine original coordinates.

\subsection{2.2 Disjointness without a finite-d
extrapolation}\label{edition-6-section-4}

The D-bone line sums i+j are respectively \texttt{-2-r}, \texttt{-1},
\texttt{0}, \texttt{1}, \texttt{2}, and \texttt{3+r}. Their integer
ranges are disjoint, and the varying lines in each family are distinct.
The H rows are \texttt{d+5-r} (at least 4), 3, 0 and 1. The V columns
are \texttt{d+1}, \texttt{d+2}, \texttt{d+3+r}. This proves disjointness
among bones of a common orientation.

All D cells have i at most d.~All V cells have i at least d+1, so D and
V cannot meet. For H versus V, J has j at least 4, whereas every V cell
has j at most 3. L and M have i at most d.~K has columns at most d+2 and
row 3: P and Q stop at row 2, while N has columns at least d+3. These
prove every H/V separation by the original coordinates.

It remains to compare D and H. A, B,C,E,F have columns at most -1 and
cannot meet K,L,M, whose columns are positive. G has row at least 2,
excluding L,M. Equality of a G cell with row 3 forces \texttt{r+u=d};
its column is then r, below K's smallest column d.

For A\_r and J\_s, equality of rows gives \texttt{s=2+2r+u}. Equality of
columns would then give \texttt{6+3r+u+v=0}, impossible for nonnegative
r,u,v. For any of B,C,E,F, write its anchor \texttt{(-d+b,d+c)}.
Equality with J\_s would give \texttt{u+v=b+2c-7}; these right sides are
respectively -5,-3,-2,-1. For G\_r versus J\_s, the two equalities give
\texttt{u+v=-1}. These contradictions exhaust the D/H pairs.

\subsection{2.3 The positive incidence
identity}\label{edition-6-section-5}

C3 contains \texttt{3d+17} bones, hence \texttt{9d+51} pairwise distinct
cells. The original area formula gives

\begin{verbatim}
|W_3|-|W_2|+9 = (9d+42)+9 = 9d+51.
\end{verbatim}

Containment, (3.3), disjointness and cardinality therefore prove the
exact partition

\begin{verbatim}
covered(C3) = (W_3\I3) disjoint-union S3 disjoint-union B31 disjoint-union B32.
\end{verbatim}

Consequently

\begin{verbatim}
T3 = (t3+T2)\{S3,B31,B32} union C3              (3.4)
\end{verbatim}

is a tiling for every integer d\textgreater=3. As an equality in the
free integer cell module,

\begin{verbatim}
partial beta3 = 1_W3 - J3 1_W2 + partial(S3+B31+B32),   (3.5)
\end{verbatim}

where beta3 is the sum of the listed bone generators, all with
coefficient +1. The number of bones changes from \texttt{6d+12} to
\path|6d+12-2+(3d+17)=9d+27|; exactly one right stone is consumed.

\section{3. Fourth collar: release eight original
bones}\label{edition-6-section-6}

Translate T3 by \texttt{t4=(-2,1)}. The original translated source I4
lies in W4 and has bounds

\begin{verbatim}
D0 in [-d-5,2d+11],
D1 in [-d-8,2d+8],
D2 in [-d-11,2d+5].                          (4.1)
\end{verbatim}

Remove the right stone and eight bones

\begin{verbatim}
S4=S(-1,4-d),
D(-1,1-d), H(0,3-d),
V(-1,-d-2), V(0,-d-3), V(1,-d-4),
V(1,-d), V(2,-d-3), V(2,-d).                 (4.2)
\end{verbatim}

Every listed source tile is present for every integer d\textgreater=4.
S4 is the base tile \texttt{q(0,d-3)} translated by the cumulative shift
\texttt{(-2,1)}. It is distinct from all three consumed corners. The D
bone is the parent's first-collar C, and the H bone is its G\_0; both
have total later shift zero at stage 4. G\_0 differs from the G\_(d-2)
removed at stage 3. The six vertical bones, in the order printed, are
the parent's second-collar K\_(d-1), L, M (each shifted by
\texttt{(-1,-1)}), first-collar E (shift zero), second-collar N\_0
(shift \texttt{(-1,-1)}), and first-collar J\_0 (shift zero). None was
removed at stage 3. This supplies their explicit preimages rather than
assuming a source configuration.

Insert C4 given by:

\begin{longtable}[]{@{}llll@{}}
\toprule\noalign{}
Label & Bone & Anchor & Index range \\
\midrule\noalign{}
\endhead
\bottomrule\noalign{}
\endlastfoot
A & D & \texttt{(-4,-d-2)} & one \\
B & D & \texttt{(-2,-d-3)} & one \\
C & D & \texttt{(-1,-d-2)} & one \\
D & D & \texttt{(-1,4-d)} & one \\
E & D & \texttt{(-1,5-d)} & one \\
F & D & \texttt{(0,-d-4)} & one \\
G & D & \texttt{(1,-d-3)} & one \\
H & D & \texttt{(1,-d-2)} & one \\
I & D & \texttt{(2,-d-5)} & one \\
J\_r & H & \path|(-d-4+2r,d+7-r)| &
\texttt{0\textless{}=r\textless{}=d+2} \\
K\_r & H & \path|(3+r,-d-3+r)| &
\texttt{0\textless{}=r\textless{}=d+1} \\
L & H & \texttt{(d+2,3)} & one \\
M & H & \texttt{(d+2,4)} & one \\
N & H & \texttt{(d+3,1)} & one \\
O & H & \texttt{(d+3,2)} & one \\
P & H & \texttt{(d+4,-1)} & one \\
Q & H & \texttt{(d+4,0)} & one \\
R & V & \texttt{(-1,-d-1)} & one \\
S & V & \texttt{(0,-d-2)} & one \\
T & V & \texttt{(1,-d-1)} & one \\
U & V & \texttt{(2,-d-2)} & one \\
V & V & \texttt{(2,1-d)} & one \\
W & V & \texttt{(4,-d-6)} & one \\
Z\_r & V & \path|(5+r,-d-6+r)| &
\texttt{0\textless{}=r\textless{}=d+2} \\
\end{longtable}

The letter in the first column is a family label, while the second
column determines the bone orientation.

\subsection{3.1 Containment and exact inner
cells}\label{edition-6-section-7}

All coordinates in this table are retained in \texttt{tables.json}.
Substitution in the three original D functions, followed by the exact
affine endpoint test detailed in section 4, proves all outer
inequalities \path|-d-11<=Di<=2d+11| for d\textgreater=4. This includes
all 432 one-sided bounds; there are no untested residue classes or upper
bounds on d.

For inner membership, J has D1=\texttt{-d-9-u}, below the lower bound in
(4.1). K has D0=\texttt{-d-6-u}, and Z has D0=\texttt{-d-11+u}, both
below the inner lower bound. A,B,F have D1=\texttt{2d+9+u}, above the
inner upper bound. I and W are below the inner D0 lower bound.
L,M,N,O,P,Q have D2 either \texttt{2d+6+2u} or \texttt{2d+7+2u}, above
the inner upper bound.

The complete remaining inner-cell table is

\begin{longtable}[]{@{}ll@{}}
\toprule\noalign{}
Replacement family & Offsets inside I4 \\
\midrule\noalign{}
\endhead
\bottomrule\noalign{}
\endlastfoot
C, D, E & 0,1,2 \\
G & 0 \\
H & 0,1 \\
R, S, T, U, V & 0,1,2 \\
\end{longtable}

Substitution proves that all these cells lie in I4 for every
d\textgreater=4. For G the D0 expression is \texttt{-d-4-2u}, and for H
it is \texttt{-d-3-2u}, giving exactly the stated cutoffs. The 27 affine
coordinate expressions in this table equal, as a multiset of affine
functions of d, the 27 cells of (4.2). \path|verify_certificates.py|
compares these expressions symbolically, not by sampling d.

\subsection{3.2 Exact all-parameter disjointness
certificate}\label{edition-6-section-8}

The certificate \path|evidence/affine-certificates.json| contains a
rational contradiction for each of the 2,484 cross-family
cell-intersection systems. The systems quantify over the unbounded
variables d,r,s, with d\textgreater=4 and the exact index ranges in the
table. Section 4 proves how every certificate establishes the required
disjointness, gives the constraint construction, and explains the
independent arithmetic replay. For the three indexed families, their
anchor step and bone direction have nonzero determinant. Thus cells
within a family have unique (r,u), completing within-family
disjointness. No optimization solver or finite parameter test is used to
accept these statements.

\subsection{3.3 Partition and tiling}\label{edition-6-section-9}

The number of inserted bones is

\begin{verbatim}
9 + (d+3)+(d+2)+6+6+(d+3) = 3d+29.
\end{verbatim}

The original region area difference is \path!|W_4|-|W_3|=9d+60!. Adding
the 27 released source cells gives \path|9d+87=3(3d+29)|. Containment,
the exact intersection and disjointness therefore prove

\begin{verbatim}
covered(C4) = (W_4\I4) disjoint-union covered(the nine tiles in (4.2)).
\end{verbatim}

Removing those nine original tiles and inserting C4 constructs T4. Its
number of bones is \path|9d+27-8+(3d+29)=12d+48|; its right-stone count
is \path|binom(d,2)-4|. With eta4 denoting the sum of the eight released
bone generators, the cochain identity is

\begin{verbatim}
partial beta4 = 1_W4 - J4 1_W3 + partial S4 + partial eta4.   (4.3)
\end{verbatim}

All beta4 coefficients are +1. This proves the second announced family
for every d\textgreater=4.

\section{4. What the affine certificate proves, and how it is
checked}\label{edition-6-section-10}

The certificate handles exact rational inequalities, not approximate LP
output. A family is stored as

\begin{verbatim}
label, bone direction v, x=(xd,xr,xc), y=(yd,yr,yc), upper=(ud,uc).
\end{verbatim}

Its cell at offset u has coordinates

\begin{verbatim}
(xd*d+xr*r+xc+u*v_x, yd*d+yr*r+yc+u*v_y),
d>=dmin, 0<=r<=ud*d+uc, u in {0,1,2}.
\end{verbatim}

These are the original table coordinates. No change of variables
replaces them.

For two different families and fixed u,v in \{0,1,2\}, equality of their
two coordinates gives two affine equations in (d,r,s). The nine
constraints, in their exact certificate order, are

\begin{enumerate}
\def\labelenumi{\arabic{enumi}.}
\tightlist
\item
  \texttt{-d\textless{}=-dmin};
\item
  \texttt{-r\textless{}=0};
\item
  \texttt{r-ud*d\textless{}=uc};
\item
  \texttt{-s\textless{}=0};
\item
  \path|s-ud'*d<=uc'|; 6--7. equality of x coordinates, written as its
  two opposite inequalities; 8--9. equality of y coordinates, written as
  its two opposite inequalities.
\end{enumerate}

Each certificate supplies nonnegative rational multipliers for these
nine inequalities. Their weighted left-hand sides sum to zero and their
weighted right-hand sides sum to a strictly negative rational number. A
collision would therefore give \texttt{0\textless{}0}. This proves
infeasibility over the reals, hence also over the original integer
parameters. Every pair and offset is enumerated and verified exactly;
missing, duplicate, extra or malformed cases fail. There are 702 cases
for the 13 third-collar families and 2,484 for the 24 fourth-collar
families.

The construction script for these witnesses uses Fourier--Motzkin
elimination with rational arithmetic, retaining the nonnegative
multiplier of every original inequality. Its search algorithm is not
trusted by the verifier: \path|verify_certificates.py| reconstructs all
nine inequalities from \texttt{tables.json}, multiplies the supplied
rationals, checks the zero left side and strict negative right side, and
verifies the complete case denominator.

For outer containment and claimed inner containment, a required
difference is \texttt{a*d+b*r+c}. Its minimum on the exact r interval
occurs at r=0 when b\textgreater=0 and at r=ud*d+uc when b\textless0.
Substitution leaves \texttt{A*d+C}. The code verifies
\texttt{A\textgreater{}=0} and \texttt{A*dmin+C\textgreater{}=0}. For an
outside cell it verifies one violated inner bound with integer gap at
least one by the same calculation. This proves each assertion on the
unbounded domain. It also verifies every range is nonempty and every
indexed family's cell map is injective by the two-direction determinant.

Finally, the checker compares the exact affine coordinate multisets for
the inner intersection and removed source cells, and compares the
complete linear cardinality polynomials. These, together with the
explicit original-source membership in sections 2 and 3, establish the
two positive partitions. The numeric replay is additional implementation
evidence, not the all-parameter argument.

\section{5. The complete unrestricted lattice bone
cokernel}\label{edition-6-section-11}

Here `unrestricted' means that placed bones at every integer translation
are included. The precise finite-support maps are retained in section 6.

Let

\begin{verbatim}
Lambda=Z[X^+-1,Y^+-1],
fX=1+X+X^2, fY=1+Y+Y^2, g=X^2+XY+Y^2,
J=(fX,fY,g), Q=Lambda/J, S=1+X+Y.
\end{verbatim}

Identify the original cell (i,j) with \texttt{X\^{}i\ Y\^{}j}.
Horizontal and vertical bone boundaries generate fX and fY with all
translations. A D-bone has polynomial \path|1+XY^-1+X^2Y^-2|, whose
product with the unit Y\^{}2 is g. Thus J is exactly the image of the
global original bone boundary map. This identifies Q with the first
cohomology of that particular two-term lattice complex.

\subsection{5.1 An explicit ring isomorphism, including its
inverse}\label{edition-6-section-12}

Define

\begin{verbatim}
E=Z[t]/(t^2+t+1),
bar:E->F_3,  bar(t)=1.
\end{verbatim}

The evaluation is well-defined because \texttt{1+1+1=0} in F\_3. On the
additive group \path|E direct-sum F_3*z|, define multiplication by

\begin{verbatim}
(a,c)*(b,e)=(ab, bar(a)e+bar(b)c).             (5.1)
\end{verbatim}

The unit is (1,0). Direct expansion using multiplicativity of bar proves
associativity and distributivity. In this ring z\^{}2=0 and t*z=z.

The original-coordinate assignment is

\begin{verbatim}
X -> (t,0),
Y -> (-1-t,1),
S -> (0,1)=z.                               (5.2)
\end{verbatim}

Both X and Y images have cube one: the E components do, and for Y the z
coefficient in its cube is 3=0. Thus (5.2) extends to the original
Laurent ring. Substitution gives zero on fX, fY and g, so it defines a
ring homomorphism from Q to (5.1).

The inverse sends

\begin{verbatim}
(a(t),c) -> a(X)+c*S in Q.                   (5.3)
\end{verbatim}

Here are the exact original polynomial identities that prove
well-definedness and multiplicativity:

\begin{verbatim}
(X-1)S = g-fY,
(Y-1)S = g-fX,
3S = S*fX - (X+2)(g-fY),
S^2 = 3S+(X-1)S+(Y-1)S.                    (5.4)
\end{verbatim}

All right sides vanish in Q. Therefore 3S=0, S\^{}2=0, and
multiplication of S by any a(X) acts through a(1) modulo 3, exactly as
in (5.1). The relation fX=0 makes a(t)'s representative immaterial.
These verify (5.3).

The composites fix X and Y because \texttt{-1-X+S=Y}, and fix t and z by
(5.2). They therefore fix every Laurent polynomial class and every
element of (5.1). This proves the ring isomorphism with both inverse
maps:

\begin{verbatim}
Q ~= E semidirect F_3*z,   z^2=0, t*z=z.     (5.5)
\end{verbatim}

In particular, the additive group is free of rank two plus one summand
Z/3, and S has exact order three: its image is the nonzero z. This is a
complete computation for the displayed unrestricted boundary map, not
merely a nonzero finite observation.

\subsection{5.2 The exact morphism to the first-order
jet}\label{edition-6-section-13}

Let

\begin{verbatim}
A=F_3[epsilon,eta]/(epsilon^2,epsilon*eta,eta^2).
\end{verbatim}

The original observation \path|X->1+epsilon, Y->1+eta| kills fX,fY,g. It
descends through Q. Write a class in (5.5) as \texttt{a+b*t+c*z}, with
a,b integers and c in F\_3. Its jet image is

\begin{verbatim}
(a+b) + (b+c)*epsilon + c*eta.              (5.6)
\end{verbatim}

Modulo three, the inverse coordinate map is

\begin{verbatim}
c = eta coefficient,
b = epsilon coefficient - eta coefficient,
a = constant coefficient - b.
\end{verbatim}

Both composites are identities. Hence the first-order jet is exactly
Q/3Q, and its kernel in Q is precisely 3Q. Equivalently,
\path|(t-1)->epsilon|, \path|z->epsilon+eta|, with inverse
\path|eta->z-(t-1)|. This identifies the whole finite jet and the
information it forgets.

\subsection{5.3 Exact benzel and collar classes in
Q}\label{edition-6-section-14}

Projection \texttt{(a,c)-\textgreater{}a} in (5.1) is the original
evaluation \texttt{X-\textgreater{}t,Y-\textgreater{}t\^{}2}. A cell
maps to \texttt{t\^{}(i+2j)}. Cyclic permutation of barycentric
coordinates induces \path|(i,j)->(j,1-i-j)| and increases i+2j by 2
modulo three. This permutation preserves W\_h, and its cell orbits have
size three: a fixed cell would have i=j=k=1/3, not integers. The
projected cell sum on each orbit is a monomial times
\texttt{1+t+t\^{}2=0}. Thus the region class lies in the exact kernel
F\_3*S of that projection.

Cyclic permutation also makes the three integer coordinate sums equal to
\texttt{\textbar{}W\_h\textbar{}/3}. The jet of the region is therefore
\path!(|W_h|/3)*(epsilon+eta)!. Substitution of the source area formula
gives \path!|W_h|/3 = binom(d,2)+(3d-1)h+3h^2!, congruent to
\path|binom(d,2)-h| modulo three. Since the jet is injective on F\_3*S
by (5.6), the equality already holds in the full Q:

\begin{verbatim}
[1_W_h] = (binom(d,2)-h)*S.                (5.7)
\end{verbatim}

The identities \path|(X-1)S=(Y-1)S=0| show that every translation acts
identically on S. Thus the actual residual of every containing collar
has the exact class

\begin{verbatim}
[1_W_(h+1)-J_h 1_W_h] = -S in Q.         (5.8)
\end{verbatim}

This class has order exactly three. A right-stone boundary represents S;
adding one cancels (5.8). Equations (3.5) and (4.3) exhibit positive,
original-support preimages for that cancellation after the explicitly
listed old bones are released.

\section{6. Finite support, the global quotient, and positive tilings:
actual maps}\label{edition-6-section-15}

For each finite cell set lambda, take the original free cell module
C\^{}1(lambda), the free module on placed bones contained in lambda, and
its cell-incidence differential. Inclusion of finite supports gives
coherent cochain maps, indexed by the join-semilattice of finite subsets
with union and empty bottom.

The cell-polynomial map embeds C\^{}1(lambda) into Lambda. It induces

\begin{verbatim}
H_B^1(lambda) -> Q,
[sum n_(i,j)[(i,j)]] -> [sum n_(i,j)X^i Y^j].     (6.1)
\end{verbatim}

Its kernel is explicitly

\begin{verbatim}
(C^1(lambda) intersect J) / im partial_B(lambda),  (6.2)
\end{verbatim}

through the same representative map. The prequotient kernel is exactly
the denominator and each killed class has a representative in the
numerator, giving both directions. The intersection uses the specified
embedding into Lambda. It retains relations whose unrestricted boundary
witnesses leave lambda rather than assuming those witnesses stay in the
original support.

The directed union of these complexes has cell module Lambda and
boundary image J: every Laurent polynomial combination of fX,fY,g uses
finitely many original placed bones and hence is contained in some
finite support. Thus the colimit cohomology maps to Q with inverse
obtained from any finite representative; enlarging the support
containing all finitely many relation tiles proves independence. This
gives the precise finite/global comparison.

Apply the upstream Split-Zero reconstruction to the full support
diagram. The maps (6.1) and (5.6) induce the labelled maps

\begin{verbatim}
(lambda,[p]) -> (lambda,[p] in Q) -> (lambda,phi(p) in A).
\end{verbatim}

They commute with the actual support inclusions and preserve each
supported zero and the global bottom separately. The actual comparison
from bone generators to bone-plus-right-stone generators has kernel
\path|(im partial_103)/(im partial_B)| via
\texttt{{[}p{]}-\textgreater{}{[}p{]}}, as in the parent's displayed
specialization of upstream D7. These are the maps used by the collar
calculation.

Positive tilings retain a further explicit source map. The inclusion of
the nonnegative integer tile monoid into its free integer module gives
the commuting square

\begin{verbatim}
N^(placed permitted tiles)  ----inclusion----> Z^(placed permitted tiles)
            | partial+                               | partial
            v                                        v
N^(original cells)          ----inclusion----> Z^(original cells).
\end{verbatim}

The positive tiling set is the fibre of partial+ over the literal
all-ones cell vector. Every coefficient in that fibre is 0 or 1, since a
tile coefficient at least 2 would make each of its cell coefficients at
least 2. Coverage and disjointness follow from the coefficient-one
target. Inclusion induces the exact map from this fibre into the signed
integer solution fibre.

The cokernel (5.5) concerns the displayed group completion. The positive
fibres for sections 2 and 3 are supplied by the explicit beta3 and
beta4, not inferred merely from their image in Q. Release of the two or
eight source bones is an actual enlargement of the movable generator
set; the positive source configurations before and after are both
retained in (3.5) and (4.3).

\section{7. Computation, unsuccessful choices, and the next exact
target}\label{edition-6-section-16}

Exploratory integer programming found the tile arrangements. The
finalized constructors use only the printed formulas. The all-parameter
verifier uses exact rational arithmetic and the complete finite list of
certificates, independently of the search algorithm. The finite replay
checks cells, original tile shapes, multiplicity-one coverage, source
membership, source/target cochain identities, reflection, and the exact
cokernel and jet observations.

Five additional exact tiling witnesses for h=5, d=4,\ldots,8 are saved
in \path|evidence/fifth-witnesses.json|. They start from the explicit
T4, translate by (1,1), and consume \texttt{S(3-d,d-1)}. The saved
witness at d=4 moves 22 old bones and the saved witnesses at
d=5,\ldots,8 each move 16. These are counts of those witnesses, not
mathematical lower bounds. The solver reached its exploration cutoff on
the d=4 optimization, but its saved feasible tiling passes the integer
verifier; no optimality conclusion is used. The h=5 data is not an
all-parameter construction.

The immediate next calculation is the fifth-collar positive repair on
the same source, with the explicit moved-bone set included, followed by
a uniform generator description through varying h. The earlier interior
range was \path|d>=3, 1<=h<=binom(d,2)-1|. The constructions now cover
h=1,2,3,4 whenever they lie in the nonnegative-invariant domain, plus
the five recorded h=5 examples. The full P7 statement has not been
proved by this continuation.

\section{Sources and status}\label{edition-6-section-17}

\begin{itemize}
\tightlist
\item
  Parent proof and source tilings: workbench commit
  \path|90f23b04fbde1e25d07342e2b8a8c80d919402c5|,
  \path|workbenches/splitzero-nonzeta/sprints/benzel-p7/PROOF.md| and
  \texttt{check.py}.
\item
  Original Propp list:
  \href{https://arxiv.org/html/2206.06472v4}{\path|https://arxiv.org/html/2206.06472v4|},
  Problem 7 and original \path|coordinates/prototiles.|
\item
  Original benzel area and invariant:
  \href{https://arxiv.org/html/2209.05717v2}{\path|https://arxiv.org/html/2209.05717v2|}.
\item
  Maintained author status:
  \href{https://faculty.uml.edu/jpropp/benzels.html}{\path|https://faculty.uml.edu/jpropp/benzels.html|};
  still labels P7 open in the retrieved page. This does not establish an
  exhaustive current literature audit.
\item
  Split-Zero source:
  \path|KokunoYumeto/zeta-function-research-reader@1c8ec52c85c173adc9f8403a8a26914955e2f5a9|,
  \path|workbenches/splitzero-tandem/tex/support_diagrams.tex|, D1--D8.
  The current proof specifies all specialized maps rather than importing
  an analytic zeta estimate.
\end{itemize}

The complete mathematics above is written work with replayable exact
certificates. It is not an independently accepted solution of a numbered
open problem. Third-party books, papers and formal-source corpora are
not redistributed in this package.

\chapter{Turn 3: fifth collar and polygon maps}\label{edition-7}

Source: \path|sources/benzel-p7-turn3/benzel-p7-turn3/PROOF.md|.

SHA-256:
\path|1145b57c01c6da29ecb3154ac831896ded42db4116d0245f061cb0e7e50f6924|.

Propp 7 focused attempt, turn 3. Recorded 15 September 2026.

\section{0. Exact scope and input lineage}\label{edition-7-section-1}

The new all-parameter result is a type-103 tiling of

\begin{verbatim}
W_5(d)=V(d+15,2d+15), for every integer d>=4,
\end{verbatim}

and its reflected region. There are binom(d,2)-5 right stones and 15d+75
bones. The proof uses 16 released source bones for all d\textgreater=4,
including d=4. The earlier d=4 witness released 22; it is preserved, not
called minimal or overwritten. In original parameters the new family is
V(n,2n-15), n\textgreater=19, through d=n-15 and its inverse n=d+15.

The constructions of collars 1--4 in tables.json come from the earlier
turn-1 and uploaded turn-2 sources. The previously supplied
fourth-collar result is included with its complete indexed table and a
fresh exact certificate replay. It is not counted as a new result of
turn 3. The current GitHub head inspected was
038ed02bb524e60e9c2bd20c3a0aa02365f4001d. Its historical turn2 and
TURN2.md are preserved. This continuation is additive; it does not
silently replace either of those source generations.

The original P7 statement is Propp,
\href{https://arxiv.org/abs/2206.06472}{\path|https://arxiv.org/abs/2206.06472|}.
The author page
\href{https://faculty.uml.edu/jpropp/benzels.html}{\path|https://faculty.uml.edu/jpropp/benzels.html|}
still labels P7 open, with an explicit June-2025 status horizon. The
bounded search on this turn did not establish a later resolution. This
is not an exhaustive novelty search. The original benzel coordinates and
invariant convention are those used in
\href{https://arxiv.org/html/2209.05717v2}{\path|https://arxiv.org/html/2209.05717v2|}.
The support reconstruction input remains pinned to
\href{mailto:KokunoYumeto/zeta-function-research-reader@1c8ec52c85c173adc9f8403a8a26914955e2f5a9}{\path|KokunoYumeto/zeta-function-research-reader@1c8ec52c85c173adc9f8403a8a26914955e2f5a9|},
\path|workbenches/splitzero-tandem/tex/support_diagrams.tex,| D1--D8.

The full P7 target remains unresolved by this work. No independent
specialist review, Lean build, priority certification, or public
resolution claim is asserted.

\section{1. Original carriers and an explicit
bijection}\label{edition-7-section-2}

A cell is (x,y) in Z\^{}2, with inverse barycentric map
(x,y)-\textgreater(x,y,1-x-y). Its three original differences are

\begin{verbatim}
D0=y-x, D1=1-x-2y, D2=2x+y-1.
\end{verbatim}

The original region V(a,b) is defined by \path|1-a<=D0,D1,D2<=b-1.|
Throughout,

\begin{verbatim}
W_h(d)=V(d+3h,2d+3h), d>=3, h>=0.
\end{verbatim}

The original tile types have offsets

\begin{verbatim}
R: (0,0),(1,0),(0,1);
D: (0,0),(1,-1),(2,-2);
H: (0,0),(1,0),(2,0);
V: (0,0),(0,1),(0,2).
\end{verbatim}

No other trihex is included. Every anchor and tile phase is retained
below.

Put e\_0=(0,0), e\_1=(1,0), e\_2=(0,1). Define

\begin{verbatim}
Phi_d(r,s,p)=(d-2-2r-s,r-s)+e_p,   p in {0,1,2}.
\end{verbatim}

The inverse on an original cell is

\begin{verbatim}
p = x+2y-(d-2) mod 3, chosen in {0,1,2};
(alpha,beta)=e_p;
r=(y-x+d-2-beta+alpha)/3;
s=(1-x-2y+d-3+alpha+2beta)/3.
\end{verbatim}

The choice of p makes both numerators divisible by 3: y-x is congruent
to -(x+2y), and 1-x-2y is congruent to 1-(x+2y). Substitution of the
inverse in Phi\_d gives (x,y); substitution of Phi\_d in the inverse
gives r,s,p, since alpha+2beta=p.~Thus these are inverse maps on the
entire original cell lattice, not just on the chosen tilings.

Write R=r+h and S=s+h. Substitution in the six benzel inequalities gives
the following exact three polygon domains:

\begin{longtable}[]{@{}lrrr@{}}
\toprule\noalign{}
phase p & upper bound on R & upper bound on S & lower bound on R+S \\
\midrule\noalign{}
\endhead
\bottomrule\noalign{}
\endlastfoot
0 & d+2h-1 & d+2h-2 & h-1 \\
1 & d+2h-1 & d+2h-1 & h \\
2 & d+2h-2 & d+2h-1 & h-1 \\
\end{longtable}

All three also have R\textgreater=0, S\textgreater=0, and
R+S\textless=d+3h-2. These are equalities of the transported integer
domains. For example, \path|D0=3r-d+2+beta-alpha| and
\path|D1=3s-d+3-alpha-2beta| give the R,S bounds with their displayed
integer rounding; \path|D2=2d-5-3(r+s)+2alpha+beta| gives the sum
bounds. Evaluating the three possible offsets gives exactly the table,
so the inverse maps restrict to the region and these polygons in both
directions.

\subsection{1.1 Area and base tiling from the same
bijection}\label{edition-7-section-3}

Each polygon is the triangle R,S\textgreater=0, R+S\textless=d+3h-2 with
three corner triangles removed. The removed sizes are two copies of
binom(h,2) and one copy of binom(h+1,2), in the order specified by the
table. The upper-R and upper-S cuts do not intersect: their least
combined sum exceeds d+3h-2 for d\textgreater=3,h\textgreater=0. Neither
can intersect the removed lower-sum triangle. For h=0 the three removed
counts are zero and all cuts are empty. Therefore each phase has exactly

\begin{verbatim}
binom(d+3h,2)-2 binom(h,2)-binom(h+1,2)
  = binom(d,2)+(3d-1)h+3h^2
\end{verbatim}

cells. This proves directly, using the original cell bijection,

\begin{verbatim}
|W_h(d)|=3 binom(d,2)+(9d-3)h+9h^2.       (1)
\end{verbatim}

At h=0 all three polygons are exactly r,s\textgreater=0,
r+s\textless=d-2. Grouping the three cells with the same (r,s) gives the
original right stone R(d-2-2r-s,r-s). These stones form T0. Their
disjointness and exhaustiveness follow from both inverse maps, proving
the base tiling without a finite enumeration.

\subsection{1.2 All original tile phases}\label{edition-7-section-4}

For a tile of kind K anchored at Phi\_d(r,s,p), transport every original
cell by Phi\_d\^{}\{-1\}. Each entry (q,u,v) below denotes the original
cell \path|Phi_d(r+u,s+v,q).|

\begin{longtable}[]{@{}
  >{\raggedright\arraybackslash}p{(\columnwidth - 4\tabcolsep) * \real{0.3000}}
  >{\raggedleft\arraybackslash}p{(\columnwidth - 4\tabcolsep) * \real{0.4000}}
  >{\raggedright\arraybackslash}p{(\columnwidth - 4\tabcolsep) * \real{0.3000}}@{}}
\toprule\noalign{}
\begin{minipage}[b]{\linewidth}\raggedright
Kind
\end{minipage} & \begin{minipage}[b]{\linewidth}\raggedleft
Anchor phase
\end{minipage} & \begin{minipage}[b]{\linewidth}\raggedright
Three transported cells (phase, delta-r, delta-s)
\end{minipage} \\
\midrule\noalign{}
\endhead
\bottomrule\noalign{}
\endlastfoot
R & 0 & \path|((0, 0, 0), (1, 0, 0), (2, 0, 0))| \\
R & 1 & \path|((1, 0, 0), (2, -1, 0), (0, 0, -1))| \\
R & 2 & \path|((2, 0, 0), (0, 0, -1), (1, 1, -1))| \\
D & 0 & \path|((0, 0, 0), (2, -1, 1), (1, -1, 1))| \\
D & 1 & \path|((1, 0, 0), (0, -1, 0), (2, -2, 1))| \\
D & 2 & \path|((2, 0, 0), (1, 0, 0), (0, -1, 0))| \\
H & 0 & \path|((0, 0, 0), (1, 0, 0), (2, -1, 0))| \\
H & 1 & \path|((1, 0, 0), (2, -1, 0), (0, -1, -1))| \\
H & 2 & \path|((2, 0, 0), (0, 0, -1), (1, 0, -1))| \\
V & 0 & \path|((0, 0, 0), (2, 0, 0), (1, 1, -1))| \\
V & 1 & \path|((1, 0, 0), (0, 0, -1), (2, 0, -1))| \\
V & 2 & \path|((2, 0, 0), (1, 1, -1), (0, 1, -2))| \\
\end{longtable}

For an original offset (u,v), the output phase is q=p+u+2v mod 3. With
\path|(alpha,beta)=e_p| and \path|(gamma,delta)=e_q,| the inverse change
of indices is

\begin{verbatim}
delta-r=(beta+v-delta-alpha-u+gamma)/3,
delta-s=(-alpha-u+gamma-2(beta+v-delta))/3.
\end{verbatim}

These integers give every entry of the table. Applying Phi\_d back
returns each original offset, proving the table and its inverse on
generators. In particular all nine bone phases and all three right-stone
phases are retained; no phase restriction is asserted for arbitrary
tilings.

The induced bijections on placed-tile generators and cell generators
extend to inverse homomorphisms on their free nonnegative monoids, and
on their free integer modules. Summing the three displayed cell images
proves the incidence square commutes. The fibre of the positive
incidence map over the original region vector is therefore mapped
bijectively to the fibre over these three exact polygons. The inverse is
the original tile map, not just an equality of counts.

For finite supports lambda, send lambda to \path|Phi_d^{-1}(lambda).|
This preserves empty support and union and has the inverse Phi\_d.~The
preceding linear maps commute with inclusions on these supports.
Consequently they give the corresponding maps on the support-indexed
complexes and, via the pinned reconstruction, their Split-Zero totals.
Supported zero at lambda is sent to supported zero at its displayed
image support. The same maps induce the \path|homology/quotient|
comparisons because they carry the actual incidence boundaries to the
actual incidence boundaries. This is the exact support and coefficient
transport used here.

The d-direction also has an exact commuting identity:

\begin{verbatim}
Phi_(d+1)(r,s,p)=Phi_d(r,s,p)+(1,0).
\end{verbatim}

It identifies the particular translated narrow-shell inclusion tested
during exploration. The three polygon equations do not supply a positive
tiling of that shell; the saved failures concern that specific
frozen-source inclusion, not other maps.

\section{2. Original source tilings for the fifth
collar}\label{edition-7-section-5}

Each stage j uses the translation and released original tiles specified
in tables.json, then inserts its listed positive bone families. In that
file a family {[}name,kind,X,Y,U{]} means the anchor

\begin{verbatim}
(X0*d+X1*r+X2, Y0*d+Y1*r+Y2), 0<=r<=U0*d+U1.
\end{verbatim}

Singleton families have U=(0,0). This convention is a literal evaluation
map, with the named source generator and r retained. The tables below
include the complete preceding construction. Each released tile is
matched with an actual original generator and checked against every
intervening removal by \path|source_membership| in certify.py. The
algorithm solves the displayed affine anchor equalities exactly; it
never chooses a representative by finite sampling.

For the fifth collar translate T4 by t5=(1,1). Its original difference
increments are (0,-3,3), so the translated inner bounds are

\begin{verbatim}
D0 in [-d-11,2d+11],
D1 in [-d-14,2d+8],
D2 in [-d-8,2d+14].                         (2)
\end{verbatim}

The outer W5 bounds are {[}-d-14,2d+14{]}, proving inclusion of the
actual translated source. Release S5=R(3-d,d-1) and these 16 bones:

\begin{longtable}[]{@{}
  >{\raggedright\arraybackslash}p{(\columnwidth - 4\tabcolsep) * \real{0.3333}}
  >{\raggedright\arraybackslash}p{(\columnwidth - 4\tabcolsep) * \real{0.3333}}
  >{\raggedright\arraybackslash}p{(\columnwidth - 4\tabcolsep) * \real{0.3333}}@{}}
\toprule\noalign{}
\begin{minipage}[b]{\linewidth}\raggedright
Original kind
\end{minipage} & \begin{minipage}[b]{\linewidth}\raggedright
Anchor
\end{minipage} & \begin{minipage}[b]{\linewidth}\raggedright
Original preimage (before its later translations)
\end{minipage} \\
\midrule\noalign{}
\endhead
\bottomrule\noalign{}
\endlastfoot
R & \texttt{(-d+3,d-1)} & base q(d-3,0) \\
D & \texttt{(-d-6,d+5)} & stage 3 A\_0 \\
D & \texttt{(-d-5,d+3)} & stage 3 A\_1 \\
D & \texttt{(-d-5,d+5)} & stage 3 B \\
D & \texttt{(-d-5,d+6)} & stage 3 C \\
D & \texttt{(-d-4,d+6)} & stage 3 E \\
D & \texttt{(-d-3,d+6)} & stage 3 F \\
D & \texttt{(-d-2,d)} & stage 2 A \\
D & \texttt{(-d-2,d+1)} & stage 2 B \\
D & \texttt{(-d-1,d+5)} & stage 3 G\_0 \\
D & \texttt{(-d+1,d+4)} & stage 3 G\_1 \\
H & \texttt{(-d-3,d+2)} & stage 2 E \\
H & \texttt{(-d-1,d+1)} & stage 2 G \\
H & \texttt{(-d,d)} & stage 2 H \\
H & \texttt{(-d,d+2)} & stage 1 F\_0 \\
H & \texttt{(-d+2,d+1)} & stage 1 F\_1 \\
V & \texttt{(-d+1,d-3)} & stage 2 K\_0 \\
\end{longtable}

The base indices (d-3,0) lie in the original triangle for all
d\textgreater=4 and differ from all four previously consumed pairs. The
stage-3 preimages have subsequent total translation (-1,2); the stage-2
preimages have total translation zero; the stage-1 preimages have total
translation (1,1). Applying those shifts to their printed anchors gives
the table above. All constant indices 0,1 are within the corresponding
ranges. None equals an earlier removed tile for any integer
d\textgreater=4. The exact affine inverse witnesses and the exclusion
calculations are emitted in all-parameter-summary.json, including d=4;
no exceptional source assumption is used.

\section{3. The fifth replacement, for all
d\textgreater=4}\label{edition-7-section-6}

Insert the following bones. The column \path|inside offsets| uses the
bone's original offset u=0,1,2 and the inner bounds (2). A blank entry
means that all three cells are outside the translated source.

\begin{longtable}[]{@{}lllll@{}}
\toprule\noalign{}
Label & Kind & Anchor & r range & Inside offsets \\
\midrule\noalign{}
\endhead
\bottomrule\noalign{}
\endlastfoot
D0 & D & \path|(-d+r-8,d-2r+3)| & \texttt{0..d+3} & \texttt{{[}{]}} \\
D1 & D & \texttt{(-d-8,d+4)} & singleton & \texttt{{[}{]}} \\
D2 & D & \texttt{(-d-4,d+2)} & singleton & \texttt{{[}0,\ 1,\ 2{]}} \\
D3 & D & \texttt{(-d-3,d+2)} & singleton & \texttt{{[}0,\ 1,\ 2{]}} \\
D4 & D & \texttt{(-d-3,d+3)} & singleton & \texttt{{[}0,\ 1,\ 2{]}} \\
D5 & D & \texttt{(-d-1,d-1)} & singleton & \texttt{{[}0,\ 1,\ 2{]}} \\
D6 & D & \texttt{(-d+1,d+1)} & singleton & \texttt{{[}0,\ 1,\ 2{]}} \\
D7 & D & \texttt{(-d+1,d+3)} & singleton & \texttt{{[}0,\ 1,\ 2{]}} \\
D8 & D & \texttt{(-d+2,d+1)} & singleton & \texttt{{[}0,\ 1,\ 2{]}} \\
D9 & D & \texttt{(-d+2,d+3)} & singleton & \texttt{{[}0,\ 1,\ 2{]}} \\
H0 & H & \texttt{(-d-9,d+5)} & singleton & \texttt{{[}{]}} \\
H1 & H & \texttt{(-d-8,d+6)} & singleton & \texttt{{[}{]}} \\
H2 & H & \texttt{(-d-7,d+4)} & singleton & \texttt{{[}2{]}} \\
H3 & H & \texttt{(-d-7,d+7)} & singleton & \texttt{{[}{]}} \\
H4 & H & \texttt{(-d-6,d+3)} & singleton & \texttt{{[}1,\ 2{]}} \\
H5 & H & \texttt{(-d-6,d+5)} & singleton & \texttt{{[}0,\ 1,\ 2{]}} \\
H6 & H & \texttt{(-d-6,d+8)} & singleton & \texttt{{[}{]}} \\
H7 & H & \texttt{(-d-5,d+6)} & singleton & \texttt{{[}0,\ 1,\ 2{]}} \\
H8 & H & \texttt{(-d-5,d+9)} & singleton & \texttt{{[}{]}} \\
H9 & H & \texttt{(-d-4,d+4)} & singleton & \texttt{{[}0,\ 1,\ 2{]}} \\
H10 & H & \texttt{(-d-3,d+5)} & singleton & \texttt{{[}0,\ 1,\ 2{]}} \\
H11 & H & \texttt{(-d-1,d+2)} & singleton & \texttt{{[}0,\ 1,\ 2{]}} \\
H12 & H & \texttt{(-d-1,d+4)} & singleton & \texttt{{[}0,\ 1,\ 2{]}} \\
V0 & V & \path|(-d+r-5,d-2r)| & \texttt{0..d+2} & \texttt{{[}{]}} \\
V1 & V & \texttt{(-d,d-1)} & singleton & \texttt{{[}0,\ 1,\ 2{]}} \\
V2 & V & \texttt{(-d+1,d-2)} & singleton & \texttt{{[}0,\ 1,\ 2{]}} \\
V3 & V & \texttt{(-2,-d-5)} & singleton & \texttt{{[}{]}} \\
V4 & V & \texttt{(-1,-d-6)} & singleton & \texttt{{[}{]}} \\
V5 & V & \texttt{(0,-d-6)} & singleton & \texttt{{[}{]}} \\
V6 & V & \texttt{(1,-d-7)} & singleton & \texttt{{[}{]}} \\
V7 & V & \texttt{(2,-d-7)} & singleton & \texttt{{[}{]}} \\
V8 & V & \texttt{(3,-d-8)} & singleton & \texttt{{[}{]}} \\
V9 & V & \texttt{(4,-d-8)} & singleton & \texttt{{[}{]}} \\
V10 & V & \path|(r+5,-d+r-9)| & \texttt{0..d+4} & \texttt{{[}{]}} \\
\end{longtable}

\subsection{3.1 Containment and exact
intersection}\label{edition-7-section-7}

Substitute every anchor plus its three original offsets into D0,D1,D2.
Each bound is affine in d and r. For an expression Ad+Br+C on
0\textless=r\textless=ud+v, its minimum in r is attained at r=0 for
B\textgreater=0 and at r=ud+v for B\textless0. This produces an affine
expression ad+c.~Its nonnegativity for d\textgreater=4 follows from
a\textgreater=0 and 4a+c\textgreater=0. The exact verifier applies these
comparisons to all 612 outer one-sided inequalities.

For every offset marked inside, the same calculation proves all six
inequalities in (2). For every offset marked outside, the verifier
supplies one fixed coordinate inequality violating (2) by at least one,
valid on the complete parameter range. There is no unresolved cutoff or
omitted residue class. The 51 inside cells are compared as a multiset of
affine coordinate pairs in d with the cells of the 17 released tiles.
They agree identically.

For example, D0\_r lies wholly outside the inner D0 or D2 bounds as
certified by those literal expressions, while H2 has only offset 2
inside and H4 has offsets 1,2 inside. The general verifier, not this
illustrative sentence, prints the exact offset table above. The right
stone and all 16 source bones are retained in the 51-cell identity.

\subsection{3.2 Every collision is ruled out on the unbounded
domain}\label{edition-7-section-8}

Cells from a single indexed family have injective (r,u) coordinates: its
anchor step and original bone direction have nonzero determinant.
Singleton bones have three distinct cells.

For two different families and offsets u,v, introduce real variables
(d,r,s). Include d\textgreater=4, both exact index ranges, and the two
coordinate equalities. Encode the latter as pairs of opposite
inequalities. There are exactly nine inequalities in each system and
5,049 such systems for the 34 fifth-collar families.

The file all-parameter-certificates.json, regenerated by certify.py,
gives a nonnegative rational combination of each system whose left side
is zero and whose right side is strictly negative. The acceptance
function recomputes the nine constraints from the original table and
checks every coefficient sum using exact fractions. Thus a simultaneous
collision would give 0\textless0. This rules it out for all real
solutions of those bounds, hence for the original integer indices. The
finite d replay is not used for this conclusion.

The generator affine.py uses exact elimination to discover the
combinations. The acceptance arithmetic does not rely on a solver
status, a tolerance, or agreement among models. Every family pair and
all nine pairs of offsets are enumerated directly, so there is no
unlisted collision system. Earlier collars are checked by the identical
argument on their own domains; the complete numbers are given in section
6.

\subsection{3.3 The positive partition and
tiling}\label{edition-7-section-9}

The replacement contains (d+13) diagonal bones, 13 horizontal bones and
(2d+17) vertical bones, totalling 3d+43. Its disjoint cells therefore
number 9d+129. Formula (1) gives

\begin{verbatim}
|W5|-|W4|=9d+78.
\end{verbatim}

The source intersection contains exactly 51 cells. Containment,
disjointness, exact source intersection and cardinality prove

\begin{verbatim}
covered(C5)=(W5 \ ((1,1)+W4)) disjoint-union covered(S5+eta5),
\end{verbatim}

where eta5 is the sum of the 16 original released bone generators. Hence

\begin{verbatim}
T5=(((1,1)+T4) \ {S5 and the 16 released bones}) union C5
\end{verbatim}

is an original type-103 tiling for every integer d\textgreater=4. Its
incidence identity in the free integer cell module is

\begin{verbatim}
partial beta5 = 1_W5 - J5 1_W4 + partial S5 + partial eta5.       (3)
\end{verbatim}

Every coefficient of beta5 is +1. T4 has binom(d,2)-4 right stones and
12d+48 bones; deleting one stone and 16 bones, then adding 3d+43 bones,
gives the announced counts. This proves the fifth-family result.

The involution (x,y)-\textgreater(x,1-x-y) sends (D0,D1,D2) to
(-D2,-D1,-D0), hence sends V(a,b) onto V(b,a). It preserves right
stones, interchanges H and D bones and reverses the V direction;
re-anchoring the reversed bone at its original opposite endpoint gives
the printed positive orientation. Applying the involution twice is the
identity. This proves the reflected family and its inverse tile
transport.

\section{4. The exact positive-fibre bijection and cohomology
relation}\label{edition-7-section-10}

Let A5 be the set of 17 released tiles in the translated source and B5
the 3d+43 replacement bones. Let U be the set of all original W4 tilings
whose translates contain A5. Let V be the set of all original W5 tilings
containing B5. Both sets have literal definitions on the original tiling
carrier.

Define

\begin{verbatim}
F(T)=(J5 T \ A5) union B5,
F^{-1}(T')=J5^{-1}((T' \ B5) union A5).
\end{verbatim}

The partition identity in section 3 proves both functions are defined.
The remaining cells of a target tiling containing B5 are precisely
J5(W4) minus the released cells, so every remaining tile lies entirely
in J5(W4). Removing B5, reinserting A5 and translating back therefore
gives an original source tiling. Both compositions remove and restore
the same disjoint tile sets and equal the identity. T4 supplies a member
of U. This proves an actual bijection of the specified positive tiling
fibres, not an assertion that F covers all W5 tilings.

On tile vectors F is the affine map J5 c-A5+B5. Differences satisfy

\begin{verbatim}
F(c)-F(c')=J5(c-c').
\end{verbatim}

The linear cochain map is J5, with partial J5=J5 partial on each
original generator. Equation (3) specifies the fixed affine correction.
Thus the affine positive-fibre map and the linear cochain map are
related by these exact formulas; neither is substituted for the other.

In the finite-support bone quotient of W5, the original representative
map gives

\begin{verbatim}
[1_W5]-J5[1_W4]=-[S5],
\end{verbatim}

since beta5 and eta5 are bone boundaries. The earlier integral global
quotient Q and its finite-jet observation receive this equality through
the cell inclusions. The source stone has the retained order-three
global class, and the correction is realized inside the specified
support by the positive partition. The observation is used together with
the original representatives and their support; no vanishing scalar
observation is treated as a positive source section.

The same formula applies on the support-indexed complexes under the
pinned Split-Zero reconstruction: inclusion of supports and translation
act on the original generators, and quotient maps take a boundary to its
support-labelled zero. The positive fibre map is the displayed affine
correction inside those same cell modules.

\section{5. What was additionally checked, and what remains
next}\label{edition-7-section-11}

Five sixth-collar positive certificates are saved for d=4,5,6,7,8. Each
names the actual translated T5 source, its removed stone and old bones,
and every new bone. The constructors for h\textless=5 contain no search.
The h=6 data are bounded integer certificates, not an all-parameter
sixth-collar formula. Solver optimality is not used or claimed. The
d=4,h=6 certificate reaches the known zero-invariant endpoint for that
fixed d.

Three bone-only annuli are also saved, for source
\path|(d,h)=(3,0),(3,1),(4,0),| using the literal zero translation into
W\_(h+d)(d+1). The direct calculation

\begin{verbatim}
binom(d+1,2)-(h+d)=binom(d,2)-h
\end{verbatim}

preserves the stone-count invariant, and (1) gives annulus area
9(2d+1)(d+h). Every saved annulus is checked as a positive bone
partition of exactly the original set difference. A uniform annulus
formula has not been proved. Timed searches on larger cases are recorded
as uncompleted searches, not infeasibility results.

The three-polygon and all-phase tile maps in section 1 provide a
concrete next calculation in both parameters. The next work is to
construct a parameter-uniform positive lift through the remaining h
values, using the actual transported incidence map, and to return that
lift by Phi\_d.~The exact sixth-collar and constant-charge annuli are
starting configurations. This is a continuation task, not a stated
conditional theorem or a completed full P7 argument.

\section{6. Replay and all-parameter source
closure}\label{edition-7-section-12}

Run \path|python check.py --max-d 100 --output evidence/normal.json| and
the corresponding
\path|python -O check.py --max-d 100 --output evidence/optimized.json|.
Both modes execute the exact all-parameter certificates, original source
membership calculations, finite tilings and reflection checks, positive
forward/inverse fibre maps, polygon inverse and tile-map tests, saved
finite certificates and rejection tests.

The five stages have 189, 495, 702, 2484 and 5049 collision
contradictions, respectively. The stage-3 construction holds for
d\textgreater=3; stages 4 and 5 hold for d\textgreater=4. All source
membership and range checks use these unbounded domains.

\subsection{Complete prior-stage tables}\label{edition-7-section-13}

For each family below the anchor is (X0*d+X1*r+X2, Y0*d+Y1*r+Y2),
0\textless=r\textless=U0*d+U1. These are the unchanged mathematical
anchors from the preceding constructions. The explicit base tiling in
section 1 and these tables supply the entire mathematical source used by
the fifth constructor.

\subsubsection{Stage 1: d\textgreater=3, translation (-2,
1)}\label{edition-7-section-14}

Released original generators:

\begin{verbatim}
R(-2, -d+3)
\end{verbatim}

\begin{longtable}[]{@{}lllll@{}}
\toprule\noalign{}
Family & Kind & X & Y & U \\
\midrule\noalign{}
\endhead
\bottomrule\noalign{}
\endlastfoot
A & D & \texttt{{[}0,\ 0,\ -2{]}} & \texttt{{[}-1,\ 0,\ 3{]}} &
\texttt{{[}0,\ 0{]}} \\
B & D & \texttt{{[}0,\ 0,\ -2{]}} & \texttt{{[}-1,\ 0,\ 4{]}} &
\texttt{{[}0,\ 0{]}} \\
C & D & \texttt{{[}0,\ 0,\ -1{]}} & \texttt{{[}-1,\ 0,\ 1{]}} &
\texttt{{[}0,\ 0{]}} \\
F & H & \texttt{{[}-1,\ 2,\ -1{]}} & \texttt{{[}1,\ -1,\ 1{]}} &
\texttt{{[}1,\ -1{]}} \\
G & H & \texttt{{[}0,\ 1,\ 0{]}} & \texttt{{[}-1,\ 1,\ 3{]}} &
\texttt{{[}1,\ -2{]}} \\
E & V & \texttt{{[}0,\ 0,\ 1{]}} & \texttt{{[}-1,\ 0,\ 0{]}} &
\texttt{{[}0,\ 0{]}} \\
J & V & \texttt{{[}0,\ 1,\ 2{]}} & \texttt{{[}-1,\ 1,\ 0{]}} &
\texttt{{[}1,\ -1{]}} \\
\end{longtable}

\subsubsection{Stage 2: d\textgreater=3, translation (1,
1)}\label{edition-7-section-15}

Released original generators:

\begin{verbatim}
R(-d+1, d)
\end{verbatim}

\begin{longtable}[]{@{}lllll@{}}
\toprule\noalign{}
Family & Kind & X & Y & U \\
\midrule\noalign{}
\endhead
\bottomrule\noalign{}
\endlastfoot
A & D & \texttt{{[}-1,\ 0,\ -2{]}} & \texttt{{[}1,\ 0,\ 0{]}} &
\texttt{{[}0,\ 0{]}} \\
B & D & \texttt{{[}-1,\ 0,\ -2{]}} & \texttt{{[}1,\ 0,\ 1{]}} &
\texttt{{[}0,\ 0{]}} \\
R & D & \texttt{{[}-1,\ 1,\ -1{]}} & \texttt{{[}1,\ -2,\ -2{]}} &
\texttt{{[}1,\ -1{]}} \\
E & H & \texttt{{[}-1,\ 0,\ -3{]}} & \texttt{{[}1,\ 0,\ 2{]}} &
\texttt{{[}0,\ 0{]}} \\
F & H & \texttt{{[}-1,\ 0,\ -2{]}} & \texttt{{[}1,\ 0,\ 3{]}} &
\texttt{{[}0,\ 0{]}} \\
G & H & \texttt{{[}-1,\ 0,\ -1{]}} & \texttt{{[}1,\ 0,\ 1{]}} &
\texttt{{[}0,\ 0{]}} \\
H & H & \texttt{{[}-1,\ 0,\ 0{]}} & \texttt{{[}1,\ 0,\ 0{]}} &
\texttt{{[}0,\ 0{]}} \\
K & V & \texttt{{[}-1,\ 1,\ 1{]}} & \texttt{{[}1,\ -2,\ -3{]}} &
\texttt{{[}1,\ -1{]}} \\
L & V & \texttt{{[}0,\ 0,\ 1{]}} & \texttt{{[}-1,\ 0,\ -2{]}} &
\texttt{{[}0,\ 0{]}} \\
M & V & \texttt{{[}0,\ 0,\ 2{]}} & \texttt{{[}-1,\ 0,\ -3{]}} &
\texttt{{[}0,\ 0{]}} \\
N & V & \texttt{{[}0,\ 1,\ 3{]}} & \texttt{{[}-1,\ 1,\ -2{]}} &
\texttt{{[}1,\ 0{]}} \\
\end{longtable}

\subsubsection{Stage 3: d\textgreater=3, translation (1,
-2)}\label{edition-7-section-16}

Released original generators:

\begin{verbatim}
R(d-2, 0)
H(d-1, 1)
H(d, 0)
\end{verbatim}

\begin{longtable}[]{@{}lllll@{}}
\toprule\noalign{}
Family & Kind & X & Y & U \\
\midrule\noalign{}
\endhead
\bottomrule\noalign{}
\endlastfoot
A & D & \texttt{{[}-1,\ 1,\ -5{]}} & \texttt{{[}1,\ -2,\ 3{]}} &
\texttt{{[}1,\ 2{]}} \\
B & D & \texttt{{[}-1,\ 0,\ -4{]}} & \texttt{{[}1,\ 0,\ 3{]}} &
\texttt{{[}0,\ 0{]}} \\
C & D & \texttt{{[}-1,\ 0,\ -4{]}} & \texttt{{[}1,\ 0,\ 4{]}} &
\texttt{{[}0,\ 0{]}} \\
E & D & \texttt{{[}-1,\ 0,\ -3{]}} & \texttt{{[}1,\ 0,\ 4{]}} &
\texttt{{[}0,\ 0{]}} \\
F & D & \texttt{{[}-1,\ 0,\ -2{]}} & \texttt{{[}1,\ 0,\ 4{]}} &
\texttt{{[}0,\ 0{]}} \\
G & D & \texttt{{[}-1,\ 2,\ 0{]}} & \texttt{{[}1,\ -1,\ 3{]}} &
\texttt{{[}1,\ -1{]}} \\
J & H & \texttt{{[}-1,\ 2,\ -3{]}} & \texttt{{[}1,\ -1,\ 5{]}} &
\texttt{{[}1,\ 1{]}} \\
K & H & \texttt{{[}1,\ 0,\ 0{]}} & \texttt{{[}0,\ 0,\ 3{]}} &
\texttt{{[}0,\ 0{]}} \\
L & H & \texttt{{[}1,\ 0,\ -2{]}} & \texttt{{[}0,\ 0,\ 0{]}} &
\texttt{{[}0,\ 0{]}} \\
M & H & \texttt{{[}1,\ 0,\ -2{]}} & \texttt{{[}0,\ 0,\ 1{]}} &
\texttt{{[}0,\ 0{]}} \\
P & V & \texttt{{[}1,\ 0,\ 1{]}} & \texttt{{[}0,\ 0,\ 0{]}} &
\texttt{{[}0,\ 0{]}} \\
Q & V & \texttt{{[}1,\ 0,\ 2{]}} & \texttt{{[}0,\ 0,\ 0{]}} &
\texttt{{[}0,\ 0{]}} \\
N & V & \texttt{{[}1,\ 1,\ 3{]}} & \texttt{{[}0,\ -2,\ 1{]}} &
\texttt{{[}0,\ 2{]}} \\
\end{longtable}

\subsubsection{Stage 4: d\textgreater=4, translation (-2,
1)}\label{edition-7-section-17}

Released original generators:

\begin{verbatim}
R(-1, -d+4)
D(-1, -d+1)
H(0, -d+3)
V(-1, -d-2)
V(0, -d-3)
V(1, -d-4)
V(1, -d)
V(2, -d-3)
V(2, -d)
\end{verbatim}

\begin{longtable}[]{@{}lllll@{}}
\toprule\noalign{}
Family & Kind & X & Y & U \\
\midrule\noalign{}
\endhead
\bottomrule\noalign{}
\endlastfoot
A & D & \texttt{{[}0,\ 0,\ -4{]}} & \texttt{{[}-1,\ 0,\ -2{]}} &
\texttt{{[}0,\ 0{]}} \\
B & D & \texttt{{[}0,\ 0,\ -2{]}} & \texttt{{[}-1,\ 0,\ -3{]}} &
\texttt{{[}0,\ 0{]}} \\
C & D & \texttt{{[}0,\ 0,\ -1{]}} & \texttt{{[}-1,\ 0,\ -2{]}} &
\texttt{{[}0,\ 0{]}} \\
D & D & \texttt{{[}0,\ 0,\ -1{]}} & \texttt{{[}-1,\ 0,\ 4{]}} &
\texttt{{[}0,\ 0{]}} \\
E & D & \texttt{{[}0,\ 0,\ -1{]}} & \texttt{{[}-1,\ 0,\ 5{]}} &
\texttt{{[}0,\ 0{]}} \\
F & D & \texttt{{[}0,\ 0,\ 0{]}} & \texttt{{[}-1,\ 0,\ -4{]}} &
\texttt{{[}0,\ 0{]}} \\
G & D & \texttt{{[}0,\ 0,\ 1{]}} & \texttt{{[}-1,\ 0,\ -3{]}} &
\texttt{{[}0,\ 0{]}} \\
H & D & \texttt{{[}0,\ 0,\ 1{]}} & \texttt{{[}-1,\ 0,\ -2{]}} &
\texttt{{[}0,\ 0{]}} \\
I & D & \texttt{{[}0,\ 0,\ 2{]}} & \texttt{{[}-1,\ 0,\ -5{]}} &
\texttt{{[}0,\ 0{]}} \\
J & H & \texttt{{[}-1,\ 2,\ -4{]}} & \texttt{{[}1,\ -1,\ 7{]}} &
\texttt{{[}1,\ 2{]}} \\
K & H & \texttt{{[}0,\ 1,\ 3{]}} & \texttt{{[}-1,\ 1,\ -3{]}} &
\texttt{{[}1,\ 1{]}} \\
L & H & \texttt{{[}1,\ 0,\ 2{]}} & \texttt{{[}0,\ 0,\ 3{]}} &
\texttt{{[}0,\ 0{]}} \\
M & H & \texttt{{[}1,\ 0,\ 2{]}} & \texttt{{[}0,\ 0,\ 4{]}} &
\texttt{{[}0,\ 0{]}} \\
N & H & \texttt{{[}1,\ 0,\ 3{]}} & \texttt{{[}0,\ 0,\ 1{]}} &
\texttt{{[}0,\ 0{]}} \\
O & H & \texttt{{[}1,\ 0,\ 3{]}} & \texttt{{[}0,\ 0,\ 2{]}} &
\texttt{{[}0,\ 0{]}} \\
P & H & \texttt{{[}1,\ 0,\ 4{]}} & \texttt{{[}0,\ 0,\ -1{]}} &
\texttt{{[}0,\ 0{]}} \\
Q & H & \texttt{{[}1,\ 0,\ 4{]}} & \texttt{{[}0,\ 0,\ 0{]}} &
\texttt{{[}0,\ 0{]}} \\
R & V & \texttt{{[}0,\ 0,\ -1{]}} & \texttt{{[}-1,\ 0,\ -1{]}} &
\texttt{{[}0,\ 0{]}} \\
S & V & \texttt{{[}0,\ 0,\ 0{]}} & \texttt{{[}-1,\ 0,\ -2{]}} &
\texttt{{[}0,\ 0{]}} \\
T & V & \texttt{{[}0,\ 0,\ 1{]}} & \texttt{{[}-1,\ 0,\ -1{]}} &
\texttt{{[}0,\ 0{]}} \\
U & V & \texttt{{[}0,\ 0,\ 2{]}} & \texttt{{[}-1,\ 0,\ -2{]}} &
\texttt{{[}0,\ 0{]}} \\
V & V & \texttt{{[}0,\ 0,\ 2{]}} & \texttt{{[}-1,\ 0,\ 1{]}} &
\texttt{{[}0,\ 0{]}} \\
W & V & \texttt{{[}0,\ 0,\ 4{]}} & \texttt{{[}-1,\ 0,\ -6{]}} &
\texttt{{[}0,\ 0{]}} \\
Z & V & \texttt{{[}0,\ 1,\ 5{]}} & \texttt{{[}-1,\ 1,\ -6{]}} &
\texttt{{[}1,\ 2{]}} \\
\end{longtable}

\chapter{Turn 4: actual Split-Zero complexes}\label{edition-8}

Source: \path|sources/benzel-p7-turn4/benzel-p7-turn4/turn4/PROOF.md|.

SHA-256:
\path|b35395dcecf2d5d60739bd8c41e3a9fbda0dcff506a007dcba9997f197034bbe|.

15 September 2026. Parent:
\path|d7cd069031536bea677ed31f7729ebdd4e9f97f6|, the fifth-collar
contribution in PR 24 of \path|KokunoYumeto/mathematics-commons-pilot|.

This continuation proves an all-parameter support statement for the
original fifth repair. It supplies a genuine support-indexed augmented
complex, computes its distinguished integral homology transitions and
their Rees defect, and gives exact dual certificates for its positive
fibres. It does not claim a sixth-family construction, a full resolution
of Propp 7, independent specialist acceptance, Lean verification, or
priority.

\section{0. Source mathematics actually used}\label{edition-8-section-1}

The coefficient and reconstruction sources are the owner's
\path|KokunoYumeto/zeta-function-research-reader|, pinned at
\path|1c8ec52c85c173adc9f8403a8a26914955e2f5a9|:

\begin{itemize}
\tightlist
\item
  \path|workbenches/splitzero-tandem/tex/split_zero_carriers.tex|,
  scalar carrier and ring reflection.
\item
  \path|workbenches/splitzero-tandem/tex/support_diagrams.tex|, D1-D8:
  reconstruction, changing support indices, coequalizer, comparison
  kernel, and the inclusion of the global-zero kernel into the
  supported-zero preimage.
\item
  \path|formal/splitzero/DERIVED_MATHEMATICS.md|, sections 1-6, and
  \path|formal/splitzero/SplitZeroComplex.lean|, especially
  \path|differential_square|, \path|cycle_condition|,
  \path|homology_coequalizes|, and \path|homology_coequalizer|.
\item
  \path|workbenches/split-support-rees-trace/RESEARCH_NOTE.md|, sections
  2-5: the two filtrations on the actual comparison image, their Rees
  inclusion, residue pairing, amplification, and boundary character.
\end{itemize}

The current guide was also read at
\path|ba360f08f3417464ca915233e3d3b54f5b76b30c|. Its same-source
\path|metric/residual-observation| programme is not silently transferred
to tilings. The core D1-D8 file has the identical blob
\path|d3493f891291ee6e94dbf2c77649f7d85d240df2| at both revisions. The
new 132-page analytic continuation was not independently audited or
rebuilt here. Full source URLs and reading scope are in
\path|SOURCES.json|.

The new constructions below are specializations and calculations on the
original benzel complexes, not additional assertions made by those
upstream sources. In particular, the positive-fibre certificate is
proved below, rather than attributed to an upstream theorem.

\section{1. Fixed original supports and the sixteen released
generators}\label{edition-8-section-2}

Throughout, d is an integer at least 4. Use the original cells
\texttt{(x,y)} with

\begin{verbatim}
y-x, 1-x-2y, 2x+y-1 in [1-a,b-1],
W_h(d)=V(d+3h,2d+3h).
\end{verbatim}

The bone anchors are \path|H(x,y)={(x,y),(x+1,y),(x+2,y)},|
\path|V(x,y)={(x,y),(x,y+1),(x,y+2)},|
\path|D(x,y)={(x,y),(x+1,y-1),(x+2,y-2)};| S(x,y) is the right stone
\{(x,y),(x+1,y),(x,y+1)\}.

The source is the parent's actual T4(d), translated by
\path|J(x,y)=(x+1,y+1).| The parent proves it tiles J W4(d) inside
W5(d). The source stone is S\_d=S(3-d,d-1). Write

\begin{verbatim}
K_d=W5(d) \ J W4(d),
Lambda_empty=K_d disjoint-union cells(S_d).
\end{verbatim}

For compact tables only, write a tile with offset anchor (u,v) for its
original anchor (u-d,v+d). The map \path|J_d(u,v)=(u-d,v+d)| has inverse
(x,y)-\textgreater(x+d,y-d). It transports every displayed tile by
translation; neither the region nor the coefficient of a cell is
changed.

The parent's sixteen released old bones B1,\ldots,B16 have these offset
anchors, in their original order:

\begin{longtable}[]{@{}rlrr@{}}
\toprule\noalign{}
i & kind & u & v \\
\midrule\noalign{}
\endhead
\bottomrule\noalign{}
\endlastfoot
1 & D & -6 & 5 \\
2 & D & -5 & 3 \\
3 & D & -5 & 5 \\
4 & D & -5 & 6 \\
5 & D & -4 & 6 \\
6 & D & -3 & 6 \\
7 & D & -2 & 0 \\
8 & D & -2 & 1 \\
9 & D & -1 & 5 \\
10 & D & 1 & 4 \\
11 & H & -3 & 2 \\
12 & H & -1 & 1 \\
13 & H & 0 & 0 \\
14 & H & 0 & 2 \\
15 & H & 2 & 1 \\
16 & V & 1 & -3 \\
\end{longtable}

For A subset \{1,\ldots,16\}, let

\begin{verbatim}
Lambda_A = Lambda_empty disjoint-union (union of cells(Bi), i in A),
v_A = sum of the original cell generators in Lambda_A.
\end{verbatim}

These unions are disjoint because the stone and bones are actual
distinct source tiles. Parent source membership is retained by exact
blob checks in \texttt{support.py}. The original cell count is

\begin{verbatim}
|Lambda_A|=9d+81+3|A|.
\end{verbatim}

All placed bones contained in Lambda\_A are permitted, not only the
bones in a proposed replacement table. Let C\_A\^{}0 be the free integer
module on those placed bones, C\_A\^{}1 the free integer module on
Lambda\_A, and C\_A\^{}2=0. The original differential is incidence:
\path|partial[b]=sum_{c| in b\}{[}c{]}. Write B\_A=im partial and
H\_A=C\_A\^{}1/B\_A.

For A subset B, let j\_AB be inclusion of the original generators and
eta\_AB=sum\_\{i in B\textbackslash A\}{[}Bi{]}. Then

\begin{verbatim}
v_B = j_AB v_A + partial eta_AB,
eta_AC = j_BC eta_AB + eta_BC.
\end{verbatim}

These are equalities of original cochains. In particular {[}v\_A{]}
transports to {[}v\_B{]}. The affine map on positive fillings is n
-\textgreater{} j\_AB n + eta\_AB, with the original incidence equality.
On the target fillings containing the specified old bones eta\_AB with
coefficient one, \path|restriction/removal| gives its inverse. The
inverse follows from the disjoint source-cell partition and
nonnegativity, not from cancellation of cell counts alone.

\section{2. The actual Split-Zero objects, including the
affine-to-linear map}\label{edition-8-section-3}

Take G(Z)=Z disjoint-union \{tau\}, with supported ring zero e. Its pair
presentation is tau-\textgreater(0,0), n-\textgreater(n,1), onto
\{(0,0)\} union (Z x \{1\}); componentwise arithmetic with Boolean
support proves the operations and inverse. Thus e and tau have the same
integer reflection but different support characters.

Use the join-semilattice

\begin{verbatim}
L={bottom} disjoint-union P({1,...,16}).
\end{verbatim}

Bottom represents empty CELL support. The empty RELEASE SET represents
Lambda\_empty, which is nonempty, and is a different label. The join of
supported labels is union. The map bottom-\textgreater empty cell set,
A-\textgreater Lambda\_A preserves bottom and joins and is injective;
its inverse on the image recovers exactly which disjoint old bone cells
were added.

At bottom put C\^{}i\_bottom=0. Reconstruct each degree by the upstream
D2 formula:

\begin{verbatim}
M^i = disjoint union_{A in L} C_A^i,
(A,x)+(B,y)=(A union B, j x+j y),
e(A,x)=(A,0), tau(A,x)=(bottom,0).
\end{verbatim}

The incidence maps commute with every j\_AB and therefore give
G(Z)-linear maps. The next differential has value (A,0) in degree 2.
Hence

\begin{verbatim}
d^1 d^0(A,n)=(A,0)=z_L(A,n).
\end{verbatim}

Both sides have the original source and target degrees. At a nonbottom
label this value is not (bottom,0), as is proved by the label
projection. The cycle equalizer d(x)=e d(x) gives all original
degree-one cell vectors here, as the outgoing amplitude is zero.

The projection q(A,p)=(A,{[}p{]}) is the coequalizer of
(A,n)-\textgreater(A,partial n) and (A,n)-\textgreater(A,0). To check
the universal property directly, let f equalize those two actual maps
into any G(Z)-semimodule N. Two representatives p,p' at A differing by
partial n have

\begin{verbatim}
f(A,p)=f(A,p')+f(A,partial n)
      =f(A,p')+f(A,0)=f(A,p').
\end{verbatim}

The last identity is the image of the original fibre-zero absorption.
Thus g(A,{[}p{]})=f(A,p) is well-defined. Lifting representatives proves
addition and scalar compatibility, and surjectivity of q proves
uniqueness. No cancellation in N is used.

Here is a linear complex which also retains the AFFINE source correction
and the positive fibre. At each supported label define

\begin{verbatim}
C_hat_A^0 = Z direct-sum C_A^0,
d_hat_A(r,n)=partial n-r v_A,
C_hat_A^1=C_A^1, C_hat_A^2=0.
\end{verbatim}

At bottom all three groups are zero. The actual transition is

\begin{verbatim}
j_hat_AB^0(r,n)=(r,j_AB n+r eta_AB),
j_hat_AB^1=j_AB.
\end{verbatim}

The two cochain identities above prove

\begin{verbatim}
d_hat_B j_hat_AB^0 = j_AB d_hat_A,
j_hat_BC^0 j_hat_AB^0 = j_hat_AC^0.
\end{verbatim}

Thus the upstream reconstruction applies to this actual diagram of
complexes. Addition transports both summands by these arrows before
adding; merely adding their tile coordinates at a union label would omit
the charge-dependent corrections.

Its charge-one nonnegative integral cycles are EXACTLY the original
patch tilings: d\_hat\_A(1,n)=0 says partial n=v\_A; integer nonnegative
counts and the all-ones target force each cell to be covered once.
Conversely every patch tiling supplies such a cycle. The inclusion of
these positive cycles into the integral cycle fibre is the identity on
all charge and tile coordinates. The calculation below supplies a signed
cycle outside that positive fibre, together with the exact dual
certificate and the later positive cycle. This supplies the relationship
rather than assuming either fibre determines the other.

\texttt{support.py} implements these reconstructed cochains, augmented
arrows, support-preserving zero action, and the computed cyclic homology
subdiagram. No Lean build is claimed.

\section{3. A signed integral lift at the fifteenth support, for every
d\textgreater=4}\label{edition-8-section-4}

Let A\_j=\{1,\ldots,j\}. Denote by beta the parent's proved positive
fifth replacement on Lambda\_A16. It consists of 3d+43 original bones.
The following nine bones rho are literal fixed rows
D2,D3,D4,D5,D6,D8,H11,V1,V2 of that table:

\begin{verbatim}
D(-4,2), D(-3,2), D(-3,3), D(-1,-1), D(1,1),
D(2,1), H(-1,2), V(0,-1), V(1,-2).
\end{verbatim}

Anchors in this section use the explicit translation J\_d of section 1.
Their cells include all three cells of B16=V(1,-3). The union of their
cells, minus B16, has 24 cells.

The following signed chain mu uses only bones within that 24-cell set:

\begin{verbatim}
+H(-4,2) +V(-3,1) -D(-3,2)
+D(-2,0) +D(-2,1) +H(-2,1)
+H(-1,0) +H(-1,2) +D(1,1) +D(2,1).
\end{verbatim}

Expansion of every three-cell generator gives

\begin{verbatim}
partial mu = partial rho - partial B16.                 (1)
\end{verbatim}

This is a single finite integer identity, retained in
\path|signed_patch.json|; translation J\_d proves it for every d.~No
coefficient field is changed. Therefore

\begin{verbatim}
gamma15=beta-rho+mu,
partial gamma15=v_A15.                                 (2)
\end{verbatim}

All bones of beta-rho avoid B16's cells, by the parent's disjoint
partition and the inclusion of those cells in rho. Every bone of mu is
in Lambda\_A15, verified on the unbounded original domain in section 4.
The coefficient of D(-d-3,d+2) is -1, and all other coefficients are
nonnegative. In particular gamma15 is an actual integral signed lift in
this original support.

Let w=v\_empty, included by the actual cell maps in every larger
support. Then

\begin{verbatim}
a15=gamma15 - sum_{i=1}^{15}[Bi],
partial a15=w.                                         (3)
\end{verbatim}

This fixed-vector preimage will determine the actual image filtration in
section 6. The original all-ones target v\_A15, its correction to w, and
all negative coefficients remain available.

\section{4. Sixteen exact duals exclude every proper positive
support}\label{edition-8-section-5}

For each i, \path|dual_certificates.json| supplies an integer cochain
y\_i on \path|Lambda_{A16{i}}.| Its weights are at cells J\_d(u,v). The
complete coefficients are part of this proof, not a request for a future
separation argument. They satisfy, for every integer d\textgreater=4,

\begin{verbatim}
partial^* y_i(b)=sum_{c in b} y_i(c) >= 0
      for EVERY placed bone b contained in that support,
y_i(v_{A16\{i}})<0.                                    (4)
\end{verbatim}

Here partial\^{}* is the transpose of the original incidence map into
the integer dual. For example y\_15 is -1 at the single cell J\_d(4,-1)
and zero elsewhere. All nine possible placed bones through that cell
leave \path|Lambda_{A16{15}}.|

The complete y\_16 is given by the following offset/weight pairs:

\begin{verbatim}
(-1,0):1, (-1,2):2, (0,-1):-1, (0,0):1, (0,2):-1,
(1,0):-2, (1,1):3, (1,2):-1, (1,3):1,
(2,0):1, (2,1):3, (2,2):2,
(3,-1):-4, (3,0):7, (3,1):-3, (4,-1):-10.
\end{verbatim}

Their sum is -1. On the signed chain mu, every positive term evaluates
to zero; the negative term D(-3,2) evaluates to +1. Thus y\_16(partial
mu)=-1 by the same original incidence pairing.

For reference, the other exact negative evaluations, in order
i=1,\ldots,16, are

\begin{verbatim}
-9,-1,-2,-2,-2,-2,-1,-1,-1,-1,-1,-1,-1,-1,-1,-1.
\end{verbatim}

\subsection{Complete verification over the unbounded parameter
range}\label{edition-8-section-6}

The checker never assumes that a sample of d determines the answer. Each
cochain has finitely many negative sites. Every bone of negative total
weight contains a negative site, so it occurs among the three
orientations and three anchor offsets through such a site. The checker
enumerates precisely those candidates.

For each candidate it must show that some cell is outside the actual
support. At an affine cell J\_d(u,v), membership in W5 and J W4 consists
of the six original affine inequalities. Membership in a released
original tile is a pair of affine equalities. Collect every nonconstant
affine expression ad+b occurring in these tests. The integer interval
endpoints floor(-b/a) and floor(-b/a)+1 separate every possible change
of sign or equality. On each resulting interval, all membership truth
values are constant. The last interval is unbounded. The verifier checks
this sign stability exactly, then checks membership using the original
inequalities, not a surrogate shape.

For the actual sixteen certificates, every test is already stable on the
single interval {[}4,infinity). There are respectively

\begin{verbatim}
51,35,41,32,27,29,33,23,21,15,18,12,12,13,9,28
\end{verbatim}

negative-weight candidate bones, and each is excluded by those original
membership tests. All nonzero weighted sites remain inside the claimed
support. All remaining contained bones have nonnegative sum by the
exhaustive negative-site argument. The same procedure proves containment
of every cell of mu on {[}4,infinity). The fixed-table matches for rho
are equality checks on the original affine generator formulas. Together
these finite exact checks prove (1)-(4) on the stated infinite domain.

\subsection{Conclusion on the positive
fibres}\label{edition-8-section-7}

A nonnegative real chain n with partial n=v\_\{A16\{i\}\} would satisfy

\begin{verbatim}
y_i(v)=<partial^*y_i,n> >= 0,
\end{verbatim}

contradicting its displayed negative integer value. Thus that fibre is
empty even over nonnegative reals.

For any proper A choose i outside A. A positive filling at A would
extend to A16\{i\} by adding the actual disjoint old bones with indices
in (A16\{i\})\textbackslash A, using the affine map of section 1. This
is impossible by (4). Consequently

\begin{verbatim}
{A : the original nonnegative real filling fibre over v_A is nonempty}
    = {A16}.                                           (5)
\end{verbatim}

At A16 the parent's beta is a positive integral filling. This excludes
all 65,535 proper subsets simultaneously, for every d\textgreater=4. It
proves inclusion-minimality WITHIN THE DISPLAYED sixteen-bone release
dictionary for the fixed source tiling, source stone and translation.
Alternative source tilings, stones, translations, or released bones
outside that dictionary remain represented by other supports; no global
minimum over those choices is asserted.

The map from positive integral cycles into all integral cycles in
section 2 has empty charge-one source at A15 and contains the explicit
charge-one signed target (1,gamma15). At A16 it contains (1,beta). Under
the original augmented arrow, (1,gamma15) goes to (1,gamma15+B16). Its
difference from (1,beta) is the explicit incidence-kernel vector
rho-mu-B16. This gives all the maps between the two stages and the exact
positive repair, not a categorical claim of unrelated objects.

\section{5. The distinguished internal homology class and its comparison
kernel}\label{edition-8-section-8}

Put \path|p_d=J_d(4,-1)=(4-d,d-1),| an actual cell of the removed source
stone. For every prefix A\_j with j\textless=14, the evaluation
\path|l_j(p)=coefficient| of p\_d annihilates all its contained bone
boundaries: its support is contained in the maximal omission of B15
treated by y\_15. But l\_j(w)=1. Therefore

\begin{verbatim}
Z -> H_Aj, n -> [n w]
\end{verbatim}

is injective with left inverse l\_j. At A15, equation (3) kills {[}w{]}
integrally; inclusion and the same preimage kill it at A16. The actual
cyclic image subdiagram is thus

\begin{verbatim}
Z --id--> ... --id--> Z --> 0 --> 0,
labels             A0,...,A14  A15  A16.                 (6)
\end{verbatim}

Each displayed 0 remains the zero AT ITS LABEL in the reconstruction.
For example the transition sends (A14,{[}w{]}) to (A15,0), whereas tau
sends it to (bottom,0). Those points have different labels. This is an
operational use of the upstream support-indexed homology, not a
numerical rank substituted for a transition map.

For the original chain inclusion C\_A0 -\textgreater{} C\_A15, let

\begin{verbatim}
Krep={p in C_A0^1 : its original inclusion lies in B_A15}.
\end{verbatim}

The map Krep/B\_A0 -\textgreater{} ker(H\_A0 -\textgreater{} H\_A15),
{[}p{]}-\textgreater{[}p{]}, is well-defined, injective, and surjective:
source boundaries map to boundaries; the prequotient kernel is exactly
B\_A0; every killed class has a source representative in Krep. The
element w belongs to Krep by (3). Evaluation at p\_d gives the split
injection Z{[}w{]} into this actual comparison kernel. No claim that it
is the entire kernel is made.

One may also keep the constant charge-line diagram E\_A=Z, at every
supported label, and the map E\_A -\textgreater{} H\_A sending n to
{[}nw{]}. On the prefix chain its supported-zero preimage is 0 through
A14 and all of Z from A15 onward. Its global-zero kernel is only its
bottom fibre: a supported source maps to a supported target label, even
when its amplitude is zero. Their inclusion is exactly the upstream D8
map, here calculated on the charge-line comparison.

\section{6. The actual Rees inclusion, integral lift, and residue
dual}\label{edition-8-section-9}

Use I\_Z=Z w inside the original terminal cell module. Define

\begin{verbatim}
V_Z={n in C_A16^0 : partial n belongs to Z w},
c:V_Z -> I_Z, n -> partial n,
F_j V_Z=V_Z intersect C_Aj^0    (0<=j<=16).
\end{verbatim}

Outside this range extend by 0 at j\textless0 and by V\_Z at
j\textgreater=16. The target filtration is G\_j I\_Z=0 at j\textless0
and I\_Z at j\textgreater=0. This is the source's comparison-image
construction on the actual incidence map, with the original integer
cutoff labels retained.

Evaluation at p\_d proves c(F\_j)=0 for j\textless15; a15 proves
c(F\_j)=I\_Z for j\textgreater=15. Its quotient-to-image isomorphism
sends {[}n{]} to partial n with inverse k w -\textgreater{} {[}k a15{]};
changing a preimage is exactly addition of an original incidence-kernel
element.

Writing the decreasing indices as a=-j gives precisely the Rees
convention of the source:

\begin{verbatim}
L_Q=T^15 Z[T] w subset L_P=Z[T] w,
D_Z=L_P/L_Q = Z[T]/(T^15) w.                            (7)
\end{verbatim}

The isomorphism sends a polynomial of degree below 15 to its original
class; the inverse takes the unique remainder modulo T\^{}15. Its
integral basis is w,Tw,\ldots,T\^{}14w. Rationalization acts
coefficientwise and has no kernel on this explicitly free abelian group.
Over the source's characteristic-zero field Q it yields a defect of
length 15. This statement concerns (7), not the torsion of every benzel
cohomology group.

The dual lattice quotient is T\textsuperscript{(-15)Z{[}T{]}w}* /
Z{[}T{]}w\^{}*. Its residue pairing is

\begin{verbatim}
([f w],[g w^*]) -> coefficient of T^(-1) in fg.
\end{verbatim}

Changing representatives changes fg by a polynomial, so the pairing is
well-defined. In the bases T\^{}i w and T\textsuperscript{(-15+j)w}*,
its matrix is 1\_\{i+j=14\}; this matrix is its own inverse over Z. Thus
all integral units, support maps and the residue coefficient are
retained.

The graded defect and boundary character are

\begin{verbatim}
Theta_D(z)=1+z+...+z^14,
A_P(z)-A_Q(z)=1-z^15=(1-z)Theta_D(z),
-[z d/dz (A_P-A_Q)]_(z=1)=15.                           (8)
\end{verbatim}

On this same rank-one comparison, both its m-th tensor and symmetric
power have inclusion T\^{}(15m)Z{[}T{]} subset Z{[}T{]}. Their
rationalized defects have length 15m, with the displayed monomial
remainder bases. This is an evaluated amplification formula, not an
unproved sublinear estimate or a claimed vanishing result.

The positive release threshold is 16 by (5), while the actual integral
image delay measured by (7) is 15. Their relationship is supplied by the
charge-one inclusion of section 2, the signed lift (2), its dual
obstruction (4), and the positive lift beta. In particular an actual
supported-zero homology class at A15 remains attached to a support with
an empty positive filling fibre.

\section{7. What was run and what changes
next}\label{edition-8-section-10}

\path|python verify.py| and \path|python -O verify.py| replay the exact
affine certificates without a solver, the 24-cell signed identity, the
source-row matches, all \path|scalar/reconstructed-module| laws on the
stated finite test sets, the original augmented cochain squares, the
noninjective homology transition, and the residue pairing. Separate
tests enumerate every contained bone at d=4,\ldots,100 and check the
original positive and signed incidence equations.
\path|recorded-checks.json| states the exact execution scope. Finite
algebra tests are not represented as a Lean proof. The preceding
all-parameter fifth construction is a pinned dependency, not a new
discovery in this turn.

The new result removes every proper subset of the current repair
dictionary from the positive search for every d, while showing
explicitly that a signed integral lift was available earlier. Progress
through other collars must therefore change the available support
dictionary or source configuration rather than search its smaller
subsets again. The augmented support complex and exact dual incidence
maps provide the data structures for doing that calculation. No uniform
variable-h positive lift has yet been supplied.

\chapter{Turn 5: variable-parameter packing}\label{edition-9}

Source: \path|sources/benzel-p7-turn5/benzel-p7-turn5/turn5/PROOF.md|.

SHA-256:
\path|60d6e9338fa71b6a7bae539d7eee249b26db44377ac8a661cb58ee0402b9d689|.

15 September 2026. Parent workbench:
\path|6fe012fd1facd564e07e1fa052ad219debdd8278| (draft PR \#25). This
continuation proves the packing below for the full displayed
two-parameter domain. It supplies complete positive tilings for three
specified one-stone regions. It does not prove the full P7 statement,
the entire centered-one-stone family, or optimality of any release set.
Independent review, Lean replay, and priority certification have not
been performed.

\section{1. Original carriers and source
comparison}\label{edition-9-section-1}

Retain the axial cell \texttt{(x,y)} and its three original differences

\begin{verbatim}
u=y-x, v=1-x-2y, w=2x+y-1; u+v+w=0.
\end{verbatim}

The original \texttt{V(a,b)} consists of cells for which all three
differences lie in \texttt{{[}1-a,b-1{]}}. The four permitted placed
tiles are

\begin{verbatim}
R(x,y)={(x,y),(x+1,y),(x,y+1)},
H(x,y)={(x,y),(x+1,y),(x+2,y)},
V(x,y)={(x,y),(x,y+1),(x,y+2)},
D(x,y)={(x,y),(x+1,y-1),(x+2,y-2)}.
\end{verbatim}

Use \path|W_h(d)=V(d+3h,2d+3h)|, integers \texttt{d\textgreater{}=3},
\path|0<=h<=d(d-1)/2|. The inverse on this parameter image is
\texttt{d=b-a}, \texttt{h=(2a-b)/3}. No parameter of the original region
is replaced.

Write \path|rho(x,y)=(y,1-x-y)| and \path|F(x,y)=(x,1-x-y)|. Direct
substitution gives

\begin{verbatim}
(u,v,w)(rho p)=(v,w,u)(p),
(u,v,w)(F p)=(-w,-v,-u)(p).
\end{verbatim}

Thus rho has inverse rho\^{}2 and preserves each original benzel. F is
its own inverse and exchanges V(a,b) with V(b,a). On tiles, rho sends H
to V, V to D, D to H, and preserves R; F exchanges H and D and preserves
V and R. These assertions follow by substituting each of the three
original offsets, not from a change of drawing convention.

The band shape below was obtained by reading the explicit Kim--Propp
sector in the pinned public P4 formalization. At \texttt{h=d(d-1)/2},
set \texttt{n=y-1}. The formulas here become exactly its
\path|kpSectorBandBase| and \path|kpSectorBandCount|. Its containment
source explicitly uses \path|V(kpB(d),kpA(d))|; the F-map above returns
that carrier to \texttt{W\_h(d)}. The endpoint is retained as a known
construction, not claimed as a discovery. The argument below establishes
the displayed extension over variable h directly, without importing an
unexecuted Lean result.

Pinned source:
\path|crabsatellite/benzel-problem4@f8f91033216ba16fa851adbb9d6a9cf005509619|,
\path|KimProppSector.lean|, \path|KimProppContainment.lean|,
\path|BoneRuns.lean|, \path|KimProppGeometry.lean| under
\path|lean4/BenzelProblem4Kernel/|.

\section{2. The two-parameter positive
packing}\label{edition-9-section-2}

Let \texttt{t=d(d-1)/2} and \texttt{Delta=t-h}. For a positive integer
y, let r(y) be the unique positive integer satisfying

\begin{verbatim}
r(y)(r(y)-1)/2 < y <= r(y)(r(y)+1)/2.
\end{verbatim}

For every integer \texttt{1\textless{}=y\textless{}=2h+d}, define a
horizontal band with first coordinate L\_y and q\_y bones:

\begin{verbatim}
y<=h:  L_y=1-2y-r(y),       q_y=y;
y>h:   L_y=y-3h-d+1,        q_y=min(h,2h+d-y).
\end{verbatim}

Its bones are \texttt{H(L\_y+3j,y)}, for
\texttt{0\textless{}=j\textless{}q\_y}. Empty bands contribute no
generators. Let E be this actual horizontal packing. Define

\begin{verbatim}
Gamma(d,h)=F(E union rho(E) union rho^2(E)).
\end{verbatim}

\textbf{Result.} For every integer \texttt{d\textgreater{}=3} and
\texttt{0\textless{}=h\textless{}=t}, Gamma is a pairwise disjoint
positive packing in the original W\_h(d). It contains exactly
\texttt{3h(h+d)} bones. The uncovered original cell set U has exactly
\texttt{3(t-h)} cells and satisfies the integer incidence identity

\begin{verbatim}
partial Gamma + 1_U = 1_(W_h(d)).                         (2.1)
\end{verbatim}

The proof follows in full.

\subsection{2.1 Exact band intervals and
containment}\label{edition-9-section-3}

E is first placed in the reflected original carrier
\path|V(2d+3h,d+3h)|. Its lower and upper difference bounds are

\begin{verbatim}
ell=1-2d-3h,  M=d+3h-1.
\end{verbatim}

In each nonempty band the cell x-coordinate ranges over the following
closed integer interval:

\begin{verbatim}
prefix, 1<=y<=h:              [1-2y-r, y-r];
middle, h<y<=h+d:             [y-3h-d+1, y-d];
tail, h+d<y<2h+d:             [y-3h-d+1, 3h+2d-2y].
\end{verbatim}

Its length is respectively \texttt{3y}, \texttt{3h}, and
\texttt{3(2h+d-y)}, so consecutive bones partition each interval.
Different bands have different y-coordinates.

For prefix bands, \texttt{1\textless{}=r\textless{}=d-1}, since
\texttt{y\textless{}=h\textless{}=t}. The exact ranges of the three
differences on the interval are

\begin{verbatim}
u in [r,3y+r-1],
v in [1-3y+r,r],
w in [1-3y-2r,3y-2r-1].
\end{verbatim}

Using \texttt{1\textless{}=y\textless{}=h} and
\texttt{1\textless{}=r\textless{}=d-1}, each lower endpoint is at least
ell and each upper endpoint is at most M. For example, the smallest w is
at least \texttt{ell+2}, and the largest u is at most \texttt{M-1}.

For the middle bands the exact ranges are

\begin{verbatim}
u in [d,3h+d-1],
v in [1-3y+d,3h+d-3y],
w in [3y-6h-2d+1,3y-2d-1].
\end{verbatim}

At \texttt{h+1\textless{}=y\textless{}=h+d}, the v-lower bound is at
least ell, the w-lower bound at least \texttt{ell+3}, and the largest u
and w are M. The other endpoints satisfy the same bounds directly. At
h=0 these bands are empty, so no unsupported cell assertion is used.

For the tail bands the ranges are

\begin{verbatim}
u in [3y-3h-2d,3h+d-1],
v in [1-3h-2d,3h+d-3y],
w in [3y-6h-2d+1,6h+4d-3y-1].
\end{verbatim}

Substitution of \path|h+d+1<=y<=2h+d-1| places these intervals in
\texttt{{[}ell,M{]}}. Thus every band cell is in the specified reflected
benzel. Rho preserves it and F returns it to the original benzel,
proving containment for every cell of Gamma.

\subsection{2.2 The three rotated sectors are
disjoint}\label{edition-9-section-4}

Every E-cell has y\textgreater=1 and x\textless y. The last inequality
follows from the upper bound \texttt{y-r} in the prefix and the bound
\texttt{\textless{}=y-d} in the middle and tail.

An E-cell \texttt{(x,y)} in rho(E) would have preimage
\texttt{(1-x-y,x)} in E. Consequently
\texttt{1\textless{}=x\textless{}y}.

When y\textless=h, both x and y lie in prefix bands. The upper bound in
row y gives \texttt{y-x\textgreater{}=r(y)}. The lower bound for the
first coordinate in row x gives \texttt{y-x\textless{}=r(x)}.
Monotonicity gives \texttt{r(x)\textless{}=r(y)}. All these inequalities
force \texttt{r(x)=r(y)=r} and \texttt{y-x=r}. But a single triangular
block contains r consecutive integers, so its two elements differ by at
most r-1. This is a contradiction.

When y\textgreater h, the upper bound in row y gives
\texttt{y-x\textgreater{}=d}. For x\textless=h, the prefix lower bound
gives \path|y-x<=r(x)<=d-1|, again a contradiction. For x\textgreater h,
the tail-or-middle lower bound in row x gives
\texttt{2x+y\textless{}=3h+d}. But \texttt{x\textgreater{}=h+1} and
\texttt{y-x\textgreater{}=d} give \path|2x+y=3x+(y-x)>=3h+3+d|, the
final contradiction.

Therefore E and rho(E) are disjoint. Rotating this equality proves
disjointness of the other two pairs. F is injective, so all bones in
Gamma are pairwise disjoint.

\subsection{2.3 Counts and the original
residual}\label{edition-9-section-5}

The first h bands have \texttt{h(h+1)/2} bones. The middle d bands have
dh bones. The final bands contribute \texttt{h(h-1)/2}. Hence E has
\texttt{h(h+d)} bones and Gamma has \texttt{3h(h+d)}.

For completeness, the original region count can be obtained through the
previously constructed three-phase cell map, with both inverse formulas
retained:

\begin{verbatim}
Phi_d(r,s,p)=(d-2-2r-s,r-s)+e_p,
e_0=(0,0), e_1=(1,0), e_2=(0,1).
\end{verbatim}

For cell \texttt{(x,y)}, choose \path|p=x+2y-(d-2) mod 3| in \{0,1,2\};
write \path|e_p=(alpha,beta)|. Then

\begin{verbatim}
r=(y-x+d-2-beta+alpha)/3,
s=(d-2-x-2y+alpha+2beta)/3.
\end{verbatim}

The residues make these integers, and substitution proves both inverse
identities. Put \texttt{R=r+h}, \texttt{S=s+h}. All phases have
\texttt{R,S\textgreater{}=0}, \texttt{R+S\textless{}=d+3h-2}. For phases
0,1,2 respectively, the additional upper bounds on (R,S) are

\begin{verbatim}
(d+2h-1,d+2h-2), (d+2h-1,d+2h-1), (d+2h-2,d+2h-1),
\end{verbatim}

and the lower bounds on R+S are \texttt{h-1,h,h-1}. These are obtained
by substituting Phi in the original three pairs of inequalities. In each
phase the big integer triangle has \path|binom(d+3h,2)| points. The
three removed corner triangles have counts
\path|binom(h,2),binom(h,2),binom(h+1,2)| in a phase-dependent order.
Their inequalities show they are disjoint: both upper caps
simultaneously would exceed \texttt{R+S\textless{}=d+3h-2}, and a
low-sum cap cannot meet an upper cap for d\textgreater=3. Empty caps at
h=0 have count zero. Thus each phase has

\begin{verbatim}
binom(d,2)+(3d-1)h+3h^2
\end{verbatim}

cells. It follows that

\begin{verbatim}
|W_h(d)|=3(t-h)+9h(h+d).
\end{verbatim}

Subtracting the number of disjoint packed cells proves
\texttt{\textbar{}U\textbar{}=3(t-h)}. Since U is the literal set
complement, (2.1) holds at every original cell with coefficient exactly
one.

At h=t, U is empty and the construction recovers a complete known
bone-only endpoint. At h=0, Gamma is empty; the phase map on
\path|r,s>=0, r+s<=d-2| groups its three phases into the original base
right stones. The interior construction is a partial packing, not a
completed tiling.

\section{3. What is retained by the
cohomology}\label{edition-9-section-6}

Let N be the number of bones in Gamma. Let C\^{}0\_Gamma be the free
integer module on these actual placed bones, and C\^{}1 the free integer
module on all original region cells. Use their actual incidence
differential and zero following differential.

Order the cells of each bone b lexicographically as
\path|(c_0,c_1,c_2)|. There is an explicit isomorphism

\begin{verbatim}
C^1 / im(partial_Gamma)  ->  Z[U] direct-sum (direct-sum over b in Gamma of Z^2),
\end{verbatim}

sending a cell coefficient vector a to

\begin{verbatim}
((a_c)_(c in U), (a_(c_1)-a_(c_0),a_(c_2)-a_(c_0))_b).       (3.1)
\end{verbatim}

The inverse puts the hole coefficients on U and sends \texttt{(u,v)} for
b to \path|u[c_1]+v[c_2]|. Composing this inverse with (3.1) subtracts
exactly \path|sum_b a_(c_0) partial[b]| from a. The other composition is
identity. These equations prove the isomorphism and identify its
complete kernel with the original incoming boundaries.

In particular, the original region vector maps to \path|(1_U,0,...,0)|.
The 2N other quotient directions remain in the construction; they are
not silently deleted because this distinguished vector has zero
coordinates there.

The inclusion of these chosen bone generators into *all* original
contained bones induces the representative-preserving map from this
quotient to the full bone cohomology. Its kernel is exactly

\begin{verbatim}
im(partial_all) / im(partial_Gamma),
\end{verbatim}

with maps taking the same original cell class and with the inverse
obtained by choosing that representative. Inclusion, image, and quotient
are taken inside the same C\^{}1. Thus the partial packing gives a
specific source representative for the remaining calculation and not an
inference from a vanishing global observation.

\subsection{3.1 Positive release fibres and the actual Split-Zero
reconstruction}\label{edition-9-section-7}

Index releases of Gamma by subsets A. Let

\begin{verbatim}
Lambda_A=U union (union of cells of b for b in A),  v_A=1_(Lambda_A).
\end{verbatim}

The nonnegative integral fillings of Lambda\_A are in bijection with
original region tilings containing every frozen bone in
Gamma\textbackslash A. The forward map adjoins those bones; the inverse
removes them. Disjointness proved in section 2 verifies both maps and
their inverse identities. At A=Gamma this is the original entire tiling
fibre. A smaller support is not asserted to have all its tilings.

Adjoin an external bottom label to the powerset of Gamma. Map bottom to
the empty support and A to Lambda\_A. This map preserves joins, even
when the residual is empty. A supported empty release label and the
external bottom are retained as labels; the support map is permitted to
be noninjective at a zero-residual endpoint, exactly as in upstream
D1--D8.

For the support-indexed complex use C\_A\^{}0=Z{[}all permitted
contained tiles{]}, \path|C_A^1=Z[Lambda_A],| C\_A\^{}2=0. The original
incoming map is incidence. Its reconstruction is the disjoint union of
fibres with addition transported to the union label. The scalar pair
\path|(amplitude,presence)| realizes

\begin{verbatim}
G(Z) = {(0,0)} union (Z times {1}).
\end{verbatim}

The supported zero acts as \path|(A,x)->(A,0)|; external tau acts as
\path|(A,x)->(bottom,0)|. The reconstructed differential square has the
first value, with the original source and target degree. Its homology
quotient coequalizes the original incidence and \path|(A,n)->(A,0)|. For
any split-linear map equalizing them, replacing a representative by a
boundary changes its image by its absorbing fibre zero. This proves
descent, and surjectivity proves uniqueness: it is the actual internal
coequalizer.

To retain the affine addition of released source tiles, use the
augmented source

\begin{verbatim}
C_hat_A^0=Z direct-sum C_A^0,
d_hat_A(r,n)=partial n-r v_A.
\end{verbatim}

For A subset B, put \path|eta_AB=sum_(b in B\A)[b]|. The original cell
partitions give

\begin{verbatim}
v_B=j_AB v_A+partial eta_AB,
eta_AC=j_BC eta_AB+eta_BC.
\end{verbatim}

Hence the literal linear transition is

\begin{verbatim}
(r,n) -> (r,j_AB n+r eta_AB).                             (3.2)
\end{verbatim}

Expanding the differential proves \path|d_hat_B j_hat_AB=j_AB d_hat_A|;
expanding the second identity proves composition. Addition in the
reconstructed G(Z)-module transports *both* charges and chains by (3.2).
The bottom fibre is zero. Charge-one nonnegative integral cycles are
exactly the positive fillings, since each original cell coefficient is
one. The inclusion of that fibre into the integral cycle fibre retains
every coefficient; negative scalar multiplication is not asserted to
preserve the positive subset.

These are the original Split-Zero reconstruction, supported
differential, coequalizer, and comparison-kernel constructions
specialized to a variable-h source configuration. Upstream pin:
\path|KokunoYumeto/zeta-function-research-reader@1c8ec52c85c173adc9f8403a8a26914955e2f5a9|,
\path|workbenches/splitzero-tandem/tex/support_diagrams.tex|, D1--D8;
\path|formal/splitzero/SplitZeroComplex.lean|.

\section{4. Exact residual on the entire one-stone
ray}\label{edition-9-section-8}

Set \texttt{h=t-1}, for every integer d\textgreater=3. Define

\begin{verbatim}
p_d=(t,t-d+1).
\end{verbatim}

Then the residual of section 2 is exactly

\begin{verbatim}
U_d={p_d,rho(p_d),rho^2(p_d)}.                            (4.1)
\end{verbatim}

Proof: in the pre-reflection carrier, one of the three candidate cells
is \texttt{z=(t-d+1,t)}. In row y=t, the middle-band upper endpoint is
\texttt{t-d}, so z is not in E. The row coordinate of rho(z) is
\texttt{d-2t\textless{}1}, so rho(z) is not in E. For rho\^{}2(z) the
row is \path|t-d+1=binom(d-1,2)|, whose triangular block is d-2; its
lower endpoint is \texttt{1+d-2t}, while the candidate first coordinate
is \texttt{d-2t}. Thus rho\^{}2(z) is not in E either. Direct
substitution places z in the reflected original benzel: its differences
are \path|d-1,d-3t,3t-2d+1|, all in \path|[4-2d-3t,d+3t-4]| for
d\textgreater=3. Rotation and reflection give three uncovered original
cells. Section 2 proves there are exactly three, establishing (4.1).

Their y-coordinate range is \texttt{3t-d\textgreater{}=6}. Every allowed
original trimer has y-range at most two. Thus the frozen residual has no
tile and its charge-one positive fibre is empty. This is an explicit
proof for this source configuration, not a nonexistence assertion about
W\_(t-1)(d). Indeed, the next section constructs positive completions
after source changes.

The centered right stone S\_0=R(0,0) is invariant under rho and lies in
every region on this ray. To study tilings containing S\_0, release at
least the parent bones meeting S\_0. For such A, the original
bone-filling support is

\begin{verbatim}
Lambda_A=(U_d union cells(A)) \ S_0.
\end{verbatim}

The forward positive map adjoins S\_0 and the frozen parent bones; its
inverse removes those same tiles. The transition (3.2) remains valid
because S\_0 is subtracted in both v\_A and v\_B. At the full release
set this parametrizes exactly the original tilings containing this
specified stone. Other stone positions are not excluded by the
construction or the problem statement.

\section{5. Completed dual-guided repairs}\label{edition-9-section-9}

The package contains complete positive tilings, with the centered stone,
at

\begin{verbatim}
d=3,h=2: V(9,12), 30 bones and one right stone;
d=4,h=5: V(19,23), 135 bones and one right stone;
d=5,h=9: V(32,37), 378 bones and one right stone.
\end{verbatim}

The first two give source controls in already treated region families.
The third goes beyond the fixed h\textless=5 family. These statements
have only the three listed parameter values; no all-d positive section
is asserted.

The parent bone order is \path|construction.packing(d,h)|, sorted by
original kind and original cell tuple. \path|evidence/repair-d.json|
identifies released indices, every replacement tile, and each original
cell. It is a proof certificate through direct incidence expansion, not
an optimization status.

For d=5, the initial support releases 30 parent bones. The first exact
cell cochain y has value -3 on its 90-cell target vector and nonnegative
value on every contained bone boundary. Eleven more specified source
bones enable original tiles with negative extended-y values. On that
41-release, 123-cell support, a second exact cochain has value -2 and
again nonnegative value on every contained bone. Ten additional releases
enable negative directions for that cochain. The resulting 51-release,
153-cell support has the recorded positive 51-bone filling. Adjoining
the 327 frozen sector bones and S\_0 gives the original full tiling of
V(32,37).

For each rejected support, evaluation of a supposed nonnegative real
filling n against its exact cochain gives

\begin{verbatim}
y(v_A)=<partial^* y,n> >= 0,
\end{verbatim}

contradicting the recorded negative value. The matrix partial is the
original full contained-bone incidence matrix. The verifier enumerates
all its columns. To choose a new support, extend y by zero to the actual
larger support and enable a bone b with \path|<y,partial b><0|; the data
identify every old parent bone released to contain b. This proves
precisely which obstruction is invalidated. It does not assume that
defeating one dual certifies feasibility. The final independent
cell-partition certificate certifies feasibility on the final support.

All label inclusions, eta\_AB vectors, and charged cochain squares are
checked against (3.2), and the final positive fibre inverse is replayed.
The result is an executed application of the support-indexed complex and
its exact dual incidence pairing.

\section{6. Bounds, failed attempts, and next
source}\label{edition-9-section-10}

A greedy inward hole-slide rule did not complete most tested residuals.
Those finite search outcomes are retained as exploration, not a theorem
about all sequences of moves. A slide itself has an exact inverse:
moving a hole p to p+3v replaces the bone on p+v,p+2v,p+3v by the bone
on p,p+v,p+2v. The two incidence differences equal {[}p{]}-{[}p+3v{]}.
The positivity assertion requires the old bone actually present, as
checked by the exploratory implementation.

For d=6 and d=7, support-expansion searches reached further supports but
did not produce accepted complete positive certificates in this session.
Optimizer infeasibility reports and time limits are not promoted to
mathematical exclusions. A preliminary rational dual extraction failed
exact checking after numerical reconstruction; it was rejected. The
accepted d=5 certificates use exact fractions and pass all incidence
inequalities.

The new all-parameter result is the two-variable packing, its literal
residual, and the complete \path|quotient/support| maps. The next
constructive calculation is a repeatable transport of the three
explicitly located holes on the one-stone ray, retaining which sector
bones each move releases. The fixed-rung collar tables are no longer the
only available source. No full numbered-problem resolution, conditional
resolution, or unproved positivity claim is part of this continuation.

\chapter{Turn 6: stable residual and invariant
families}\label{edition-10}

Source: \path|sources/benzel-p7-turn6/benzel-p7-turn6/turn6/PROOF.md|.

SHA-256:
\path|271a3bc14d627c37f5d8e7933f6d6c13a757bcc22b6b3e0913d8078232228a8c|.

15 September 2026. Written constructive research proofs, with exact
replay and finite certificates. Independent mathematical review and a
Lean build have not been performed. Priority of these constructions has
not been established. The full original Propp Problem 7 is not claimed
resolved.

This continuation retains the original axial cells, the four original
permitted prototiles, integer incidence coefficients, every support, and
the positive filling fibre. An exact all-parameter argument is
distinguished from its finite replay by explicitly giving the argument;
no experimental cutoff replaces a quantifier. The finite core
certificates in §7 produce infinite families through the proved
identity-on-generators transport, not by extrapolation.

\section{1. Original objects and notation}\label{edition-10-section-1}

A cell is (x,y) in Z². Its original differences are

\begin{verbatim}
u=y-x,   v=1-x-2y,   w=2x+y-1.
\end{verbatim}

They sum to zero. The benzel V(a,b) is the set where all three
differences belong to {[}1-a,b-1{]}. Throughout the principal lane,

\begin{verbatim}
t_d=d(d-1)/2,  d>=3,  0<=h<=t_d,
W_h(d)=V(d+3h,2d+3h),  Δ=t_d-h.
\end{verbatim}

The inverse parameter map is d=b-a, h=(2a-b)/3. Its domain is retained:
these formulas concern this original integral lane, not every real pair
(a,b). The original right stone is

\begin{verbatim}
R(x,y)={(x,y),(x+1,y),(x,y+1)}.
\end{verbatim}

H(x,y), V(x,y), D(x,y) are the three bones with cell increments (1,0),
(0,1), (1,-1), respectively, each of length three. A placed tile
generator maps under ∂ to the sum of its three original cell generators.
A nonnegative integral preimage of the all-ones region vector is a
tiling: every cell sum is one, so no tile multiplicity can exceed one
and no overlap is possible. A tiling gives the reverse preimage by its
actual placed tiles.

Use the original bijections

\begin{verbatim}
ρ(x,y)=(y,1-x-y),       F(x,y)=(x,1-x-y).
\end{verbatim}

Their inverses are ρ² and F. On differences they give (v,w,u) and
(-w,-v,-u). Thus ρ preserves each original benzel and F exchanges V(a,b)
with V(b,a). On tile orientations ρ cycles H,V,D, and F interchanges H,D
and fixes V. Both preserve the original right-stone shape. All
translations are literal translations of the original integer cells.

For n\textgreater=1 let r(n) be the unique integer r\textgreater=1
satisfying

\begin{verbatim}
t_r < n <= t_(r+1).
\end{verbatim}

The interval for a fixed r has exactly r integers. The integer inverse
is implemented with isqrt and its defining inequalities are checked.

The original area identity is

\begin{verbatim}
|W_h(d)|=3t_d+(9d-3)h+9h²=3Δ+9h(h+d).              (1)
\end{verbatim}

For an explicit counting derivation, use

\begin{verbatim}
Φ_d(r,s,p)=(d-2-2r-s,r-s)+e_p,
e_0=(0,0), e_1=(1,0), e_2=(0,1).
\end{verbatim}

For an original cell recover p as the representative 0,1,2 of x+2y-(d-2)
modulo 3. With e\_p=(α,β), its inverse is

\begin{verbatim}
r=(y-x+d-2-β+α)/3,
s=(d-2-x-2y+α+2β)/3.
\end{verbatim}

The phase choice proves integrality; substitution proves both inverse
laws. Set R=r+h,S=s+h and n=d+2h-1. All phases satisfy R,S\textgreater=0
and R+S\textless=d+3h-2. The remaining bounds are

\begin{verbatim}
phase 0: R<=n,   S<=n-1, R+S>=h-1;
phase 1: R<=n,   S<=n,   R+S>=h;
phase 2: R<=n-1, S<=n,   R+S>=h-1.
\end{verbatim}

These follow by substituting Φ in the six original inequalities. In each
phase the big triangle loses three disjoint corner triangles. Its
remaining cardinality is C(d+3h,2)-2C(h,2)-C(h+1,2), equal to
t\_d+(3d-1)h+3h². Their sum proves (1). At h=0 all three domains
coincide; the three cells at each (r,s) are the original right stone.
Hence

\begin{verbatim}
T_0(d)={R(d-2-2r-s,r-s): r,s>=0,r+s<=d-2}          (2)
\end{verbatim}

is a complete base tiling, with t\_d stones.

\section{2. A shifted positive packing on the entire two-parameter
domain}\label{edition-10-section-2}

Fix an integer shift 0\textless=s\textless=Δ. In the reflected benzel
define a horizontal sector E\_s(d,h) by its bands. For
1\textless=y\textless=2h+d set

\begin{verbatim}
L_y=1-2y-r(y+s), q_y=y                      (y<=h),
L_y=y-3h-d+1,    q_y=min(h,2h+d-y)           (y>h).
\end{verbatim}

The band consists of H(L\_y+3j,y), 0\textless=j\textless q\_y. Let
Γ\_s(d,h) be the images of these actual placed bones under F, Fρ and
Fρ². This specifies every tile, not only a region whose tileability is
asserted.

\subsection{2.1 Containment}\label{edition-10-section-3}

The reflected bounds are A=1-2d-3h and B=d+3h-1. For a prefix cell
x=1-2y-r+z, 0\textless=z\textless=3y-1, its differences are

\begin{verbatim}
u=3y+r-1-z,  v=r-z,  w=1-3y-2r+2z.
\end{verbatim}

Because y+s\textless=h+Δ=t\_d, one has 1\textless=r\textless=d-1. Thus u
ranges from r to 3y+r-1\textless=B; v ranges from r-3y+1\textgreater=A
to r\textless=B; and w ranges from 1-3y-2r\textgreater=A to
3y-2r-1\textless=B. These are the original six bounds.

For a tail cell write y=h+ell, 1\textless=ell\textless=h+d,
\path|q=min(h,h+d-ell),| x=ell-2h-d+1+z, 0\textless=z\textless=3q-1.
Then

\begin{verbatim}
u=3h+d-1-z,
v=d-3ell-z,
w=3ell-3h-2d+1+2z.
\end{verbatim}

The inequalities q\textless=h and ell+q\textless=h+d give
A\textless=u,v,w\textless=B directly. In particular
u\textgreater=d.~Empty bands have no cell to check.

\subsection{2.2 Sector disjointness}\label{edition-10-section-4}

Each sector cell has y\textgreater=1 and x\textless y:
u\textgreater=r\textgreater=1 in a prefix and u\textgreater=d in a tail.
A cell in E\_s intersect ρE\_s would therefore satisfy
1\textless=x\textless y and both (x,y) and (1-x-y,x) in E\_s.

For prefix rows x,y, the lower endpoint constraints give

\begin{verbatim}
y-x>=r(y+s),        y-x<=r(x+s).
\end{verbatim}

Monotonicity forces equality r(x+s)=r(y+s)=r and y-x=r. But x+s and y+s
lie in the same r-element triangular block, whose maximum difference is
r-1. This is a contradiction.

For a tail row y and a prefix row x, those inequalities instead give
y-x\textgreater=d and \path|y-x<=r(x+s)<=d-1.| For two tail rows, the
lower endpoint of row x gives 2x+y\textless=3h+d.~The original row y
gives y\textgreater=x+d; with x\textgreater=h+1 this implies
2x+y\textgreater=3h+d+3. These cases exhaust 1\textless=x\textless y.
The other pairs of rotated sectors reduce to this pair by ρ. Reflection
preserves disjointness. Within one sector the bands have different y,
and within a band the original length-three intervals are disjoint.

\subsection{2.3 Count and original incidence
equation}\label{edition-10-section-5}

One sector has

\begin{verbatim}
sum_(y=1)^h y + sum_(ell=1)^(h+d) min(h,h+d-ell)
   = h(h+1)/2 + hd+h(h-1)/2 = h(h+d)
\end{verbatim}

bones. Thus Γ\_s has 3h(h+d) disjoint bones. By (1) its complement has
3Δ cells. The original integer incidence equation is

\begin{verbatim}
∂Γ_s + 1_(W_h\cells Γ_s)=1_(W_h).                 (3)
\end{verbatim}

The known Kim--Propp zero-invariant construction is the endpoint of this
band scheme. That endpoint retains its source attribution. The argument
above proves the larger packing domain directly; it does not assume
positive completion of the complement.

\section{3. The full residual has a parameter-independent inverse
description}\label{edition-10-section-6}

Take s=Δ. Define the locally finite infinite sector E\_∞(Δ) by the bands

\begin{verbatim}
H(1-2y-r(y+Δ)+3j,y), y>=1, 0<=j<y.
\end{verbatim}

Let Σ\_Δ=F(E\_∞ union ρE\_∞ union ρ²E\_∞). Its finite prefix through row
h is literally the prefix of Γ\_Δ(d,h); the tile generators have
identical coordinates, not merely equivalent counts.

\subsection{3.1 Exact rank map on the original cell
lattice}\label{edition-10-section-7}

The three differences of an original cell are pairwise unequal: their
cyclic differences are all 2 modulo 3. Since they sum to zero, their
unique minimum is -k for an integer k\textgreater=1. Rotate to put that
minimum in coordinate u. The cell then has the unique form

\begin{verbatim}
p_(k,m)=(k-m,-m),     0<=m<k.
\end{verbatim}

Indeed v=1-k+3m and w=2k-3m-1 exceed -k exactly in that range. Thus

\begin{verbatim}
(k,m,j) -> ρ^j p_(k,m),  j=0,1,2                 (4)
\end{verbatim}

is a bijection. Its inverse uses the position of the unique minimum,
k=-min(u,v,w), and m=-y after that exact rotation. Put

\begin{verbatim}
n(k,m)=t_k+m+1,
Ω_Δ={ρ^j p_(k,m): n(k,m)<=Δ, j=0,1,2}.          (5)
\end{verbatim}

Equation (4) proves \textbar Ω\_Δ\textbar=3Δ. It also proves

\begin{verbatim}
Ω_(t_k)=W_0(k),                                  (6)
\end{verbatim}

including the empty k=1 case: the lower difference bounds retain exactly
the shells of minimum -1,\ldots,-(k-1); the upper bounds follow from sum
zero. This supplies the original carrier bijection, not a picture-based
identification.

\subsection{3.2 Which original cells the infinite bands
cover}\label{edition-10-section-8}

For the canonical p=p\_(k,m), set B=k-m\textgreater=1 and
a=1-2k+3m\textless=k-2. The three possible pre-reflection coordinates
are

\begin{verbatim}
Fp=(B,B+a),   ρFp=(B+a,-m),   ρ²Fp=(-m,B).
\end{verbatim}

The middle one has nonpositive second coordinate and is outside E\_∞.
Membership of (-m,B) is exactly

\begin{verbatim}
a<=r(B+Δ)<=k.                                    (7)
\end{verbatim}

Membership of (B,B+a) is exactly a\textgreater=1 and

\begin{verbatim}
r(B+a+Δ)<=a,
\end{verbatim}

or equivalently B+Δ\textless=t\_a. Its lower endpoint inequality is
automatic in that range. For a\textgreater=1, (7)'s lower inequality is
B+Δ\textgreater t\_a; its upper inequality is B+Δ\textless=t\_(k+1). For
a\textless=0 the lower inequality is automatic and the other possible
point is absent. Since a\textless=k-2, these alternatives are disjoint
and exhaust exactly

\begin{verbatim}
B+Δ<=t_(k+1),
\end{verbatim}

which is Δ\textless=t\_k+m=n(k,m)-1. Therefore an original cell has
exactly one Σ\_Δ owner when its rank n exceeds Δ, and has no owner
exactly when it belongs to Ω\_Δ. This proves both disjointness of the
infinite packing and its entire finite complement.

The code's stable\_owner follows this inverse calculation and returns
the sector index, the original band row, the bone index, and the
original placed tile. All four are checked against its three cell
generators.

\subsection{3.3 The finite residual is exactly the same
Ω\_Δ}\label{edition-10-section-9}

The prefix of Γ\_Δ is a subset of Σ\_Δ and avoids Ω\_Δ. Also Ω\_Δ is
contained in W\_0(d), since Δ\textless=t\_d.~Reflection puts it in
V(2d,d), where every difference is \textless=d-1. A finite tail sector
has u\textgreater=d; its rotations and reflection therefore avoid Ω\_Δ
as well. Hence Ω\_Δ is disjoint from the entire finite packing and is
inside W\_h. The complement has 3Δ cells by (3), exactly the cardinality
of Ω\_Δ. Consequently

\begin{verbatim}
∂Γ_Δ(d,h)+1_(Ω_Δ)=1_(W_h(d))                    (8)
\end{verbatim}

for all original integer d\textgreater=3 and
0\textless=h\textless=t\_d.~No positive completion has been stipulated
to obtain this identity.

\section{4. Complete families at every triangular
invariant}\label{edition-10-section-10}

For d\textgreater=3 and 1\textless=k\textless=d put h=t\_d-t\_k.
Equation (8) leaves precisely Ω\_(t\_k)=W\_0(k), and (2) tiles it by
original right stones. Thus

\begin{verbatim}
Γ_(t_k)(d,t_d-t_k) union T_0(k)                   (9)
\end{verbatim}

is a complete positive type-103 tiling of V(d+3h,2d+3h). It has t\_k
right stones and 3h(h+d) bones. The placement sets are disjoint by §3,
and both parts' original cell equations add to 1\_W. Reflection supplies
the reflected pairs with an explicit inverse.

In particular k=2 proves the entire one-stone ray

\begin{verbatim}
h=t_d-1,
(a,b)=(d+3(t_d-1),2d+3(t_d-1)),  d>=3,          (10)
\end{verbatim}

with its right stone exactly R(0,0). At k=1 this retains the known
zero-invariant bone tiling; at k=d it retains the original base tiling.
Neither endpoint is counted as a new discovery here.

\section{5. The one-stone family is obtained by explicit positive chain
homotopies}\label{edition-10-section-11}

There is a direct relationship between (10) and the original unshifted
packing from turn 5. It explains the actual transport of the three
previously separated residual cells.

Set \path|P_r=(t_(r+2),t_(r+1))| and Q\_r=(1-r²,t\_(r+1)). For
r=d-2,d-3,\ldots,1, use the old horizontal band

\begin{verbatim}
H(1-r²+3j,t_(r+1)), 0<=j<t_(r+1),
\end{verbatim}

and replace it by

\begin{verbatim}
H(2-r²+3j,t_(r+1)), 0<=j<t_(r+1).               (11)
\end{verbatim}

The old union is the horizontal interval from Q\_r to P\_r-(1,0), and
the new union is from Q\_r+(1,0) to P\_r. Thus the signed tile vector
κ\_r=new-old has the exact original boundary

\begin{verbatim}
∂κ_r=[P_r]-[Q_r].                               (12)
\end{verbatim}

Also P\_r=p\_(r+2), Q\_r=ρ²p\_(r+1), where \path|p_j=(t_j,t_j-j+1).|
Include all three rotations of (11).

These old bands are actual bones in Γ\_0(d,t\_d-1). In its defining
sector they are the rows y=t\_(r+1), whose triangular block is r; Fρ²
maps the row to the displayed original horizontal interval. Different r
or sector uses disjoint old bones. Each simultaneous triple of moves
covers the three current holes and uncovers the three starts. The
process consequently changes the hole orbit of p\_(r+2) to that of
p\_(r+1). No original band awaiting use is changed by an earlier move,
because each new cell is a current hole and all such holes lie outside
those untouched original bones. At the end the hole set is

\begin{verbatim}
orbit(p_2)={(1,0),(0,0),(0,1)}=R(0,0).
\end{verbatim}

This is a proof of positivity at every step, not just a signed endpoint
equation. Each band move can also be performed bone by bone, from the
rightmost bone to the leftmost, shifting an original bone into the
adjacent hole and transferring the hole three cells to the left. Its
inverse shifts those same bones in reverse order.

The total number of moved original bones is

\begin{verbatim}
3 sum_(r=1)^(d-2) t_(r+1)=3 C(d,3).             (13)
\end{verbatim}

This is the cost of this specified construction, not a global minimum.
Since r(y+1) differs from r(y) precisely at these triangular rows, the
resulting bone set is exactly Γ\_1(d,t\_d-1), the k=2 construction of
§4, on the original placed generators.

\subsection{5.1 Actual cochain maps and their
homotopy}\label{edition-10-section-12}

Let M=Z\^{}3, concentrated in degree one, with basis e\_j for j=0,1,2.
Let f\_start(e\_j)={[}ρ\^{}j p\_d{]} and f\_end(e\_j)={[}ρ\^{}(j+2(d-2))
p\_2{]}. For decreasing r choose the rotation j+2(d-2-r) of κ\_r in the
jth path. Let H(e\_j) be its sum. Equations (12) telescope to

\begin{verbatim}
∂H=f_start-f_end.                              (14)
\end{verbatim}

Thus H is a displayed degree-minus-one cochain homotopy for the two maps
into the original bone/cell complex. Their equal induced classes are
represented by the original cells; summing (14) gives the positive
region-minus-stone equation after (11).

For a subset A of the three marked paths, let the target cell fibre be
the free module on all original cells used by those paths and their
endpoints. Let the degree-zero fibre be the free module on the original
old and new bones of those paths. Both sets are explicit finite unions.
All transitions for A⊂B are inclusions of these original generators;
H,f\_start,f\_end and ∂ commute with them. The source fibre is Z{[}A{]}.
The empty-path label has the zero modules. Consequently (14) is an
identity on the entire join-indexed diagram, not only at its top fibre.

The pinned Split-Zero reconstruction makes the total carrier a disjoint
union of those fibres. Its supported scalar e sends (A,x) to (A,0), and
τ sends it to (empty,0). Addition uses the union of path labels and the
displayed inclusions. D1--D5 apply to these very modules and maps. The
internal quotient coequalizes the original incoming boundary and its
supported-zero companion. For the comparison of two supports, D7
identifies killed original representatives modulo original boundaries
with the kernel of the induced homology map. Here (14) supplies the
particular representatives and boundaries, and (11) supplies their
positive realization. Neither is assumed from the other.

\section{6. A second complete family: triangular invariant minus
one}\label{edition-10-section-13}

For m\textgreater=3 and d\textgreater=m put Δ=t\_m-1 and h=t\_d-t\_m+1.
The following original core replacement is valid throughout this
unbounded parameter domain. At d=m use the previous complete
first-collar tiling T\_1(m) of W\_1(m). Its entire earlier proof is
retained in \path|sources/turn1-PROOF.md,| §3. For direct
reproducibility the same actual constructor is printed here. Translate
T\_0(m) by (-2,1), remove R(-2,3-m), and add

\begin{verbatim}
D(-2,3-m), D(-2,4-m), D(-1,1-m), V(1,-m),
H(-m-1+2r,m+1-r)      0<=r<m,
H(r,3-m+r)           0<=r<m-1,
V(r+2,r-m)           0<=r<m.
\end{verbatim}

The retained source proof proves original containment, all pairwise
orientation intersections, the exact three-cell inner intersection, and
complete coverage for every m\textgreater=3. This dependency is a proved
earlier construction, not an unproved tileability premise.

For d\textgreater=m+1, one has h\textgreater=m+1. Let K=V(2m+3,m+3), the
reflected core. Any cell of K with y\textgreater=1 has y\textless=m+1:
3y=u-v+1\textless=3m+5. All outer sector rows which can meet K are
therefore prefix rows.

At y=1, r(y+Δ)=m-1 and the whole band has x=-m,\ldots,2-m: one original
bone inside K. For 2\textless=y\textless=m+1, r(y+Δ)=m. The sector
interval is

\begin{verbatim}
1-2y-m <= x <= y-m.
\end{verbatim}

Its intersection with K is exactly the last three cells

\begin{verbatim}
y-m-2 <= x <= y-m.                              (15)
\end{verbatim}

The lower endpoint follows from u\textless=m+2. At
\path|x=y-m-2,y-m-1,y-m,| substitution gives all six core bounds (their
extrema occur at 2\textless=y\textless=m+1), so all three cells are
present. They form one actual bone of the sector, with index j=y-1
because \path|1-2y-m+3(y-1)=y-m-2.|

No bone crosses the core boundary: each of those rows has precisely one
last bone wholly inside, with all preceding bones outside. The three
sectors and F preserve K's original region. The finite tail rows
y\textgreater h\textgreater=m+1 miss it. Exactly 3(m+1) outer-source
bones lie wholly inside W\_1(m). Their disjoint cells, together with
Ω\_Δ, exhaust it:

\begin{verbatim}
|W_1(m)|=3(t_m-1)+9(m+1).
\end{verbatim}

Hence the actual original bones of Γ\_Δ(d,h) outside W\_1(m) partition
its complement in W\_h(d). Replacing the core by the explicit T\_1(m)
produces the complete original tiling

\begin{verbatim}
{b in Γ_Δ(d,h): cells(b) intersect W_1(m)=empty}
                   union T_1(m).               (16)
\end{verbatim}

It contains Δ right stones and 3h(h+d) bones. This proves every
triangular-minus-one invariant ray (Δ=2,5,9,14,20,\ldots) in full,
including all of their lowest admissible d.

\section{7. Fixed original patch certificates give further infinite
families}\label{edition-10-section-14}

This section contains completed finite positive certificates and the
proved original-generator transport of those certificates. It does not
postulate a positive patch for arbitrary Δ.

For each actual recorded certificate choose its finite set A of Σ\_Δ
bones. Every such bone has a unique stable owner label (sector,y,j). Let
Y=max y over A, and let D be the smallest integer d\textgreater=3 with
t\_d-Δ\textgreater=Y. At every d\textgreater=D, all those placed bones
are literally in the prefix of Γ\_Δ(d,t\_d-Δ). The support

\begin{verbatim}
Λ=Ω_Δ union cells(A)
\end{verbatim}

is therefore the same finite set of original integer cells for every
such d.~The sets of *all* original placed R,H,V,D tiles contained in Λ
are identical. Consequently the two original incidence complexes at any
two parameters d,d'\textgreater=D are isomorphic by the identity on
every cell and every tile generator; the inverse is the same identity.
This preserves integer coefficients, nonnegative monoids, kernel,
cokernel, and each positive filling fibre exactly. The ambient inclusion
maps are the displayed original cell inclusions, and commute with
incidence.

The release-subset versions use the join lattice P(A) with a separate
bottom for absence. Its empty release set still has support Ω\_Δ. For
S⊂T⊂A, let η\_ST be the sum of the newly released original bones. The
exact region vectors satisfy v\_T=jv\_S+∂η\_ST. The augmented complex
and transitions are

\begin{verbatim}
d_hat(r,n)=∂n-rv_S,
j_hat_ST(r,n)=(r,jn+rη_ST).                     (17)
\end{verbatim}

Expansion proves d\_hat j\_hat=j d\_hat and composition follows from
η\_SU=jη\_ST+η\_TU. These are genuine linear cochain maps. Charge-one
nonnegative integral cycles are the original positive patch tilings. The
pinned G(Z)-reconstruction retains (S,0) under supported zero and sends
τ to the external bottom. Internal homology uses the actual incidence
and supported-zero pair, with its original representative kernel map.
The identity-on-generators stabilization above preserves this whole
supported diagram, not merely the charge-one endpoint.

Each record in \path|evidence/stable-patches.json| prints Δ, Y, D, every
released original bone anchor and every replacement tile anchor.
verify.py independently expands these anchors, checks the stable owners,
checks distinct release bones, and checks

\begin{verbatim}
∂(replacement)=1_(Ω_Δ)+∂A.                     (18)
\end{verbatim}

No optimizer is used by that acceptance step. Equations (8) and (18)
then give the explicit complete constructor

\begin{verbatim}
(Γ_Δ(d,t_d-Δ) minus A) union replacement        (19)
\end{verbatim}

for every integer d\textgreater=D in that record. All coefficients are
+1, the support intersection is exactly the released source cells, and
the positive inverse removes that same replacement and restores A on the
specified patch-containing fibres.

The remaining d\textless D are a finite, explicitly listed set for each
record. For every claimed completed invariant the file
\path|evidence/base-cases.json| contains each such original whole-region
tiling. The verifier enumerates \textbf{all} d=3,\ldots,D-1 with
t\_d\textgreater=Δ and rejects a missing base case. It checks every cell
multiplicity and every original tile shape. The table in COVERAGE.md is
generated from these accepted certificates; only a row with no missing d
is marked COMPLETE\_RAY. Larger-parameter certificates for other records
retain their exact threshold and missing lower cases.

Together with the completed triangular and triangular-minus-one
families, these literal records prove the completed invariant ranges
stated in COVERAGE.md. Their all-d quantifier follows from the exact
transport above and the exhaustive finite base-case list, not from an
experimental sequence matching initial values.

\section{8. An exact obstruction to freezing the central stone at
invariant two}\label{edition-10-section-15}

The stable residual construction must allow the positive source fibre to
change; it cannot retain a central stone in every parameter case. For
W\_1(3)=V(6,9), prescribe R(0,0) and remove its three cells. On the 39
remaining cells define y=2 at

\begin{verbatim}
(-3,4),(-2,2),(-1,-1),(-1,0),(-1,3),(0,-3),
(0,2),(1,-2),(2,-1),(2,1),(3,-1),(4,0),
\end{verbatim}

and y=-1 at all remaining cells. Its region sum is 12*2-27=-3. The exact
verifier enumerates every contained original R,H,V,D tile and checks its
y-sum is nonnegative. Thus no nonnegative real tile combination fills
that punctured region: pairing ∂n=1 with y would give
-3=\textless∂*y,n\textgreater\textgreater=0. This concerns the image of
the central-stone- containing positive fibre under the explicit puncture
map. It does not exclude other right-stone placements. T\_1(3), fully
constructed in §6, is a positive tiling with the required two stones and
supplies the opposite, whole-fibre witness. No solver nonexistence
status is used.

\section{9. Scope, source provenance and
continuation}\label{edition-10-section-16}

The earlier turn-5 work supplied a positive packing and finite
dual-guided repairs. This turn supplies an exactly ranked stable
residual, the complete one-stone transport, two unbounded families of
invariant values, and finite positive patches with proved all-parameter
embedding maps. The parent workbench sources are preserved. The original
problem and the known Kim--Propp endpoint retain their source
attribution.

What remains outside these completed families is the positive completion
at general invariant Δ and, for a certificate with a high embedding
threshold, its explicitly listed lower parameter cases. The stable
support and full tile-incidence diagram now specify those objects
exactly. No unproved uniform positive section, external analytic
estimate, or vanishing of the entire comparison kernel is presented as a
finished result.

The target is still the full original P7 existence statement. These
subfamily results are not a declaration of its resolution, and the role
of Split-Zero is identified by the actual supported complexes, original
cochain homotopy and positive fibres in §§5 and 7. The combinatorial
partition proofs and certificates supply the positive lifts; they have
not been inferred from a support label or an ordinary zero alone.

\chapter{Turn 8: unbounded deletion count}\label{edition-11}

Source: \path|sources/benzel-p7-turn8/benzel-p7-turn8/turn8/PROOF.md|.

SHA-256:
\path|f8662d8f26971accf8c27e377b2ba77b72a25018402ba70f52703869158d6ed9|.

16 September 2026. This is a written constructive research argument with
exact integer-Laurent-polynomial replay. Independent mathematical
review, a new Lean build, and novelty certification have not been
performed. The complete original Propp Problem 7 is not declared
resolved. All original cell coordinates, tile orientations, integer
coefficients, support labels and positive filling fibres are retained.
The two remaining nonzero deletion residue classes are recorded as
unfinished research, not as premises of the result below.

\section{1. The completed three-parameter
family}\label{edition-11-section-1}

Write t\_n=n(n-1)/2. For every three integers

\begin{verbatim}
k>=1,  m>=3k+2,  d>=m,
q=3k,  Delta=t_m-3k,  h=t_d-t_m+3k,
\end{verbatim}

this note constructs a type-103 tiling of the original benzel

\begin{verbatim}
W_h(d)=V(d+3h,2d+3h).
\end{verbatim}

It contains exactly Delta right stones and 3h(h+d) bones. Its reflected
tiling is constructed too. No search or optimization is called by the
constructor. Thus the deletion count is unbounded: q ranges over every
positive multiple of three in its full canonical domain q\textless=m-2.
This is not just a list of fixed-q examples. The preceding
triangular-invariant construction supplies q=0.

The original parameter map has inverse d=b-a and h=(2a-b)/3 on the
displayed integer lane. The original axial cell is (x,y), with
differences

\begin{verbatim}
u=y-x,  v=1-x-2y,  w=2x+y-1.
\end{verbatim}

V(a,b) consists of the cells with each difference in {[}1-a,b-1{]}. The
four permitted original tiles and their exact cell offsets are

\begin{verbatim}
R: (0,0),(1,0),(0,1);
H: (0,0),(1,0),(2,0);
V: (0,0),(0,1),(0,2);
D: (0,0),(1,-1),(2,-2).
\end{verbatim}

Their incidence boundary sends a placed tile to its three original cell
generators. Rotation and reflection are

\begin{verbatim}
rho(x,y)=(y,1-x-y),     F(x,y)=(x,1-x-y).
\end{verbatim}

The inverses are rho\^{}2 and F. They induce (u,v,w)-\textgreater(v,w,u)
and (u,v,w)-\textgreater(-w,-v,-u). On tiles rho cycles H,V,D and
preserves R; F exchanges H,D and preserves V,R. Consequently F carries
the constructed tiling of V(a,b) to a tiling of V(b,a), with its inverse
unchanged.

\subsection{Retained parent construction}\label{edition-11-section-2}

The complete preceding constructive proof is preserved in
\path|sources/turn6-PROOF.md|, §§1--3, not replaced by a tileability
assumption. Its code is retained as \texttt{packing.py} with original
bytes (the filename alone changes). For clarity, all input maps needed
here are stated explicitly.

Let r(n) be the unique positive integer with
t\_r\textless n\textless=t\_(r+1). In reflected coordinates the parent
packing consists of horizontal bands

\begin{verbatim}
H(L_y+3j,y), 0<=j<b_y, 1<=y<=2h+d,
(L_y,b_y)=(1-2y-r(y+Delta),y)                    for y<=h,
(L_y,b_y)=(y-3h-d+1,min(h,2h+d-y))             for y>h.
\end{verbatim}

Apply F, F rho and F rho\^{}2 to the bands. The resulting matching is
\path|Gamma_Delta(d,h).| The parent proof verifies all six original
containment inequalities, pairwise sector disjointness and the exact
complement:

\begin{verbatim}
partial Gamma_Delta(d,h)+1_(Omega_Delta)=1_(W_h(d)).                 (1)
\end{verbatim}

The rank map defining that complement is a bijection on the original
cells:

\begin{verbatim}
(r,s,p) -> rho^p(r-s,-s),  r>=1, 0<=s<r, p=0,1,2;
rank=t_r+s+1.
\end{verbatim}

The inverse takes the unique minimum among u,v,w, rotates it to u, and
recovers r=-u and s=-y. The cyclic differences are 2 modulo 3, so the
minimum is unique. Omega\_Delta consists precisely of ranks at most
Delta, and has 3Delta cells. In particular \path|Omega_(t_m)=W_0(m).|
The exact base tiling is

\begin{verbatim}
T_0(m)={R(m-2-2r-s,r-s): r,s>=0, r+s<=m-2}.                       (2)
\end{verbatim}

The phase/cell bijection and original region count in the retained proof
give

\begin{verbatim}
|W_h(d)|=3Delta+9h(h+d).                                         (3)
\end{verbatim}

No assertion concerning all positive fillings is inferred from (1).

\section{2. The actual source bones, with their receiving
indices}\label{edition-11-section-3}

Fix k,m in the result's domain. A source label has sector p=0,1,2, row y
and depth ell from the right end of a band. Set

\begin{verbatim}
r_y=m-1 for y<=3k,    r_y=m for y>3k,
a(p,y,ell)=F rho^p H(y-r_y-2-3ell,y).                             (4)
\end{verbatim}

The following list specifies the entire source matching A\_(m,k);
include all three sectors for every listed (y,ell):

\begin{itemize}
\tightlist
\item
  for 0\textless=ell\textless=k-2: \path|ell+1<=y<=m+2k-4,| and also the
  two rows y=m+2k-2+ell and y=m+2k-1+ell;
\item
  for ell=k-1: k\textless=y\textless=m+3k-2;
\item
  for k\textless=ell\textless=2k-2: ell+1\textless=y\textless=2k.
\end{itemize}

Empty ranges contribute no tiles. Every row is in 1..m+3k, and
\path|0<=ell<min(y,3k).| The two added rows are strictly above the first
range. Thus no source labels are repeated. Counting these ranges gives

\begin{verbatim}
|A_(m,k)|=3km+6k^2-3.                                           (5)
\end{verbatim}

These are actual generators of \path|Gamma_Delta(d,h),| even at d=m.
Here is the exact source-to-target comparison, including its inverse
index map.

The equality \path|t_(m-1)<y+Delta<=t_m| holds for
1\textless=y\textless=3k, because 3k\textless=m-2;
\path|t_m<y+Delta<=t_(m+1)| holds for 3k\textless y\textless=m+3k. These
prove the two branches of r\_y.

For d=m, h=3k. Rows y\textless=3k are prefix bones with index j=y-1-ell.
Rows y\textgreater3k are finite tail bones with index j=3k-1-ell and
tail start L=y-9k-m+1. Then

\begin{verbatim}
L+3j=y-m-2-3ell,
\end{verbatim}

which is exactly the original anchor in (4). The inverse to the tail
index map is \path|j_stable=j_tail+y-3k,| recovering
\path|j_stable=y-1-ell.| All indices lie in the original ranges, not in
enlarged fictitious bands.

For d\textgreater=m+1, \path|h=t_d-t_m+3k>=m+3k.| Every source row is a
prefix row, with j=y-1-ell. The original anchor is again (4), and the
index inverse is the identity. Hence A\_(m,k) is a subset of the
original parent matching at every exterior d\textgreater=m. In
particular its cells are disjoint from Omega\_Delta by (1).

This comparison preserves the sector, original row, tile orientation and
all three cells. Its linear extension commutes with the original cell
incidence map. No finite tail is silently identified with a prefix.

\section{3. Stones, long strips and explicit corner
replacements}\label{edition-11-section-4}

\subsection{The actual retained stones}\label{edition-11-section-5}

Remove from T\_0(m) the following original 3k stones:

\begin{verbatim}
D_(m,k)={rho^p R(m-2-a,-a): 0<=a<k, p=0,1,2}.                    (6)
\end{verbatim}

The base indices for the three rotations are obtained by permuting
(0,a,m-2-a). Since m-2\textgreater=3k, the corner ranges are disjoint;
all tiles in (6) are distinct members of (2). Let R\_(m,k)=T\_0(m) minus
this set. Its cardinality is t\_m-3k=Delta.

The last 3k cell orbits of W\_0(m) are

\begin{verbatim}
M_(m,k)= union_(i=1)^(3k) Orb_rho(i,i+1-m).                       (7)
\end{verbatim}

Indeed these are the \path|rank-(t_m-3k+1)..t_m| cells under the
displayed inverse. Consequently \path|Omega_Delta=W_0(m)| minus
M\_(m,k), as actual cell sets.

\subsection{The new bones}\label{edition-11-section-6}

First insert the three rotated copies of the original horizontal strips

\begin{verbatim}
H(y+m+3ell,y),
0<=ell<k,   1-m+k-ell<=y<=-2k.                                  (8)
\end{verbatim}

The row count m-3k+ell is at least 2+ell. These ranges therefore require
no extension of an index or limiting argument at the smallest m.

The corner table P\_k below is written in coordinates (x,z). Its actual
original embedding is (x,z)-\textgreater(x,z-m), followed by rho\^{}p,
p=0,1,2. The inverse first applies rho\^{}(-p), then sends
(x,y)-\textgreater(x,y+m). Every expression below is a placed original
tile, not a new prototile.

Horizontal corner families:

\begin{verbatim}
H(2+3j,2),                  0<=j<k-1;
H(y+3j,y),                  3<=y<=k, 0<=j<=k-y;
H(3+2r+3j,1-r),            r,j>=0, r+j<=k-2.
\end{verbatim}

Diagonal corner families:

\begin{verbatim}
D(1-k+3j,4-k-3j),           0<=j<k;
D(2-n+3j,2-n-3j),          1<=n<k, 0<=j<n;
D(2-n+3j,3-n-3j),          1<=n<k, 0<=j<n.
\end{verbatim}

Vertical corner families:

\begin{verbatim}
V(1-2k+i,k-2i+3j),          0<=i,j<k;
V(1-k,5-k+3j),              0<=j<k-1;
V(2-k+i,4-k+i+3j),          i,j>=0, i+j<=k-2;
V(x,2-3k-2x),               1-k<=x<=-1;
V(2i,3-3k-i),               0<=i<k;
V(2i+1,2-3k-i),             0<=i<k;
V(2i+2,2-4i+3j),            1<=i<=k-2, 0<=j<i;
V(2i+3,1-4i+3j),            1<=i<=k-2, 0<=j<i;
V(2k,6-4k+3j),              0<=j<k-1.
\end{verbatim}

At k=1 the table is the explicit four-tile set

\begin{verbatim}
P_1={V(-1,1), V(0,0), V(1,-1), D(0,3)}.                         (9)
\end{verbatim}

It is handled separately in the symbolic proof, so an expression with an
outer length k-2 is never interpreted as a negative number of terms.

Let B\_(m,k) be (8) and the three actual \path|embedded/rotated| corner
tables. For the corner there are k(k-1) H bones, k\^{}2 D bones and
(5k\^{}2+3k-2)/2 V bones. The strip has km-(5k\^{}2+k)/2 bones per
sector. Thus

\begin{verbatim}
|B_(m,k)|=3km+6k^2-3=|A_(m,k)|.                                (10)
\end{verbatim}

Cardinality is not used by itself to claim coverage. The full
coefficient identity is proved next.

\section{4. Exact identity with both parameters
retained}\label{edition-11-section-7}

For a finite original cell vector, the map {[}x,y{]}-\textgreater X\^{}x
Y\^{}y identifies its free integer module with the Laurent polynomials
of finite support. The inverse is coefficient extraction on every
original lattice cell. The four tile boundaries are represented exactly
by

\begin{verbatim}
r=1+X+Y, h=1+X+X^2, v=1+Y+Y^2, d=1+XY^(-1)+X^2Y^(-2).
\end{verbatim}

Work first in the integral domain

\begin{verbatim}
Z[X^±1,Y^±1,Mx^±1,My^±1,Kx^±1,Ky^±1].                           (11)
\end{verbatim}

The parameter specialization fixes X,Y and sends

\begin{verbatim}
(Mx,My,Kx,Ky) -> (X^m,Y^m,X^k,Y^k).
\end{verbatim}

Its kernel is the ideal generated by Mx-X\^{}m, My-Y\^{}m, Kx-X\^{}k,
Ky-Y\^{}k. The quotient isomorphism has inverse induced by the inclusion
of X,Y; reducing the four relations proves both inverse identities. No
evaluation of X,Y at numbers occurs in the proof.

A family with anchors

\begin{verbatim}
x=x_m m+x_k k+x_i i+x_j j+x_0,
y=y_m m+y_k k+y_i i+y_j j+y_0,
0<=i<I, 0<=j<J_0+g i
\end{verbatim}

has its exact cell polynomial obtained from the tile boundary, anchor
monomial, and the finite double sum. Put
R=X\textsuperscript{(x\_i)Y}(y\_i), S=X\textsuperscript{(x\_j)Y}(y\_j).
Its index factor is

\begin{verbatim}
[ Geo(R,I) - S^(J_0) Geo(R S^g,I) ]/(1-S),
Geo(Z,N)=(1-Z^N)/(1-Z).                                         (12)
\end{verbatim}

Equation (12) follows by summing 1+S+\ldots+S\^{}(J\_0+gi-1) first, then
summing in i. All actual index lengths in the domain
k\textgreater=2,m\textgreater=3k+2 are nonnegative. The powers
containing m or k remain the independent formal monomials in (11). Every
denominator contains X,Y only and is nonzero. Single-index families use
Geo directly. The separate k=1 calculation uses (9) and nonnegative
index lengths.

\subsection{Shared exact family data}\label{edition-11-section-8}

\path|families.json| is the single literal family table consumed by both
the finite constructor and the symbolic proof. An anchor row has
coefficients (m,k,i,j,constant), I has coefficients (m,k,constant), and
J has coefficients (m,k,constant,i). The exact affine expansions of the
corner and strips in §3 are those rows. The six source blocks below
exhibit the bijection with §2:

\begin{longtable}[]{@{}lllll@{}}
\toprule\noalign{}
Block & Depth ell & Original row y & Outer index & Inner index \\
\midrule\noalign{}
\endhead
\bottomrule\noalign{}
\endlastfoot
A1 & i & i+1+j & 0\textless=i\textless k-1 &
0\textless=j\textless3k-i \\
A2 & i & 3k+1+j & 0\textless=i\textless k-1 &
0\textless=j\textless m-k-4 \\
A3 & i & m+2k-2+i+j & 0\textless=i\textless k-1 &
0\textless=j\textless2 \\
A4 & k-1 & k+i & 0\textless=i\textless2k+1 & one term \\
A5 & k-1 & 3k+1+i & 0\textless=i\textless m-2 & one term \\
A6 & k+i & k+i+1+j & 0\textless=i\textless k-1 &
0\textless=j\textless k-i \\
\end{longtable}

In A1,A4,A6 one has y\textless=3k and the source x is y-m-1-3ell. In
A2,A3,A5 one has y\textgreater3k and x=y-m-2-3ell. Substitution gives
the source anchors in \path|families.json| term by term. Splitting the
first row interval at 3k and applying the inverse formulas i=ell or
ell-k, j=y-(displayed initial row) recovers every source label exactly
once. This proves the table correspondence without extrapolating
samples.

Rotation acts on cell polynomials by

\begin{verbatim}
Rho(f)=Y f(Y^(-1),X Y^(-1)),
\end{verbatim}

with the same linear exponent change on (Mx,My) and (Kx,Ky). Reflection
is Y f(XY\textsuperscript{(-1),Y}(-1)). These are additive, invertible
maps on the original cell modules; the factor Y is applied once to a
whole cell polynomial, not once per factor. Denominators receive the
linear exponent change without that affine translation. This is
precisely the original rho,F action, and explains the
module-versus-unital-algebra typing in the evaluator.

Let F\_A,F\_B be the original source and replacement cell polynomials,
F\_Del the sum of the deleted stone boundaries, and F\_M the missing
cell indicator (7). Applying (12) and the displayed affine actions to
the finite family table gives

\begin{verbatim}
F_B-F_A-F_Del+F_M=0.                                            (13)
\end{verbatim}

Here is a finite, independently replayable exact certificate for (13).
Multiply by the product

\begin{verbatim}
(1-Y^3)(1-XY^-2)(1-XY)(1-X^2Y^-4)(1-X^2Y^-1)
   (1-X^2Y^2)(1-X^3Y^-3)(1-X^3)(1-X^4Y^-2).                    (14)
\end{verbatim}

The resulting numerator has every integer coefficient zero in all six
independent Laurent variables. \path|bivariate_certificate.py| expands
the shared families, forms (13), applies the stated affine maps,
multiplies by (14), and verifies coefficient cancellation exactly. The
main calculation has 222 rational terms and 28,128 expanded numerator
terms before collection; no surviving coefficient remains. At k=1, a
separate 39-term identity has 510 expanded terms and zero surviving
coefficients. That calculation uses only m\textgreater=5.

When a denominator exponent is negative, the code applies the exact
identity

\begin{verbatim}
1/(1-X^-a Y^-b)=-X^a Y^b/(1-X^a Y^b).
\end{verbatim}

Both sides are elements of the same explicitly localized ring. This
retains its monomial factor and sign. It is not a parameter substitution
or discarded unit.

The localization of the integral domain (11) by the nonzero factors in
(14) is injective. Parameter specialization leaves those same nonzero
X,Y factors, and the target Laurent ring is also an integral domain.
Thus the zero numerator returns to the original finite cell identity,
for every parameter in the domain:

\begin{verbatim}
partial B_(m,k)+partial R_(m,k)
   =1_(Omega_Delta)+partial A_(m,k).                            (15)
\end{verbatim}

One can inspect/replay (13) as finite integer arithmetic from the
printed index table; no optimization infeasibility status or finite
range of m,k supplies its proof. \texttt{verify.py} separately checks
the family/index maps and original cells.

\section{5. Positivity and the full original
tiling}\label{edition-11-section-9}

The source A\_(m,k) is a matching inside the original Gamma by §2, and
its cells are disjoint from Omega by (1). Hence the right side of (15)
is exactly the all-ones vector of

\begin{verbatim}
Lambda=Omega_Delta disjoint-union cells(A_(m,k)).
\end{verbatim}

Every term on the left of (15) is an original placed tile with
coefficient +1. Equality of original coefficients forces every tile to
stay inside Lambda, forces every cell multiplicity to equal one, and
excludes duplicate tiles and overlaps. It therefore supplies the
positive patch tiling, not just a signed solution or cardinality
equality.

Adjoin the frozen original bones Gamma minus A. The explicitly
constructed full tiling is

\begin{verbatim}
(Gamma_Delta(d,h) minus A_(m,k)) union B_(m,k) union R_(m,k).       (16)
\end{verbatim}

Equation (1) proves the frozen cells and patch cells are disjoint and
exhaust the original W\_h(d). The counts follow from (5),(10) and
\textbar R\textbar=Delta. Applying F proves the reflected result. This
completes the construction asserted in §1.

The construction is invariant under rho on the original placed tiles.
Both A and B are unions of three rotated sectors, and the deleted stones
are full rho-orbits. The parent Gamma and base T0 have the same
property. Thus (16) gives an actual rho-invariant positive tiling on
this family.

\section{6. Full integral comparison kernel and explicit cell
duals}\label{edition-11-section-10}

Use the actual free cell module C=Z{[}Lambda{]}, old boundary image
\path|B_A=im(partial_A),| and new bone image \path|B_B=im(partial_B).|
The original generator inclusion induces

\begin{verbatim}
H_A=C/B_A -> H_(A+B)=C/(B_A+B_B).
\end{verbatim}

The representative-preserving map gives the exact isomorphism

\begin{verbatim}
B_B/(B_A intersect B_B) -> ker(H_A -> H_(A+B)).                   (17)
\end{verbatim}

It sends the class of b to {[}b{]} in H\_A. Its kernel is B\_A intersect
B\_B. A kernel representative decomposes as a+b, so b is an inverse
preimage; changing that decomposition changes b by precisely the
intersection. This proves both inverse maps, retaining the actual old
boundaries. It is the original killed-class comparison of Split-Zero D7,
not a replacement by dimensions.

There is a complete integer basis construction for (17). Label old and
new bones as separate generators even when a placed shape occurs in
both. Make the bipartite graph with those vertices and each shared
original cell as its edge. Mark a vertex having a cell with no owner in
the other matching. Call a component closed when it contains no marked
vertex.

An equality sum a\_i partial A\_i = sum b\_j partial B\_j implies
a\_i=b\_j at every shared cell and coefficient zero at a singly owned
cell. Propagation gives zero on each marked component and a single
constant integer on each closed component. Conversely each closed
component gives the actual identity

\begin{verbatim}
sum_(old vertices) partial A_i = sum_(new vertices) partial B_j.
\end{verbatim}

These generators have disjoint vertex supports and are primitive. Hence
they are an integral basis of the intersection. Choose one new-bone
pivot in every closed component. If there are b new bones and c closed
components, an explicit basis of (17) consists of {[}partial B\_j{]} for
every nonpivot new bone, and

\begin{verbatim}
K is isomorphic to Z^(b-c).                                    (18)
\end{verbatim}

No claim that c vanishes for all parameters is needed or made. For a
new-bone coefficient vector, subtract its pivot coefficient on each
closed component, then keep the nonpivot coefficients. The inverse
inserts zero at each pivot. Their difference is exactly the printed
intersection combination above.

The dual left inverse uses original cell cochains. Add a ground vertex
for singly owned cells, orient shared edges old-\textgreater new,
old-only edges old-\textgreater ground, and new-only edges
ground-\textgreater new. Root a spanning forest at ground in the open
part and at the selected new pivot in each closed component. For a
nonpivot new vertex, send one unit along the root-to-vertex forest path.
Assign each original cell edge +1 or -1 according to its orientation on
that path and zero otherwise. Vertex divergence is zero at every old
vertex, +1 at the chosen new vertex and -1 at a closed pivot (or
ground). Thus its cell pairing annihilates all old boundaries and is the
Kronecker delta on the basis (18). This explicitly constructs both
integral coordinate maps and their duals; no numerical rank or rational
nullspace computation is used. \path|cohomology.MatchingKernel|
implements these original graph, path and cell maps, including
intersection preimages.

The particular repaired class is

\begin{verbatim}
z=1_(Omega_Delta)-partial R_(m,k)
  =partial(B_(m,k)-A_(m,k)).                                   (19)
\end{verbatim}

The original cell p\_m=(m-2,0) belongs to Omega (its rank is
\path|t_(m-2)+1<=Delta),| is covered by the deleted R(m-2,0), and is not
covered by a retained stone. Every old bone lies outside Omega.
Coefficient evaluation at p\_m therefore annihilates B\_A and sends z to
1. The map n-\textgreater{[}nz{]} is a split injection of Z into H\_A.
Under the actual comparison it becomes zero with the original boundary
preimage (19). This summand is retained together with the complete
kernel (17), not used as a substitute for the latter.

\section{7. Actual Split-Zero supports, augmented cycles, and exterior
maps}\label{edition-11-section-11}

Let I index the actual old bones A\_(m,k), and take

\begin{verbatim}
L={bottom} disjoint-union P(I).
\end{verbatim}

The bottom cell support is empty. The nonbottom empty-release label has
cell support Omega\_Delta. For S subset I put

\begin{verbatim}
Lambda_S=Omega_Delta union union_(i in S) cells(A_i).
\end{verbatim}

The disjoint source partition makes S-\textgreater Lambda\_S a
join-preserving injection, with inverse reading the released old
generators; bottom goes to the empty cell set. Define original module
complexes

\begin{verbatim}
C_S^0=Z[all original contained bones], C_S^1=Z[Lambda_S], C_S^2=0,
d^0=partial, d^1=0.
\end{verbatim}

All transitions are original generator inclusions. Applying the pinned
D1--D8 reconstruction gives \path|Total^i=disjoint| union\_S C\_S\^{}i
over G(Z)=Z disjoint-union\{tau\}. For e=0\_Z supported in G(Z), its
operations are

\begin{verbatim}
(S,x)+(T,y)=(S union T, jx+jy),
r^bullet(S,x)=(S,rx),  e(S,x)=(S,0), tau(S,x)=(bottom,0).
\end{verbatim}

The differential square is (S,x)-\textgreater(S,0) in its actual target
degree. Its label projection retains S; the global absent zero has
bottom label. At each support the cycles and original boundaries give
the internal quotient q(S,p)=(S,{[}p{]}). It coequalizes the original
incidence map and its supported-zero companion. For a split-linear map f
equalizing them, p-p' in the old boundary module gives

\begin{verbatim}
f(S,p)=f(S,p')+f(S,p-p')=f(S,p')+f(S,0)=f(S,p').
\end{verbatim}

The last equality is the image of fibre-zero absorption, so no
cancellation in an arbitrary target is assumed. This defines the unique
descended split-linear map; lifting original representatives verifies
its scalar and addition laws. The comparison (17) is natural for the
same support transitions. Thus (19) becomes its own supported zero, not
tau. The nonempty support is retained.

Positive repairs have a genuinely linear augmented model. Let
\path|v_S=1_(Lambda_S)| and eta\_ST=sum\_(i in T minus S){[}A\_i{]}. The
original partitions give

\begin{verbatim}
v_T=jv_S+partial eta_ST,
eta_SU=j eta_ST+eta_TU.
\end{verbatim}

Use all original contained permitted tile generators, including R, and
put

\begin{verbatim}
Chat_S^0=Z direct-sum Z[contained permitted tiles],
dhat_S(r,n)=partial n-r v_S,
jhat_ST(r,n)=(r,jn+r eta_ST).                                   (20)
\end{verbatim}

At bottom the entire module is zero. Direct expansion gives

\begin{verbatim}
dhat_T jhat_ST=j dhat_S,    jhat_TU jhat_ST=jhat_SU.
\end{verbatim}

Charge-one nonnegative integral cycles are exactly original positive
patch tilings: the cell sums are one. Inclusion into the signed integral
cycle fibre keeps the same charge, tiles and coefficients, with inverse
on its image given by those same coordinates. No vanishing class
supplies positivity by itself; (15) supplies the actual positive cycle
at full release.

The exterior source-to-tail maps of §2 intertwine (20) on every original
generator. Adjoining frozen bones \path|f_d=Gamma_Delta(d,h)| minus A
defines

\begin{verbatim}
(r,n)->(r,n+r f_d).
\end{verbatim}

Since partial \path|f_d=1_(W_h)-v_I,| this is the actual augmented
cochain map to the full benzel. At charge one it adjoins the original
frozen tiling. Removing those same frozen bones is inverse on the full
tilings containing them. This does not assert that every tiling of the
full benzel contains the displayed frozen set.

\section{8. Rotation, the integer norm, and the remaining two
residues}\label{edition-11-section-12}

The threefold symmetry of this construction is explicit. On original
right-stone anchors,

\begin{verbatim}
rho R(x,y)=R(y,-x-y).
\end{verbatim}

Equality with R(x,y) forces x=y=-x-y, hence x=y=0. Thus the central
right stone is the only fixed placed right stone; every other orbit has
three distinct tiles. For any rho-invariant original tiling, the
right-stone count is therefore 0 or 1 modulo 3. This proves a specific
obstruction for invariant Delta=2 modulo 3: there is no rho-invariant
positive tiling there. It does not prohibit asymmetric positive tilings.
The actual inclusion of fixed tilings into all tilings and the original
action rho give the comparison, rather than an assertion of
unrelatedness.

On integer tile and cell modules retain the norm map N=1+rho+rho\^{}2.
It commutes with partial term by term, and N\^{}2=3N. No division by
three is performed. On support labels rho induces the actual
join-preserving permutation of the old bone orbits, and these cochain
maps lift to G(Z) with e and tau retained.

For the still-unfinished q=3k+1 and q=3k+2 constructions, the following
actual bone-band homotopies were calculated. They are complete
identities, not a general positive completion theorem. In the original
Laurent cell module, put Z=XY for the first identity and Z=X\^{}2/Y for
the second:

\begin{verbatim}
sum_(i=0)^(n-1) [partial H(x+i-1,y+i+1)-partial V(x+i,y+i)]
  = X^xY^y ((XY)^n-1)(1-X^-1Y),

sum_(i=0)^(n-1) [partial D(x+2i,y-i+1)-partial H(x+2i,y-i)]
  = X^xY^y (1-(X^2/Y)^n)(Y-1).                                (21)
\end{verbatim}

These follow by multiplying the literal three-cell polynomials and
summing a finite geometric progression. The endpoint cell vectors and
every intermediate placed bone are explicit, and changing the sign
reverses the signed homotopy. A two-strip version, relevant to the first
remaining residue, is

\begin{verbatim}
sum_(i=0)^(n-1) [partial D(x+i-2,y+i+2)+partial D(x+i-2,y+i+3)
                -partial V(x+i,y+i)-partial V(x+i,y+i+3)]
 =X^xY^y (1-(XY)^n) C(X,Y),
C=(XY^4+XY^3+XY^2+XY+Y^3+Y^2)/X^2.                           (22)
\end{verbatim}

These identities transport endpoint obstructions between corners. Actual
finite supports and positive corner completions found in exploration are
separately recorded with their finite scope. A parameter-uniform
positive endpoint formula for the other two residues has not been
completed in this contribution. Solver statuses and zero jet values are
not promoted to such a theorem.

\section{9. Provenance, reproduction and acceptance
boundary}\label{edition-11-section-13}

The source input is the user's Split-Zero workbench, pinned at
\path|KokunoYumeto/zeta-function-research-reader@1c8ec52c85c173adc9f8403a8a26914955e2f5a9|,
\path|workbenches/splitzero-tandem/tex/support_diagrams.tex| D1--D8, and
\path|formal/splitzero/SplitZeroComplex.lean| (supported differential,
cycle condition, internal coequalizer). This note constructs its stated
benzel diagrams, maps, kernels and positive cycles; it does not claim to
rebuild the original Lean library or apply an unrelated analytic
estimate. The preceding parent source proof and code identities are
preserved in \path|SOURCES.json|.

The original problem is Propp's Problem 7 in the trimer list. A bounded
status search did not establish a later complete resolution; it is not a
complete literature or novelty audit. The formerly known endpoint
packing, base tiling, and previously saved fixed-q controls retain their
attribution and source history.

From this directory:

\begin{verbatim}
python bivariate_certificate.py
python verify.py --max-k 15 --output evidence/normal.json
python -O verify.py --max-k 15 --output evidence/optimized.json
\end{verbatim}

The symbolic proof checks both independent parameters before evaluation.
The finite replay separately checks positive cells, exterior source
indices and their inverses, independent region enumeration, integer
matching-kernel coordinates, cell-dual paths, supported scalar
operations, augmented maps and deliberate bad inputs. The stated
execution receipts list exactly the completed parameter windows. No
independent specialist acceptance, priority determination, full P7
resolution, or new Lean certificate is claimed. The remaining positive
endpoint calculations are research tasks, not premises advertised as
finished results.

\chapter{Receiving-morphism addendum}\label{edition-12}

Source:
\path|sources/benzel-p7-turn8/benzel-p7-turn8/MORPHISM_ADDENDUM.md|.

SHA-256:
\path|14b3afc66ef0d453f7ef845396223099d19649c929c6946e557810ea5a571255|.

16 September 2026. This records explicitly the inclusion maps between
the two complexes used in PROOF §§6--7. It does not identify their
homology groups or assign the matching-kernel rank to the full
contained-bone quotient.

For a release label S, put A\_S=\{A\_i:i in S\} and let B\_S be those
bones of the new matching whose three original cells lie in Lambda\_S.
Let U\_S be the set of \textbf{all} original bones contained in
Lambda\_S. The source, intermediate and full complexes have degree-zero
terms

\begin{verbatim}
Z[A_S],  Z[A_S] direct-sum Z[B_S],  Z[U_S],
\end{verbatim}

and the common degree-one term \path|C_S=Z[Lambda_S].| Their degree-zero
maps are

\begin{verbatim}
a -> (a,0),
(a,b) -> i_A(a)+i_B(b),
\end{verbatim}

where i\_A and i\_B send each original placed generator to that same
placed generator in U\_S. When a placed bone belongs to both matchings,
the second map adds its two tagged coefficients; it is not asserted
injective. Its kernel is explicitly the subgroup generated by
(e\_t,-e\_t) over t in A\_S intersect B\_S. Distinct original generators
are a basis of Z{[}U\_S{]}, so equality of their coefficients proves
this kernel formula.

At degree one all arrows are the identity of the original cell module.
Each arrow commutes with incidence on every generator. The support
transitions S subset T use actual generator inclusions in A, B, U and C,
so these arrows form morphisms of the original join-indexed diagrams and
hence genuine G(Z)-linear cochain maps after reconstruction.

At the full release support, the first induced homology arrow is exactly

\begin{verbatim}
C/im(partial_A) -> C/(im(partial_A)+im(partial_B)).
\end{verbatim}

Its kernel is the matching comparison computed in PROOF §6. The next
arrow is

\begin{verbatim}
C/(im(partial_A)+im(partial_B)) -> C/im(partial_U).
\end{verbatim}

Its kernel is exactly
\path|im(partial_U)/(im(partial_A)+im(partial_B)),| by the same
original-representative quotient map. No vanishing of that additional
kernel is claimed. The distinguished class has its explicit nonzero
representative and evaluation left inverse before the first arrow; its
explicit boundary B-A kills it there. Additional contained bones do not
invalidate this calculation, and their further quotient effect remains
in the displayed second kernel.

The augmented positive-cycle diagram permits all original R,H,V,D tiles
contained in Lambda\_S. The separate inclusions of the displayed
positive patch generators into that module retain their coordinates and
commute with the original incidence map. The complete tiling certificate
belongs to the charge-one nonnegative fibre of this larger diagram. Thus
the matching calculation, the full contained-bone cohomology and the
positive tiling fibre are connected by explicit maps, rather than
substituted for one another.


\clearpage
\part{Turn 9: asymmetric positive endpoint completion}
\input{turn9.tex}
\end{document}
